\documentclass[11pt,a4paper]{article}
\usepackage[spanish]{babel}
\usepackage[utf8]{inputenc}
\usepackage[T1]{fontenc}
\usepackage{amsmath,amssymb}
\usepackage{booktabs}
\usepackage{longtable}
\usepackage{tabularx}
\usepackage{array}
\usepackage[numbers]{natbib}
\usepackage{hyperref}
\usepackage{enumitem}
\usepackage{algorithm}
\usepackage{algpseudocode}
\usepackage{geometry}
\usepackage{tikz}
\usepackage{graphicx}
\usepackage{adjustbox}
%Includes "References" in the table of contents
\usepackage[nottoc]{tocbibind}

\usetikzlibrary{arrows.meta,positioning,fit,calc,backgrounds,shapes.geometric,shapes.multipart}
\geometry{margin=1in}
\hypersetup{colorlinks=true,linkcolor=blue,citecolor=blue,urlcolor=blue}

\tikzset{
box/.style={draw, rounded corners, align=center, minimum height=1cm, inner sep=6pt},
line/.style={-Latex, thick},
group/.style={draw, dashed, rounded corners, inner sep=8pt}
}

\newcommand{\fitfigure}[1]{\adjustbox{max width=0.98\linewidth,max totalheight=0.42\textheight,keepaspectratio}{#1}}

\title{Frankestein Transformer:\\Biblioteca Unificada de Codificador-Decodificador, CLI y Notas de Diseño Basadas en Investigación}
\author{
Erick F. Merino M. \\
Este trabajo no está afiliado \\
\texttt{erickfmm@gmail.com}
}
\date{Febrero 2026}

\begin{document}

\maketitle

\begin{abstract}
Frankestein Transformer presenta un conjunto de herramientas unificado y basado en configuración para la experimentación sistemática con arquitecturas modernas de transformers, abarcando diecisiete variantes de mezcladores de secuencia y veintitrés familias de optimizadores. El sistema admite tanto el modelado de lenguaje enmascarado (MLM) tipo codificador como la predicción autorregresiva (AR) del siguiente token tipo decodificador mediante una configuración flexible de clase de modelo y modo, con flujos de trabajo de ajuste fino especializados para ambas arquitecturas. Las contribuciones de investigación son tres: (i) un contrato de configuración estricto basado en esquemas que permite la experimentación reproducible a través de diversos mecanismos de atención, incluyendo atención softmax estándar, atención sigmoide, redes retentivas, modelos selectivos de espacio de estados, transformers de profundidad continua, enrutamiento adaptativo de profundidad (Mixture-of-Depths) \citep{zhu2026mixtureofdepthsattention}, aumento de memoria condicional (Engram) \citep{cheng2026conditionalmemoryscalablelookup}, patrones de atención dispersa y mecanismos con compuerta; (ii) un marco integral de enrutamiento de optimizadores que admite métodos de reducción de varianza (MARS, Adan, AdEMAMix), variantes eficientes en memoria (Adafactor, GaLore, Lion), enfoques sin programación de tasa de aprendizaje, precondicionadores de segundo orden (Shampoo, SOAP, Sophia) y optimizadores de bajo rango de la familia APOLLO (Apollo, Apollo-Mini, Q-Apollo) \citep{zhu2025apollosgdlikememoryadamwlevel}; y (iii) flujos de trabajo de extremo a extremo que abarcan el despliegue cuantizado mediante empaquetado ternario de pesos y el entrenamiento de incrustaciones de oraciones inspirado en SBERT, respaldados por una pila de aseguramiento de calidad ampliada con amplia cobertura de pruebas unitarias, validación de ejemplos YAML y ejecución de integración continua. El conjunto de herramientas implementa una interfaz de configuración basada en web que proporciona renderizado de formularios impulsado por esquemas con documentación en línea y validación en tiempo real. Este documento de referencia técnica incluye diagramas arquitectónicos, visualizaciones de flujo de ejecución, tablas de decisión y apéndices completos que sintetizan la literatura sobre arquitecturas de transformers, mecanismos de atención dispersa, arquitecturas con compuerta y familias de optimizadores.
\end{abstract}

\tableofcontents
\newpage

\section{Introducción}

\subsection{Motivación y Planteamiento del Problema}

Las arquitecturas de transformers han transformado fundamentalmente el panorama del modelado de secuencias en el procesamiento del lenguaje natural \citep{vaswani_attention_2017}, la visión por computadora y la biología computacional. El éxito de modelos como BERT \citep{devlin_bert_2019}, GPT y sus variantes ha establecido al transformer como el estándar de facto para el aprendizaje de representaciones en el aprendizaje profundo. Sin embargo, la rápida proliferación de innovaciones arquitectónicas presenta desafíos prácticos significativos para investigadores y profesionales.

Los últimos años han sido testigos de una explosión de mecanismos alternativos de atención y mezcla de secuencias, cada uno abordando limitaciones específicas de la atención softmax estándar: complejidad computacional cuadrática \citep{child_sparse_transformer_2019}, requisitos de inferencia eficiente en memoria \citep{sun_retentive_2023,gu_mamba_2023}, gestión selectiva de estados basada en contenido \citep{behrouz_titans_2025} y optimizaciones conscientes del hardware \citep{ramapuram_theory_2024}. Simultáneamente, la literatura sobre optimización se ha diversificado más allá del AdamW clásico, introduciendo técnicas de reducción de varianza \citep{xie_adan_2022,pagliardini_ademamix_2024,yuan_mars_2024}, variantes eficientes en memoria \citep{shazeer_adafactor_2018,zhao_galore_2024}, enfoques sin programación de tasa de aprendizaje \citep{defazio_road_2024} y métodos de precondicionamiento de segundo orden \citep{gupta_shampoo_2018,vyas_soap_2024}.

La consecuencia práctica es un ecosistema de investigación fragmentado donde la comparación experimental entre arquitecturas y optimizadores requiere un esfuerzo de ingeniería significativo. Los investigadores deben implementar y depurar múltiples variantes desde cero, garantizar tuberías de entrenamiento consistentes y gestionar espacios de hiperparámetros complejos. Esta fragmentación dificulta la reproducibilidad, ralentiza el progreso científico y aumenta la barrera de entrada para nuevos investigadores.

Este trabajo aborda estos desafíos mediante un conjunto de herramientas de experimentación unificado y basado en configuración, disponible en \url{https://github.com/erickfmm/frankestein-transformer}, que proporciona:
\begin{enumerate}[leftmargin=*]
    \item \textbf{Diseño Basado en Esquemas}: Un contrato de configuración estricto y validado que garantiza la reproducibilidad mientras admite diecisiete arquitecturas de mezcladores de secuencia distintas y veintitrés familias de optimizadores.
    \item \textbf{Entrenamiento Agnóstico a la Arquitectura}: Infraestructura de entrenamiento común que soporta líneas base de atención densa (estándar, sigmoide), alternativas recurrentes (RetNet, Mamba, bloques tipo EDO), enrutamiento adaptativo de profundidad (Mixture-of-Depths) \citep{zhu2026mixtureofdepthsattention}, bloques de memoria condicional (Engram) \citep{cheng2026conditionalmemoryscalablelookup}, patrones de atención dispersa (Sparse Transformer, Longformer, BigBird, SparseK, NSA, SpargeAttn, FASA) y mecanismos con compuerta (GLA, DeltaNet, Gated DeltaNet, Gated DeltaNet-2, HGRN2, FoX, Gated Softmax).
    \item \textbf{Marco de Enrutamiento de Optimizadores}: Grupos de hiperparámetros con prefijos que permiten un control detallado sobre incrustaciones, capas de normalización, bloques recurrentes, bloques de atención y otros subconjuntos de parámetros a través de diversos optimizadores, incluyendo la familia APOLLO (Apollo, Apollo-Mini, Q-Apollo) \citep{zhu2025apollosgdlikememoryadamwlevel}.
    \item \textbf{Flujos de Trabajo de Extremo a Extremo}: Despliegue integrado mediante cuantización (empaquetado ternario de pesos, activaciones INT8) y capacidades de incrustación de oraciones inspiradas en SBERT \citep{reimers_sentence-bert_2019}.
    \item \textbf{Configuración Interactiva}: Interfaz basada en web que proporciona renderizado de formularios impulsado por esquemas, validación en tiempo real y generación de comandos CLI.
    \item \textbf{Herramientas de Confiabilidad}: Pruebas automatizadas integrales con suites de pruebas unitarias, validación de presets YAML y automatización de CI en versiones de Python compatibles.
\end{enumerate}

La superficie de comandos del proyecto es:
\begin{center}
\texttt{frankestein-transformer}
\end{center}
con subcomandos para flujos de trabajo de codificador y decodificador:
\texttt{train}, \texttt{finetune}, \texttt{deploy}, \texttt{quantize}, \texttt{infer}, \texttt{sbert-train}, \texttt{sbert-infer}, \texttt{web-server}.

El comando \texttt{web-server} lanza un constructor de configuración basado en Streamlit que proporciona:
\begin{itemize}[leftmargin=*]
    \item Campos de formulario impulsados por esquemas con títulos de parámetros y descripciones detalladas
    \uta Información sobre herramientas en tiempo real y texto de ayuda para cada opción de configuración
    \item Vista previa YAML en vivo y funcionalidad de descarga
    \item Comandos CLI generados para entrenamiento, despliegue, inferencia y flujos de trabajo SBERT
\end{itemize}

Esta interfaz interactiva sirve como alternativa a la edición manual de YAML, mejorando la usabilidad para usuarios que exploran las opciones de configuración disponibles y comprenden su impacto en el comportamiento del modelo y la dinámica de entrenamiento.

\subsection{Contribuciones}

Este trabajo realiza las siguientes contribuciones principales:
\begin{enumerate}[leftmargin=*]
    \item \textbf{Esquema de Configuración Unificado}: Un contrato de esquema basado en YAML con validación estricta que soporta diecisiete arquitecturas de mezcladores de secuencia distintas en cuatro categorías (líneas base densas, bloques recurrentes/retentivos, patrones de atención dispersa y mecanismos con compuerta), junto con controles de profundidad adaptativa y memoria condicional, mientras garantiza la reproducibilidad mediante restricciones \texttt{additionalProperties: false}.
    \item \textbf{Soporte Integral de Arquitecturas}: Implementación de variantes modernas de transformers incluyendo atención softmax estándar \citep{vaswani_attention_2017}, atención sigmoide \citep{ramapuram_theory_2024}, RetNet \citep{sun_retentive_2023}, Mamba \citep{gu_mamba_2023}, transformers continuos tipo EDO \citep{zhang_continuous_2021}, atención con aumento de memoria Titans \citep{behrouz_titans_2025}, capas de memoria condicional Engram \citep{cheng2026conditionalmemoryscalablelookup}, enrutamiento de tokens Mixture-of-Depths \citep{zhu2026mixtureofdepthsattention}, mecanismos de atención dispersa \citep{child_sparse_transformer_2019,beltagy_longformer_2020,zaheer_bigbird_2020,lou_sparsek_2024,yuan_nsa_2025,zhang_spargeattn_2025,wang_fasa_2026} y arquitecturas de atención con compuerta \citep{yang_gla_2023,yang_deltanet_2024,yang_gated_deltanet_2024,hatamizadeh_gated_deltanet2_2026,qin_hgrn2_2024,lin_forgetting_transformer_2025,qiu_gated_attention_2025}. El sistema soporta tanto entrenamiento en modo codificador con atención bidireccional para modelado de lenguaje enmascarado (MLM) como entrenamiento en modo decodificador con enmascaramiento causal para predicción autorregresiva del siguiente token.
    \item \textbf{Marco de Enrutamiento de Optimizadores}: Sistema de hiperparámetros con prefijos que permite control por grupo de parámetros a través de veintitrés optimizadores, incluyendo Anon con adaptabilidad ajustable \citep{anon2026anon}, métodos de reducción de varianza (MARS, Adan, AdEMAMix, Cautious AdamW), variantes eficientes en memoria (Adafactor, GaLore, Lion), enfoques sin programación de tasa de aprendizaje (Schedule-Free AdamW), métodos conscientes de curvatura (Sophia), precondicionadores de segundo orden (Shampoo, SOAP), optimizadores orientados a ortogonalidad (Muon, Turbo-Muon), escalado de lotes grandes (LAMB) y la familia APOLLO (Apollo, Apollo-Mini, Q-Apollo) \citep{zhu2025apollosgdlikememoryadamwlevel}.
    \item \textbf{Cuantización y Despliegue}: Tubería de despliegue integrada que soporta empaquetado ternario de pesos y cuantización de activaciones INT8 con estimaciones de tamaño siguiendo la aproximación de almacenamiento de $1.58$ bits.
    \item \textbf{Flujos de Trabajo de Incrustación de Oraciones}: Tuberías de entrenamiento e inferencia inspiradas en SBERT que soportan puntuación de similitud, recuperación, agrupamiento y exportación persistente de incrustaciones.
    \item \textbf{Interfaz de Configuración Interactiva}: Servidor web basado en Streamlit que proporciona generación de formularios impulsada por esquemas, validación en tiempo real, documentación en línea y generación de comandos CLI.
    \item \textbf{Tubería de Validación Automatizada}: Integración continua y pruebas unitarias ampliadas que verifican las rutas de entrenamiento del modelo, la integración optimizador/esquema y la compatibilidad de ejemplos YAML.
\end{enumerate}

\subsection{Guía de Lectura}

Este documento está organizado como una referencia técnica que aborda cuatro preocupaciones operativas:
\begin{enumerate}[leftmargin=*]
    \item \textbf{Requisitos Previos}: El Apéndice~\ref{sec:annex-tutorial} proporciona una introducción accesible a los transformers y mecanismos de atención para lectores nuevos en el campo.
    \item \textbf{Contrato de Configuración}: La Sección~\ref{sec:config-centric-architecture} describe el esquema YAML que garantiza experimentos válidos y la Sección~\ref{sec:schema-validation} explica las reglas de validación.
    \item \textbf{Selección de Arquitectura}: La Sección~\ref{sec:normalization} cubre las opciones de normalización; la Sección~\ref{sec:attention-families} proporciona una comparación exhaustiva de las familias de mezcladores de secuencia; los Apéndices~\ref{sec:annex-transformer}, \ref{sec:annex-sparse} y \ref{sec:annex-gated} sintetizan la literatura de respaldo.
    \item \textbf{Dinámica de Optimización}: La Sección~\ref{sec:optimizer-families} detalla el enrutamiento de optimizadores y la dinámica de entrenamiento; el Apéndice~\ref{sec:annex-optimizers} proporciona un análisis exhaustivo de las familias de optimizadores.
    \item \textbf{Despliegue e Inferencia}: Las Secciones~\ref{sec:quantization} y~\ref{sec:sbert} describen el despliegue cuantizado y los flujos de trabajo de incrustación de oraciones.
\end{enumerate}

\subsection{Constructor de Configuración Basado en Web}

Además de la edición directa de YAML, este proyecto proporciona una interfaz web basada en Streamlit (accesible mediante el comando \texttt{web-server}) que mejora la accesibilidad y descubribilidad de la configuración. La interfaz presenta campos de esquema con:

\begin{itemize}[leftmargin=*]
    \item \textbf{Renderizado de formularios impulsado por esquemas} --- Todos los campos se generan dinámicamente a partir del esquema autoritativo, garantizando consistencia y validación.
    \item \textbf{Documentación de parámetros en línea} --- Cada campo del formulario muestra un \emph{título} del esquema como su etiqueta, con la \emph{descripción} mostrada como información sobre herramientas al pasar el cursor.
    \item \textbf{Vista previa de configuración en tiempo real} --- Los usuarios ven la salida YAML en vivo mientras modifican los campos del formulario, permitiendo retroalimentación de validación inmediata.
    \item \textbf{Generación de comandos CLI} --- La interfaz genera comandos CLI completos para entrenamiento, despliegue, inferencia y flujos de trabajo SBERT basados en la configuración actual.
    \item \textbf{Mejoras de accesibilidad} --- La información sobre herramientas y los formularios estructurados facilitan la comprensión de las opciones de configuración, especialmente para usuarios nuevos en el proyecto o que exploran arquitecturas novedosas.
\end{itemize}

Este enfoque basado en web aborda barreras de usabilidad comunes en la experimentación impulsada por configuración:
\begin{itemize}[leftmargin=*]
    \item Reduce la necesidad de memorizar la estructura YAML y los nombres de los campos
    \item Previene errores tipográficos mediante la validación del esquema
    \item Proporciona contexto educativo a través de documentación en línea
    \uta Permite la experimentación rápida con ajuste guiado de parámetros
    \item Sirve tanto como herramienta de configuración como recurso de aprendizaje
\end{itemize}

La implementación de la interfaz web utiliza widgets de formulario de Streamlit con:
\begin{itemize}[leftmargin=*]
    \item \texttt{st.checkbox()} para alternativas binarias con texto de ayuda
    \item \texttt{st.number\_input()} para campos numéricos con tamaño de paso y formato
    \item \texttt{st.selectbox()} para opciones de enumeración con visualización de opciones
    \item \texttt{st.multiselect()} para selecciones de arreglo a partir de opciones definidas
    \item \texttt{st.text\_input()} para campos de cadena
    \item \texttt{st.info()}, \texttt{st.caption()} para información complementaria
\end{itemize}

Los metadatos del esquema (campos de título y descripción) se extraen y renderizan sistemáticamente en todas las secciones del formulario, incluyendo:
\begin{itemize}[leftmargin=*]
    \item Parámetros de arquitectura del modelo (tamaño oculto, capas, cabezas de atención, etc.)
    \item Configuraciones de ejecución de entrenamiento (tamaño de lote, acumulación, programador)
    \item Configuración del optimizador con hiperparámetros por grupo de parámetros
    \item Opciones de despliegue y cuantización
    \item Parámetros específicos de entrenamiento e inferencia SBERT
\end{itemize}

\section{Trabajo Relacionado}

Este trabajo se sitúa en la intersección de tres áreas activas de investigación: arquitecturas de atención alternativas, mecanismos de atención dispersa y algoritmos de optimización avanzados. Esta sección proporciona un estudio conciso; las formulaciones matemáticas detalladas y las descripciones algorítmicas se difieren a los apéndices.

\subsection{Arquitecturas de Mezcladores de Secuencia}

\subsubsection{Líneas Base de Atención Densa}
La atención softmax estándar \citep{vaswani_attention_2017} logra un contexto global completo mediante interacciones por pares de $n^2$, pero incurre en costos prohibitivos de memoria y cómputo para secuencias largas. La atención sigmoide \citep{ramapuram_theory_2024} reemplaza la softmax por filas con una sigmoide elemento a elemento, logrando una convergencia más rápida y hasta un 17\% de aceleración del núcleo, aunque requiere estabilización con normalización híbrida a escala. Ambas sirven como líneas base densas en este conjunto de herramientas.

\subsubsection{Arquitecturas Recurrentes y Retentivas}
RetNet \citep{sun_retentive_2023} resuelve el ``triángulo imposible'' de entrenamiento paralelo, inferencia $\mathcal{O}(1)$ y rendimiento sólido mediante una formulación dual de atención-retención. Mamba \citep{gu_mamba_2023} introduce selectividad dependiente de la entrada en modelos de espacio de estados, logrando complejidad lineal con algoritmos de exploración conscientes del hardware. Los transformers tipo EDO \citep{zhang_continuous_2021} tratan la profundidad como integración numérica sobre un sistema dinámico continuo. Titans \citep{behrouz_titans_2025} aumenta la inferencia con memorización neuronal en tiempo de prueba para contextos extremadamente largos. En el Apéndice~\ref{sec:annex-transformer} se proporcionan formulaciones detalladas.

\subsubsection{Mecanismos de Atención Dispersa}
Los métodos de atención dispersa reducen el cuello de botella cuadrático restringiendo las interacciones entre tokens: Sparse Transformer \citep{child_sparse_transformer_2019} utiliza patrones factorizados con paso y fijos ($\mathcal{O}(n\sqrt{n})$); Longformer \citep{beltagy_longformer_2020} emplea ventanas deslizantes con tokens globales ($\mathcal{O}(n \cdot w)$); BigBird \citep{zaheer_bigbird_2020} combina caminos locales, aleatorios y globales; SparseK \citep{lou_sparsek_2024} aplica selección diferenciable top-$k$; NSA \citep{yuan_nsa_2025} utiliza un diseño jerárquico de tres ramas; SpargeAttn \citep{zhang_spargeattn_2025} realiza poda a nivel de bloques en dos etapas; y FASA \citep{wang_fasa_2026} aprovecha características de frecuencia de RoPE para la selección de tokens. Las descripciones completas se encuentran en el Apéndice~\ref{sec:annex-sparse}.

\subsubsection{Mecanismos de Atención con Compuerta}
Las arquitecturas con compuerta controlan el flujo de información mediante compuertas aprendibles: GLA \citep{yang_gla_2023} agrega compuertas diagonales dependientes de datos a la atención lineal; DeltaNet \citep{yang_deltanet_2024} aplica la regla de aprendizaje delta a las actualizaciones de estado recurrente; Gated DeltaNet \citep{yang_gated_deltanet_2024} sintetiza ambos mecanismos; Gated DeltaNet-2 \citep{hatamizadeh_gated_deltanet2_2026} desacopla la compuerta delta escalar en compuertas de borrado canalizadas (eje de claves) y escritura canalizada (eje de valores); HGRN2 \citep{qin_hgrn2_2024} introduce compuertas de olvido jerárquicas con expansión de estado mediante producto exterior; FoX \citep{lin_forgetting_transformer_2025} incrusta compuertas de olvido directamente en la atención softmax; y Gated Softmax \citep{qiu_gated_attention_2025} aplica compuertas de canal posteriores a SDPA. Las derivaciones matemáticas completas aparecen en el Apéndice~\ref{sec:annex-gated}.

\subsection{Algoritmos de Optimización}

La optimización de arquitecturas de transformers altamente parametrizadas requiere navegar paisajes de pérdida no convexos con espectros de Hessiana heterogéneos. La literatura se ha diversificado en varias familias más allá de la línea base clásica de AdamW \citep{loshchilov_decoupled_2017}:

\begin{itemize}[leftmargin=*]
    \item \textbf{Reducción de varianza y momento}: RAdam \citep{liu_variance_2019}, Adan \citep{xie_adan_2022}, ADOPT \citep{taniguchi_adopt_2024}, AdEMAMix \citep{pagliardini_ademamix_2024}, MARS \citep{yuan_mars_2024} y optimizadores Cautious \citep{liang_cautious_2024} abordan la inestabilidad en pasos iniciales, las garantías de convergencia y la mezcla de historial con EMA dual.
    \item \textbf{Eficientes en memoria}: Adafactor \citep{shazeer_adafactor_2018}, GaLore \citep{zhao_galore_2024}, Lion \citep{chen_symbolic_2023} y APOLLO \citep{zhu2025apollosgdlikememoryadamwlevel} reducen la memoria del estado del optimizador mediante factorización, proyección de bajo rango o actualizaciones basadas en signo.
    \item \textbf{Sin programación de tasa y libres de parámetros}: Schedule-Free AdamW \citep{defazio_road_2024} y Prodigy \citep{mishchenko_prodigy_2023} absorben el programador o el ajuste de la tasa de aprendizaje en la dinámica del optimizador.
    \item \textbf{De segundo orden y conscientes de curvatura}: Shampoo \citep{gupta_shampoo_2018}, SOAP \citep{vyas_soap_2024} y Sophia \citep{liu_sophia_2023} incorporan información aproximada de segundo orden.
    \item \textbf{Orientados a geometría}: Muon \citep{shen_convergence_2025} y Turbo-Muon \citep{boissin_turbo-muon_2025} remodelan la geometría de actualización mediante ortogonalización.
\end{itemize}

En el Apéndice~\ref{sec:annex-optimizers} se proporcionan descripciones algorítmicas detalladas, pseudocódigo y una comparación exhaustiva de complejidad.

\section{Diseño y Arquitectura del Sistema}
\label{sec:system-design}

\subsection{Arquitectura Centrada en Configuración}
\label{sec:config-centric-architecture}

La decisión de diseño principal en este repositorio es que la experimentación es \emph{esquema primero}. En lugar de exponer una gran cantidad de indicadores laxamente verificados, el proyecto obliga a la topología del modelo, la familia de optimizadores, los límites de entrenamiento y las opciones de telemetría a través de un único documento de configuración validado. Esto reduce la ambigüedad al reproducir resultados y permite comparar muchas arquitecturas bajo una interfaz operativa consistente.

El contrato autoritativo es \texttt{configs/schema.yaml}. Este impone tres objetos de primer nivel:
\begin{itemize}[leftmargin=*]
    \item \texttt{model\_class}
    \item \texttt{model}
    \item \texttt{training}
\end{itemize}

\paragraph{Selección de Clase de Modelo.}
El campo \texttt{model\_class} determina la variante arquitectónica instanciada por el pipeline de entrenamiento. Se soportan tres opciones:
\begin{itemize}[leftmargin=*]
    \item \textbf{frankenstein}: Modelos encoder de arquitectura mixta que soportan diversos mecanismos de atención (estándar, sigmoide, retentiva, espacio de estados, dispersa y mezcladores con compuerta) con MoE (Mezcla de Expertos) y características avanzadas. Optimizado para entrenamiento bidireccional tipo encoder con objetivos de modelado de lenguaje enmascarado (MLM).
    \item \textbf{mini}: Variante de encoder simplificada diseñada para escenarios de entrenamiento a menor escala. Proporciona una sobrecarga de parámetros reducida e iteración más rápida para experimentación y creación de prototipos.
    \item \textbf{frankesteindecoder}: Decoder causal autorregresivo para generación de siguiente token estilo LLM. Permite el enmascaramiento de atención causal para tareas de generación secuencial de texto. Cuando se selecciona esta clase, en tiempo de ejecución se fuerza \texttt{mode='decoder'}.
\end{itemize}

\paragraph{Selección de Modo de Entrenamiento.}
El campo \texttt{model.mode} controla el comportamiento de enmascaramiento de atención en todo el modelo:
\begin{itemize}[leftmargin=*]
    \item \textbf{encoder}: Utiliza atención bidireccional donde todos los tokens atienden a todos los demás tokens en la secuencia. Adecuado para tareas de preentrenamiento de modelado de lenguaje enmascarado (MLM) donde el modelo aprende a predecir tokens enmascarados aleatoriamente basándose en el contexto completo.
    \item \textbf{decoder}: Utiliza enmascaramiento causal donde cada token solo puede atender a tokens anteriores en la secuencia. Requerido para tareas autorregresivas (AR) de predicción del siguiente token, como modelado de lenguaje y generación de texto. Cuando \texttt{model\_class='frankesteindecoder'}, el sistema fuerza automáticamente \texttt{mode='decoder'} en tiempo de ejecución.
\end{itemize}

Este soporte de arquitectura dual permite al sistema manejar tanto preentrenamiento tipo encoder (MLM en contextos bidireccionales) como generación tipo decoder (predicción causal autorregresiva) a través de una interfaz de configuración unificada.

El \texttt{model.layer\_pattern} soporta bloques heredados, dispersos y con compuerta:
\begin{itemize}[leftmargin=*]
    \item \textbf{Retentive Network (RetNet)} --- referencia interna: \texttt{sun\_retentive\_2023} --- nombre de código: \texttt{retnet}, \texttt{retnet\_attn}
    \item \textbf{Mamba (Selective State Space Model)} --- referencia interna: \texttt{gu\_mamba\_2023} --- nombre de código: \texttt{mamba}
    \item \textbf{Bloque de Profundidad Continua estilo ODE} --- referencia interna: \texttt{zhang\_continuous\_2021} --- nombre de código: \texttt{ode}
    \item \textbf{Atención con Memoria Aumentada Titans} --- referencia interna: \texttt{behrouz\_titans\_2025} --- nombre de código: \texttt{titan\_attn}
    \item \textbf{Atención Softmax Estándar} --- referencia interna: \texttt{vaswani\_attention\_2017} --- nombre de código: \texttt{standard\_attn}
    \item \textbf{Auto-Atención Sigmoide} --- referencia interna: \texttt{ramapuram\_theory\_2024} --- nombre de código: \texttt{sigmoid\_attn}
    \item \textbf{Transformer Disperso} --- referencia interna: \texttt{child\_sparse\_transformer\_2019} --- nombre de código: \texttt{sparse\_transformer\_attn}
    \item \textbf{Longformer} --- referencia interna: \texttt{beltagy\_longformer\_2020} --- nombre de código: \texttt{longformer\_attn}
    \item \textbf{BigBird} --- referencia interna: \texttt{zaheer\_bigbird\_2020} --- nombre de código: \texttt{bigbird\_attn}
    \item \textbf{Atención SparseK} --- referencia interna: \texttt{lou\_sparsek\_2024} --- nombre de código: \texttt{sparsek\_attn}
    \item \textbf{Atención Dispersa Nativa (NSA)} --- referencia interna: \texttt{yuan\_nsa\_2025} --- nombre de código: \texttt{nsa\_attn}
    \item \textbf{SpargeAttn} --- referencia interna: \texttt{zhang\_spargeattn\_2025} --- nombre de código: \texttt{sparge\_attn}
    \item \textbf{FASA (Frequency-aware Sparse Attention)} --- referencia interna: \texttt{wang\_fasa\_2026} --- nombre de código: \texttt{fasa\_attn}
    \item \textbf{Atención Lineal con Compuerta (GLA)} --- referencia interna: \texttt{yang\_gla\_2023} --- nombre de código: \texttt{gla\_attn}
    \item \textbf{DeltaNet} --- referencia interna: \texttt{yang\_deltanet\_2024} --- nombre de código: \texttt{deltanet\_attn}
    \item \textbf{DeltaNet con Compuerta} --- referencia interna: \texttt{yang\_gated\_deltanet\_2024} --- nombre de código: \texttt{gated\_deltanet\_attn}
    \item \textbf{Gated DeltaNet-2} --- referencia interna: \texttt{hatamizadeh\_gated\_deltanet2\_2026} --- nombre de código: \texttt{gated\_deltanet2\_attn}
    \item \textbf{HGRN2} --- referencia interna: \texttt{qin\_hgrn2\_2024} --- nombre de código: \texttt{hgrn2\_attn}
    \item \textbf{Forgetting Transformer (FoX)} --- referencia interna: \texttt{lin\_forgetting\_transformer\_2025} --- nombre de código: \texttt{fox\_attn}
    \item \textbf{Atención Softmax con Compuerta} --- referencia interna: \texttt{qiu\_gated\_attention\_2025} --- nombre de código: \texttt{gated\_softmax\_attn}
\end{itemize}

que corresponde a la literatura actual de atención y mezcladores de secuencias \citep{sun_retentive_2023,gu_mamba_2023,zhang_continuous_2021,behrouz_titans_2025,ramapuram_theory_2024,vaswani_attention_2017,child_sparse_transformer_2019,beltagy_longformer_2020,zaheer_bigbird_2020,lou_sparsek_2024,yuan_nsa_2025,zhang_spargeattn_2025,wang_fasa_2026,yang_gla_2023,yang_deltanet_2024,yang_gated_deltanet_2024,hatamizadeh_gated_deltanet2_2026,qin_hgrn2_2024,lin_forgetting_transformer_2025,qiu_gated_attention_2025}

El \texttt{training.optimizer.optimizer\_class} soporta una amplia familia de optimizadores:
\texttt{sgd\_momentum}, \texttt{adamw}, \texttt{adafactor}, \texttt{galore\_adamw}, \texttt{prodigy}, \texttt{lion}, \texttt{sophia}, \texttt{muon}, \texttt{turbo\_muon}, \texttt{radam}, \texttt{adan}, \texttt{adopt}, \texttt{ademamix}, \texttt{mars\_adamw}, \texttt{cautious\_adamw}, \texttt{lamb}, \texttt{schedulefree\_adamw}, \texttt{shampoo}, \texttt{soap}, \texttt{anon}, \texttt{apollo}, \texttt{apollo\_mini}, y \texttt{q\_apollo}.

\begin{figure}[t]
\centering
\fitfigure{\begin{tikzpicture}[
font=\scriptsize,
box/.append style={inner sep=4pt, minimum height=0.62cm}
]
% Top level
\node[box, fill=green!10, minimum width=1.9cm] (yaml) at (0,0) {\textbf{YAML Config}};
\node[box, fill=green!8, minimum width=2.0cm] (validation) at (0,-1.55) {\textbf{Validation}\\\texttt{schema + rules}};
\node[box, fill=blue!10, minimum width=1.9cm] (web) at (-3.9,-2.95) {\textbf{Web}\\\texttt{Streamlit}};
\node[box, fill=blue!8, minimum width=2.1cm] (cli) at (3.9,-2.95) {\textbf{CLI}\\\texttt{train/deploy/infer}};
% Main blocks
\node[box, fill=orange!12, minimum width=1.9cm] (model) at (-2.6,-4.3) {\textbf{Model}\\17 mixers};
\node[box, fill=yellow!15, minimum width=2.0cm] (train) at (0,-4.3) {\textbf{Training}\\AMP + scheduler};
\node[box, fill=red!12, minimum width=1.9cm] (opt) at (2.6,-4.3) {\textbf{Optimizer}\\22 families};
% Outcomes
\node[box, fill=purple!12, minimum width=2.0cm] (modelout) at (-2.6,-5.95) {\textbf{Model Build}\\pattern + loops};
\node[box, fill=purple!10, minimum width=2.0cm] (trainout) at (0,-5.95) {\textbf{Train Loop}\\checks + logs};
\node[box, fill=purple!10, minimum width=2.0cm] (optout) at (2.6,-5.95) {\textbf{Opt Step}\\group routing};
% Deployment
\node[box, fill=cyan!12, minimum width=2.0cm] (deploy) at (0,-7.75) {\textbf{Deploy}\\quantization};
\node[box, fill=cyan!10, minimum width=2.0cm] (sbert) at (2.6,-7.75) {\textbf{SBERT}\\search/cluster};
% Arrows
\draw[line] (yaml) -- (validation);
\draw[line] (validation) -- (web);
\draw[line] (validation) -- (cli);
\draw[line] (validation) -- (model);
\draw[line] (validation) -- (train);
\draw[line] (validation) -- (opt);
\draw[line] (model) -- (modelout);
\draw[line] (train) -- (trainout);
\draw[line] (opt) -- (optout);
\draw[line] (trainout) -- (deploy);
\draw[line] (optout) -- (sbert);
\draw[line, dashed] (modelout) -- (trainout);
\draw[line, dashed] (optout.south) |- (deploy.east);
\end{tikzpicture}}
\caption{Arquitectura del sistema: la configuración fluye a través de la validación, se divide en modelo/entrenamiento/optimizador, ejecuta en tiempo real y produce artefactos de despliegue/SBERT.}
\label{fig:system-architecture}
\end{figure}

\subsection{Alcance del Esquema y Reglas de Validación}
\label{sec:schema-validation}

El esquema es estricto: los objetos de primer nivel y anidados establecen \texttt{additionalProperties: false}. Esto garantiza que las claves desconocidas fallen rápidamente en lugar de ser ignoradas silenciosamente. El objeto \texttt{training.optimizer.parameters} está adicionalmente restringido por reglas de prefijo específicas del optimizador mediante verificaciones de patrón \texttt{allOf}+\texttt{if/then}.

Los valores de normalización actualmente aceptados por el esquema son:
\[
\texttt{norm\_type} \in \{\texttt{layer\_norm}, \texttt{dynamic\_tanh}, \texttt{derf}\}
\]
Por lo tanto, \texttt{rms\_norm} \textbf{no} es un valor de esquema válido en el contrato actual.

\subsection{Inventario Completo de Características del Modelo}

\small
\begin{longtable}{p{5.1cm}p{2.4cm}p{7.0cm}}
\toprule
\textbf{Campo} & \textbf{Tipo/Rango} & \textbf{Significado} \\
\midrule
\endhead
\texttt{model\_class} & enum & Variante arquitectónica: \texttt{frankenstein} (encoder de arquitectura mixta), \texttt{mini} (encoder simplificado), \texttt{frankesteindecoder} (decoder autorregresivo). \\
\texttt{model.mode} & enum & Modo de atención: \texttt{encoder} (bidireccional, para MLM) o \texttt{decoder} (causal, para AR). Cuando \texttt{model\_class='frankesteindecoder'}, fuerza \texttt{mode='decoder'}. \\
\texttt{vocab\_size} & int $\ge 1$ & Tamaño del vocabulario. \\
\texttt{hidden\_size} & int $\ge 1$ & Dimensión oculta. \\
\texttt{num\_layers} & int $\ge 1$ & Número de capas físicas. \\
\texttt{num\_loops} & int $\ge 1$ & Número de bucles lógicos (bloques en bucle). \\
\texttt{num\_heads} & int $\ge 1$ & Cabezas de atención. \\
\texttt{retention\_heads} & int $\ge 1$ & Cabezas de retención para mezcladores estilo RetNet. \\
\texttt{num\_experts} & int $\ge 1$ & Número de expertos en MoE. \\
\texttt{top\_k\_experts} & int $\ge 1$ & Enrutamiento top-$k$ de expertos en MoE. \\
\texttt{dropout} & float $[0,1]$ & Dropout global. \\
\texttt{layer\_pattern} & array enum & Lista ordenada de bloques: heredados (\texttt{retnet}, \texttt{retnet\_attn}, \texttt{mamba}, \texttt{ode}, \texttt{titan\_attn}, \texttt{standard\_attn}, \texttt{sigmoid\_attn}), dispersos (\texttt{sparse\_transformer\_attn}, \texttt{longformer\_attn}, \texttt{bigbird\_attn}, \texttt{sparsek\_attn}, \texttt{nsa\_attn}, \texttt{sparge\_attn}, \texttt{fasa\_attn}), y con compuerta (\texttt{gla\_attn}, \texttt{deltanet\_attn}, \texttt{gated\_deltanet\_attn}, \texttt{gated\_deltanet2\_attn}, \texttt{hgrn2\_attn}, \texttt{fox\_attn}, \texttt{gated\_softmax\_attn}). \\
\texttt{ode\_solver} & enum & \texttt{rk4} o \texttt{euler}. \\
\texttt{ode\_steps} & int $\ge 1$ & Pasos de integración ODE. \\
\texttt{use\_bitnet} & bool & Habilitar ruta BitLinear de bajo bit. \\
\texttt{norm\_type} & enum & \texttt{layer\_norm}, \texttt{dynamic\_tanh}, \texttt{derf}. \\
\texttt{use\_factorized\_embedding} & bool & Habilitar embeddings factorizados. \\
\texttt{factorized\_embedding\_dim} & int $\ge 1$ & Dimensión reducida de embedding para factorización. \\
\texttt{use\_embedding\_conv} & bool & Habilitar Conv1d sobre el flujo de embeddings. \\
\texttt{embedding\_conv\_kernel} & int $\ge 1$ & Tamaño del kernel Conv1d. \\
\texttt{hope\_base} & float $\ge 0$ & Valor base de HoPE (opcional en el esquema). \\
\texttt{hope\_damping} & float $\ge 0$ & Amortiguamiento de HoPE (opcional en el esquema). \\
\texttt{use\_hope} & bool & Aplicar HoPE en \texttt{titan\_attn}. \\
\texttt{use\_moe} & bool & Habilitar ruta FFN con MoE. \\
\texttt{ffn\_hidden\_size} & int $\ge 1$ & Ancho intermedio de la FFN. \\
\texttt{ffn\_activation} & enum & \texttt{silu} o \texttt{gelu}. \\
\bottomrule
\end{longtable}
\normalsize

La profundidad en bucle inducida por el esquema es:
\[
L_{\text{logical}} = \texttt{num\_layers} \times \texttt{num\_loops}
\]
que es la definición a nivel de configuración de los bloques en bucle.

\subsection{Tipos de Tarea de Entrenamiento}

El campo \texttt{training.task} determina el objetivo de entrenamiento, trabajando en conjunto con \texttt{model.mode} para definir cómo aprende el modelo:

\paragraph{Modelado de Lenguaje Enmascarado (MLM).}
\begin{itemize}[leftmargin=*]
    \item \textbf{Modo}: Requiere \texttt{mode='encoder'} para atención bidireccional.
    \item \textbf{Objetivo}: Enmascarar aleatoriamente tokens en la secuencia de entrada (típicamente 15\%) y entrenar al modelo para predecir los tokens enmascarados basándose en el contexto bidireccional completo.
    \item \textbf{Caso de Uso}: Preentrenamiento de encoders para aprendizaje de representaciones, siguiendo la metodología BERT \citep{devlin_bert_2019}. El modelo aprende representaciones bidireccionales que capturan contexto tanto de izquierda como de derecha.
    \item \textbf{Configuración}: Utiliza el parámetro \texttt{mlm\_probability} para controlar la fracción de enmascaramiento.
\end{itemize}

\paragraph{Predicción Autorregresiva (AR) del Siguiente Token.}
\begin{itemize}[leftmargin=*]
    \item \textbf{Modo}: Requiere \texttt{mode='decoder'} para enmascaramiento causal.
    \item \textbf{Objetivo}: Entrenar al modelo para predecir el siguiente token en la secuencia dados solo los tokens anteriores.
    \item \textbf{Caso de Uso}: Generación de lenguaje y tareas estilo LLM siguiendo la metodología GPT. El modelo aprende dependencias secuenciales con atención causal donde cada token solo puede atender a tokens precedentes.
    \item \textbf{Clase de Modelo}: Típicamente usa \texttt{model\_class='frankesteindecoder'} para arquitecturas de decoder autorregresivo.
\end{itemize}

Este soporte de tareas dual permite la experimentación unificada tanto en preentrenamiento tipo encoder (MLM para comprensión bidireccional) como en generación tipo decoder (AR para producción secuencial de texto) dentro del mismo código base.

\subsection{Inventario Completo de Características de Entrenamiento}

\small
\begin{longtable}{p{5.1cm}p{2.4cm}p{7.0cm}}
\toprule
\textbf{Campo} & \textbf{Tipo/Rango} & \textbf{Significado} \\
\midrule
\endhead
\texttt{batch\_size} & int $\ge 1$ & Tamaño de lote del cargador. \\
\texttt{dataloader\_workers} & int $\ge 0$ & Trabajadores del dataloader de PyTorch. \\
\texttt{max\_length} & int $\ge 1$ & Límite de longitud de secuencia. \\
\texttt{task} & enum & Objetivo de entrenamiento: \texttt{mlm} (modelado de lenguaje enmascarado) o \texttt{sbert} (embedding de oraciones). \\
\texttt{mlm\_probability} & float $[0,1]$ & Probabilidad de enmascaramiento MLM (aplica solo cuando \texttt{task='mlm'}). \\
\texttt{max\_samples} & int $\ge 1$ & Máximo de muestras transmitidas. \\
\texttt{dataset\_batch\_size} & int $\ge 1$ & Tamaño interno de fragmento del conjunto de datos en streaming. \\
\texttt{num\_workers} & int $\ge 0$ & Trabajadores del conjunto de datos en streaming. \\
\texttt{cache\_dir} & string & Directorio de caché del conjunto de datos. \\
\texttt{local\_parquet\_dir} & string & Ruta opcional local a parquet. \\
\texttt{prefer\_local\_cache} & bool & Preferir caché local cuando esté disponible. \\
\texttt{stream\_local\_parquet} & bool & Modo de transmisión desde parquet local. \\
\texttt{use\_amp} & bool & Alternar precisión mixta. \\
\texttt{gradient\_accumulation\_steps} & int $\ge 1$ & Lote efectivo mediante acumulación. \\
\texttt{optimizer} & object & Contiene \texttt{optimizer\_class} y \texttt{parameters} con prefijo. \\
\texttt{scheduler\_total\_steps} & int $\ge 1$ & Horizonte del scheduler. \\
\texttt{scheduler\_warmup\_ratio} & float $[0,1]$ & Proporción de calentamiento. \\
\texttt{scheduler\_type} & enum & \texttt{cosine}, \texttt{constant}, \texttt{linear\_warmup\_then\_constant}. \\
\texttt{grad\_clip\_max\_norm} & float $\ge 0$ & Umbral de recorte de norma global. \\
\texttt{inf\_post\_clip\_threshold} & float $\ge 0$ & Guarda contra gradientes explosivos después del recorte. \\
\texttt{max\_nan\_retries} & int $\ge 0$ & Presupuesto de reintentos para inestabilidad NaN/Inf. \\
\texttt{checkpoint\_every\_n\_steps} & int $\ge 1$ & Frecuencia de checkpoint rodante. \\
\texttt{max\_rolling\_checkpoints} & int $\ge 1$ & Número de checkpoints rodantes a mantener. \\
\texttt{num\_best\_checkpoints} & int $\ge 1$ & Número de mejores checkpoints rastreados. \\
\texttt{nan\_check\_interval} & int $\ge 1$ & Cadencia de verificación NaN/Inf. \\
\texttt{log\_gradient\_stats} & bool & Habilitar registro de estadísticas de gradientes. \\
\texttt{gradient\_log\_interval} & int $\ge 1$ & Cadencia de registro de gradientes. \\
\texttt{csv\_log\_path} & string & Ruta de salida CSV a nivel de paso. \\
\texttt{csv\_rotate\_on\_schema\_change} & bool & Rotar CSV si el esquema de registro cambia. \\
\texttt{gpu\_metrics\_backend} & enum & \texttt{nvml} o \texttt{none}. \\
\texttt{nvml\_device\_index} & int $\ge 0$ & Índice de dispositivo para telemetría NVML. \\
\texttt{enable\_block\_grad\_norms} & bool & Incluir telemetría de norma de gradiente por bloque. \\
\texttt{telemetry\_log\_interval} & int $\ge 1$ & Intervalo de telemetría pesada (pasos de optimizador). \\
\texttt{use\_galore} & bool & Habilitar estrategia GaLore. \\
\texttt{galore\_rank} & int $\ge 1$ & Dimensión de proyección de bajo rango de GaLore. \\
\texttt{galore\_update\_interval} & int $\ge 1$ & Intervalo de actualización de proyección. \\
\texttt{galore\_scale} & float $\ge 0$ & Escalado de gradientes en el espacio proyectado. \\
\texttt{galore\_max\_dim} & int $\ge 1$ & Dimensión máxima del tensor para proyección GaLore. \\
\bottomrule
\end{longtable}
\normalsize

\subsection{Contrato de Prefijo del Optimizador (Completo)}

Los valores soportados de \texttt{optimizer\_class} son:
\texttt{sgd\_momentum}, \texttt{adamw}, \texttt{adafactor}, \texttt{galore\_adamw}, \texttt{prodigy}, \texttt{lion}, \texttt{sophia}, \texttt{muon}, \texttt{turbo\_muon}, \texttt{radam}, \texttt{adan}, \texttt{adopt}, \texttt{ademamix}, \texttt{mars\_adamw}, \texttt{cautious\_adamw}, \texttt{lamb}, \texttt{schedulefree\_adamw}, \texttt{shampoo}, \texttt{soap}, \texttt{apollo}, \texttt{apollo\_mini}, \texttt{q\_apollo}.

Las familias de sufijos compartidos por grupo (todos prefijados por el nombre del optimizador) son:
\begin{itemize}[leftmargin=*]
    \item Grupos de LR: \texttt{lr\_embeddings}, \texttt{lr\_norms}, \texttt{lr\_ode}, \texttt{lr\_retnet}, \texttt{lr\_mamba}, \texttt{lr\_attention}, \texttt{lr\_other}
    \item Grupos de decaimiento de peso: \texttt{wd\_embeddings}, \texttt{wd\_norms}, \texttt{wd\_ode}, \texttt{wd\_retnet}, \texttt{wd\_mamba}, \texttt{wd\_attention}, \texttt{wd\_other}
    \item Grupos de betas: \texttt{betas\_embeddings}, \texttt{betas\_norms}, \texttt{betas\_ode}, \texttt{betas\_retnet}, \texttt{betas\_mamba}, \texttt{betas\_attention}, \texttt{betas\_other}
    \item Grupos de épsilon: \texttt{eps\_embeddings}, \texttt{eps\_norms}, \texttt{eps\_ode}, \texttt{eps\_retnet}, \texttt{eps\_mamba}, \texttt{eps\_attention}, \texttt{eps\_other}
\end{itemize}

Sufijos globales específicos del optimizador:
\begin{itemize}[leftmargin=*]
    \item \texttt{sgd\_momentum}: \texttt{momentum}, \texttt{nesterov}
    \item \texttt{adafactor}: \texttt{beta2\_decay}, \texttt{clip\_threshold}, \texttt{eps1}, \texttt{eps2}
    \item \texttt{galore\_adamw}: \texttt{rank}, \texttt{update\_proj\_gap}
    \item \texttt{prodigy}: \texttt{d\_coef}
    \item \texttt{sophia}: \texttt{rho}, \texttt{update\_k}
    \item \texttt{muon} / \texttt{turbo\_muon}: \texttt{momentum}, \texttt{nesterov}, \texttt{ns\_steps}, \texttt{ns\_eps}
    \item \texttt{cautious\_adamw}: \texttt{cautious\_clip}
    \item \texttt{apollo}: \texttt{rank}, \texttt{update\_proj\_gap}, \texttt{scale}, \texttt{scale\_type}, \texttt{proj\_type}, \texttt{scale\_front}, \texttt{disable\_nl}
    \item \texttt{apollo\_mini}: \texttt{update\_proj\_gap}, \texttt{scale}, \texttt{proj\_type}, \texttt{scale\_front}, \texttt{disable\_nl}
    \item \texttt{q\_apollo}: \texttt{rank}, \texttt{update\_proj\_gap}, \texttt{scale}, \texttt{scale\_type}, \texttt{proj\_type}, \texttt{scale\_front}, \texttt{disable\_nl}, \texttt{quant\_bits}
\end{itemize}
Todas las demás clases en la lista anterior aceptan solo grupos compartidos con prefijo.

\subsection{Seguridad de Entrenamiento y Semántica de Ejecución}

Las características de seguridad a nivel de esquema incluyen acumulación, recorte, verificación de explosión post-recorte y reintentos por NaN:
\[
g_{\text{acc}}=\frac{1}{K}\sum_{i=1}^{K} g_i,\quad K=\texttt{gradient\_accumulation\_steps}
\]
\[
g_{\text{clip}} = g_{\text{acc}}\cdot
\min\left(1,\frac{\tau}{\|g_{\text{acc}}\|_2+\epsilon}\right),\quad
\tau=\texttt{grad\_clip\_max\_norm}
\]
luego los guardas de desbordamiento usan \texttt{inf\_post\_clip\_threshold} y lógica de reintentos limitada por \texttt{max\_nan\_retries}.

\begin{algorithm}[t]
\caption{Paso de Entrenamiento Guiado por Esquema con Controles de Estabilidad}
\begin{algorithmic}[1]
\Require Flujo de lotes, configuración $C$
\State Inicializar contador de reintentos $r\gets 0$
\For{cada paso del optimizador}
    \State Acumular gradientes para $K=C.\texttt{gradient\_accumulation\_steps}$ micro-lotes
    \State Aplicar recorte de norma global con $\tau=C.\texttt{grad\_clip\_max\_norm}$
    \If{el gradiente post-recorte excede $C.\texttt{inf\_post\_clip\_threshold}$ o se detecta NaN/Inf}
        \If{$r < C.\texttt{max\_nan\_retries}$}
            \State restaurar estado seguro / saltar paso; $r \gets r+1$
            \State \textbf{continue}
        \Else
            \State detener entrenamiento con estado de fallo
        \EndIf
    \EndIf
    \State ejecutar paso del optimizador seleccionado por \texttt{optimizer\_class}
    \State actualizar scheduler (\texttt{cosine}, \texttt{constant}, o \texttt{linear\_warmup\_then\_constant})
    \If{paso mod \texttt{checkpoint\_every\_n\_steps}$=0$}
        \State guardar checkpoint rodante y podar a \texttt{max\_rolling\_checkpoints}
    \EndIf
    \State actualizar mejores checkpoints hasta \texttt{num\_best\_checkpoints}
    \State emitir CSV + telemetría según \texttt{gradient\_log\_interval} y \texttt{telemetry\_log\_interval}
\EndFor
\end{algorithmic}
\end{algorithm}

\subsection{Variantes de Normalización: RMSNorm, Tanh Dinámico y Erf Dinámico}
\label{sec:normalization}

La normalización determina cómo se controla la escala de activación a través de la profundidad. En este repositorio, la normalización no es solo una elección de modelado sino también una cuestión de compatibilidad con el esquema, porque solo ciertos valores son actualmente aceptados por \texttt{norm\_type}. Las tres formulaciones más relevantes para este código base son:

\subsection{RMSNorm}
RMSNorm elimina la centralización de la media y solo reescala por la magnitud de la raíz cuadrada media \citep{zhang_root_2019}:
\[
\mathrm{RMS}(x)=\sqrt{\frac{1}{d}\sum_{i=1}^{d}x_i^2+\epsilon},\qquad
y_i=\gamma_i\frac{x_i}{\mathrm{RMS}(x)}
\]
En comparación con LayerNorm, RMSNorm es computacionalmente más simple (sin resta de la media de características) y se usa a menudo cuando es importante reducir la sobrecarga de normalización.

\subsection{Tanh Dinámico (DyT)}
El Tanh Dinámico propone reemplazar la normalización explícita con un mapeo elemento a variable acotado \citep{zhu_transformers_2025}:
\[
\mathrm{DyT}(x)=\tanh(\alpha x)
\]
donde $\alpha$ se aprende. La idea central es que la contracción no lineal acotada puede proporcionar un escalado de señal estable sin calcular explícitamente estadísticas de normalización por token.

\subsection{Erf Dinámico (Derf)}
Derf extiende la misma dirección libre de normalización usando un mapeo basado en la función error \citep{chen_stronger_2025}:
\[
\mathrm{Derf}(x)=\mathrm{erf}(\alpha x+s)
\]
con escala/desplazamiento aprendibles. Los resultados reportados en el trabajo citado indican un rendimiento superior al de DyT y las líneas base de normalización comunes en múltiples dominios.

\subsection{Implicaciones del Esquema}
El contrato de configuración actual en este repositorio permite:
\[
\texttt{norm\_type}\in\{\texttt{layer\_norm},\texttt{dynamic\_tanh},\texttt{derf}\}
\]
por lo que DyT y Derf están directamente disponibles en ejecuciones guiadas por esquema, mientras que RMSNorm no es actualmente un valor de enumeración aceptado y requeriría extensión de código/esquema.

\scriptsize
\begin{longtable}{p{2.2cm}p{3.4cm}p{2.0cm}p{4.0cm}}
\toprule
\textbf{Método} & \textbf{Fórmula} & \textbf{Estadísticas Necesarias} & \textbf{Notas} \\
\midrule
\endhead
RMSNorm & $\gamma_i x_i/\sqrt{\frac{1}{d}\sum_j x_j^2+\epsilon}$ & Solo RMS & Menor sobrecarga; línea base ampliamente usada \citep{zhang_root_2019}. \\
Tanh Dinámico & $\tanh(\alpha x)$ & ninguna & Transformación acotada libre de normalización; reemplazo directo \citep{zhu_transformers_2025}. \\
Erf Dinámico & $\mathrm{erf}(\alpha x+s)$ & ninguna & Alternativa libre de normalización; mejora sobre DyT \citep{chen_stronger_2025}. \\
\bottomrule
\end{longtable}
\normalsize

\section{Taxonomía de Arquitectura e Implementación}
\label{sec:architecture-taxonomy}

\subsection{Familias de Atención y Mezcladores de Secuencias}
\label{sec:attention-families}

Este sistema implementa diecisiete arquitecturas de mezcladores de secuencias distintas organizadas en cinco categorías funcionales que reflejan las tendencias de investigación en el diseño de modelado de secuencias. La organización taxonómica refleja la comprensión evolutiva de cómo equilibrar expresividad, eficiencia computacional y restricciones de memoria.

\begin{enumerate}[leftmargin=*]
    \item \textbf{Líneas Base de Atención Densas}: La atención softmax estándar y la atención sigmoide proporcionan contextualización global completa a costo computacional cuadrático, sirviendo como líneas base de referencia para comparación con alternativas más eficientes.
    \item \textbf{Arquitecturas Recurrentes y Retentivas}: RetNet, Mamba, bloques estilo ODE y Titans mantienen representaciones de estado que permiten un costo de inferencia $\mathcal{O}(1)$ mientras preservan la expresividad mediante dinámicas recurrentes, parámetros selectivos o adaptación de memoria en tiempo de prueba.
    \item \textbf{Patrones de Atención Dispersa}: Siete variantes dispersas (Sparse Transformer, Longformer, BigBird, SparseK, NSA, SpargeAttn, FASA) reducen la complejidad cuadrática mediante dispersión estructurada, selección de tokens o estrategias de poda sin entrenamiento.
    \item \textbf{Mecanismos de Memoria con Compuerta}: Siete arquitecturas con compuerta (GLA, DeltaNet, Gated DeltaNet, Gated DeltaNet-2, HGRN2, FoX, Gated Softmax) introducen control dependiente de datos sobre la retención de memoria, el olvido y la fuerza de actualización.
\end{enumerate}

\begin{figure}[t]
\centering
\fitfigure{\begin{tikzpicture}[
    font=\scriptsize,
    node distance=0.38cm and 1.0cm,
    box/.append style={inner sep=4pt, minimum height=0.55cm}
]
% Root
\node[box, fill=blue!15, minimum width=3.2cm] (root) {\textbf{Sequence Mixer Registry (17 variants)}};

% Category headers
\node[box, fill=red!10, below left=1.0cm and 5.4cm of root, minimum width=2.2cm] (dense) {\textbf{Dense (2)}};
\node[box, fill=green!10, below left=1.0cm and 1.8cm of root, minimum width=2.2cm] (recur) {\textbf{Recurrent (4)}};
\node[box, fill=yellow!12, below right=1.0cm and 1.8cm of root, minimum width=2.2cm] (sparse) {\textbf{Sparse (7)}};
\node[box, fill=purple!12, below right=1.0cm and 5.4cm of root, minimum width=2.2cm] (gated) {\textbf{Gated (7)}};

% Dense column
\node[box, fill=red!5, below=of dense] (d1) {\texttt{standard\_attn}};
\node[box, fill=red!5, below=of d1] (d2) {\texttt{sigmoid\_attn}};

% Recurrent column
\node[box, fill=green!5, below=of recur] (r1) {\texttt{retnet / retnet\_attn}};
\node[box, fill=green!5, below=of r1] (r2) {\texttt{mamba}};
\node[box, fill=green!5, below=of r2] (r3) {\texttt{ode}};
\node[box, fill=green!5, below=of r3] (r4) {\texttt{titan\_attn}};

% Sparse column
\node[box, fill=yellow!6, below=of sparse] (s1) {\texttt{sparse\_transformer\_attn}};
\node[box, fill=yellow!6, below=of s1] (s2) {\texttt{longformer\_attn}};
\node[box, fill=yellow!6, below=of s2] (s3) {\texttt{bigbird\_attn}};
\node[box, fill=yellow!6, below=of s3] (s4) {\texttt{sparsek\_attn}};
\node[box, fill=yellow!6, below=of s4] (s5) {\texttt{nsa\_attn}};
\node[box, fill=yellow!6, below=of s5] (s6) {\texttt{sparge\_attn}};
\node[box, fill=yellow!6, below=of s6] (s7) {\texttt{fasa\_attn}};

% Gated column
\node[box, fill=purple!6, below=of gated] (g1) {\texttt{gla\_attn}};
\node[box, fill=purple!6, below=of g1] (g2) {\texttt{deltanet\_attn}};
\node[box, fill=purple!6, below=of g2] (g3) {\texttt{gated\_deltanet\_attn}};
\node[box, fill=purple!6, below=of g3] (g4) {\texttt{gated\_deltanet2\_attn}};
\node[box, fill=purple!6, below=of g4] (g5) {\texttt{hgrn2\_attn}};
\node[box, fill=purple!6, below=of g5] (g6) {\texttt{fox\_attn}};
\node[box, fill=purple!6, below=of g6] (g7) {\texttt{gated\_softmax\_attn}};

% Root-to-category links
\draw[line] (root) -- (dense);
\draw[line] (root) -- (recur);
\draw[line] (root) -- (sparse);
\draw[line] (root) -- (gated);

% Category-to-list links
\draw[line] (dense) -- (d1);
\draw[line] (d1) -- (d2);

\draw[line] (recur) -- (r1);
\draw[line] (r1) -- (r2);
\draw[line] (r2) -- (r3);
\draw[line] (r3) -- (r4);

\draw[line] (sparse) -- (s1);
\draw[line] (s1) -- (s2);
\draw[line] (s2) -- (s3);
\draw[line] (s3) -- (s4);
\draw[line] (s4) -- (s5);
\draw[line] (s5) -- (s6);
\draw[line] (s6) -- (s7);

\draw[line] (gated) -- (g1);
\draw[line] (g1) -- (g2);
\draw[line] (g2) -- (g3);
\draw[line] (g3) -- (g4);
\draw[line] (g4) -- (g5);
\draw[line] (g5) -- (g6);
\draw[line] (g6) -- (g7);
\end{tikzpicture}}
\caption{Taxonomía completa de las diecisiete arquitecturas de mezcladores de secuencias soportadas. Las líneas base densas proporcionan enrutamiento global completo a costo cuadrático. Las arquitecturas recurrentes permiten inferencia en tiempo constante mediante compresión de estado. Las variantes dispersas reducen la complejidad mediante patrones estructurados o selección de tokens. Los mecanismos con compuerta introducen control dependiente de datos sobre la retención y el olvido de memoria.}
\label{fig:comprehensive-taxonomy}
\end{figure}

\subsection{Atención Estándar}
Dadas las matrices proyectadas $(Q,K,V)$:
\[
\text{Attn}(Q,K,V)=\text{softmax}\left(\frac{QK^\top}{\sqrt{d_k}}\right)V
\]
Este es el mecanismo base para el enrutamiento de contenido \citep{vaswani_attention_2017}.

\subsection{Atención Sigmoide}
La atención sigmoide elimina la normalización de probabilidad por fila:
\[
\text{SigmoidAttn}(Q,K,V)=\sigma\left(\frac{QK^\top}{\sqrt{d_k}}+b\right)V
\]
y tiene diferentes requisitos de estabilidad de entrenamiento, a menudo con normalización adicional \citep{ramapuram_theory_2024}.

\subsection{Formulación Retentiva}
RetNet usa retención con matriz de decaimiento $D$:
\[
\text{Retention}(Q,K,V)=\left(QK^\top \odot D\right)V
\]
con forma recurrente:
\[
S_n=\gamma S_{n-1}+k_n^\top v_n,\qquad o_n=q_n S_n
\]
permitiendo inferencia recurrente de bajo costo \citep{sun_retentive_2023,yang_survey_2025}.

\subsection{SSM Selectivo (Mamba)}
La recurrencia discreta selectiva de espacio de estados es:
\[
h_t=\bar{A}_t h_{t-1}+\bar{B}_t x_t,\qquad y_t=C_t h_t
\]
donde $(\bar{A}_t,\bar{B}_t,C_t)$ dependen de la entrada, preservando escalado en tiempo lineal con escaneo consciente del hardware \citep{gu_mamba_2023,hwang_hydra_2024}.

\subsection{Actualizaciones Continuas estilo ODE}
El marco de profundidad continua:
\[
\frac{dh(t)}{dt}=f_\theta(h(t),t)
\]
con integradores RK prácticos para ejecución discreta \citep{zhang_continuous_2021}.

\subsection{Memoria en Tiempo de Prueba (Titans)}
Una actualización aumentada con memoria puede escribirse:
\[
M_t=(1-\alpha_t)M_{t-1}+S_t,\qquad
S_t=\eta_t S_{t-1}-\theta_t\nabla \ell(M_{t-1};x_t)
\]
para adaptar la memoria en tiempo de inferencia \citep{behrouz_titans_2025,behrouz_titans_2024}.

\subsection{Extensiones Dispersas y con Compuerta en el Código Base Actual}
El registro de mezcladores ahora incluye bloques dispersos:
\texttt{sparse\_transformer\_attn}, \texttt{longformer\_attn}, \texttt{bigbird\_attn},
\texttt{sparsek\_attn}, \texttt{nsa\_attn}, \texttt{sparge\_attn}, y \texttt{fasa\_attn};
y bloques con compuerta:
\texttt{gla\_attn}, \texttt{deltanet\_attn}, \texttt{gated\_deltanet\_attn},
\texttt{gated\_deltanet2\_attn}, \texttt{hgrn2\_attn}, \texttt{fox\_attn}, y \texttt{gated\_softmax\_attn}.

La implementación impone una política de ejecución explícita para métodos dispersos sin entrenamiento:
\texttt{fasa\_attn} y \texttt{sparge\_attn} son solo para evaluación/inferencia y generan errores en tiempo de ejecución si se usan mientras el modelo está en modo de entrenamiento.

\begin{figure}[t]
\centering
\fitfigure{\begin{tikzpicture}[node distance=0.9cm and 0.9cm]
\node[box, fill=gray!10] (tokens) {Token sequence};
\node[box, fill=blue!8, below left=1.1cm and 2.5cm of tokens] (local) {Local / windowed\\Longformer};
\node[box, fill=yellow!10, below=1.1cm of tokens] (hybrid) {Hybrid sparse graph\\BigBird / Sparse Transformer};
\node[box, fill=orange!12, below right=1.1cm and 2.5cm of tokens] (select) {Selective pruning\\SparseK / NSA / FASA / SpargeAttn};
\draw[line] (tokens) -- (local);
\draw[line] (tokens) -- (hybrid);
\draw[line] (tokens) -- (select);
\node[group, fit=(local)(hybrid)(select), label=below:{\small Three sparse design strategies: locality, structured graph sparsity, and learned or predicted selection}] {};
\end{tikzpicture}}
\caption{Mapa conceptual de las estrategias de diseño de atención dispersa utilizadas en el código base. Diferentes métodos reducen el costo restringiendo vecindarios, construyendo grafos dispersos o seleccionando solo tokens/bloques de alto valor.}
\label{fig:sparse-map}
\end{figure}

\subsection{Bloques de Atención Dispersa Implementados (Detallado)}
Este código base incluye siete familias de atención dispersa \citep{child_sparse_transformer_2019,beltagy_longformer_2020,zaheer_bigbird_2020,lou_sparsek_2024,yuan_nsa_2025,zhang_spargeattn_2025,wang_fasa_2026}.

\paragraph{Sparse Transformer (\texttt{sparse\_transformer\_attn}).}
Utiliza máscaras dispersas factorizadas (con zancada + fijas) para aproximar la conectividad densa a menor costo que la atención completa $\mathcal{O}(n^2)$:
\[
\text{Attn}_i=\text{softmax}\!\left(\frac{q_iK_{A_i}^\top}{\sqrt{d_k}}\right)V_{A_i}
\]
donde $A_i$ es la vecindad dispersa inducida por reglas de zancada/fijas.
\citep{child_sparse_transformer_2019}

\paragraph{Longformer (\texttt{longformer\_attn}).}
Utiliza localidad de ventana deslizante con tokens globales opcionales:
\[
A_i=\{j:\lvert i-j\rvert \le w/2\}\cup\mathcal{G}
\]
produciendo escalado lineal en longitud de secuencia para ventana fija $w$.
\citep{beltagy_longformer_2020}

\paragraph{BigBird (\texttt{bigbird\_attn}).}
Combina ventanas locales, enlaces aleatorios y tokens globales:
\[
A_i=A_i^{\text{window}}\cup A_i^{\text{random}}\cup A_i^{\text{global}}
\]
para preservar una fuerte conectividad de contexto largo con cómputo disperso.
\citep{zaheer_bigbird_2020}

\paragraph{SparseK (\texttt{sparsek\_attn}).}
Utiliza una proyección diferenciable tipo top-$k$ sobre puntuaciones de importancia antes de la atención, de modo que solo los pares KV seleccionados participan en la ruta costosa de producto punto.
\citep{lou_sparsek_2024}

\paragraph{NSA (\texttt{nsa\_attn}).}
Implementa un diseño disperso de tres ramas: rama comprimida, rama seleccionada y rama de ventana local, y luego las combina con compuertas aprendidas:
\[
o_t=\sum_{c\in\{\text{cmp},\text{sel},\text{win}\}} g_t^c\,\text{Attn}(q_t,\tilde K_t^c,\tilde V_t^c)
\]
\citep{yuan_nsa_2025}

\paragraph{SpargeAttn (\texttt{sparge\_attn}).}
Filtrado de bloques en dos etapas sin entrenamiento: primero predice interacciones de bloques insignificantes, luego aplica poda consciente de softmax para eliminar bloques de baja contribución.
\citep{zhang_spargeattn_2025}

\paragraph{FASA (\texttt{fasa\_attn}).}
Atención sin entrenamiento consciente de frecuencia: utiliza fragmentos de frecuencia RoPE dominantes para la predicción de importancia de tokens, luego aplica atención completa solo en los tokens seleccionados.
\citep{wang_fasa_2026}

\small
\begin{longtable}{p{2.5cm}p{1.6cm}p{1.9cm}p{2.4cm}p{4.1cm}p{1.5cm}}
\toprule
\textbf{Bloque} & \textbf{Entrenable} & \textbf{Tendencia Asintótica} & \textbf{Unidad de Dispersión Principal} & \textbf{Notas de Integración Actual} & \textbf{Ref} \\
\midrule
\endhead
Sparse Transformer & Sí & sub-cuadrática & patrón de máscara (a nivel de token) & Máscaras factorizadas con zancada/fijas dentro del pipeline SDPA. & \citep{child_sparse_transformer_2019} \\
Longformer & Sí & lineal en $n$ ($w$ fijo) & ventana deslizante + tokens globales & Máscara de ventana con índices globales opcionales. & \citep{beltagy_longformer_2020} \\
BigBird & Sí & casi lineal & bordes ventana + aleatorios + globales & Máscara dispersa aleatorizada más rutas local/global. & \citep{zaheer_bigbird_2020} \\
SparseK & Sí & similar a lineal (KV seleccionados) & selección diferenciable top-$k$ de KV & Red de puntuación aprendida + proyección SparseK + atención con KV agrupados. & \citep{lou_sparsek_2024} \\
NSA & Sí & multi-rama con tokens reducidos & bloques comprimidos + bloques seleccionados + ventana local & Tres ramas dispersas combinadas con compuerta en un tensor de salida. & \citep{yuan_nsa_2025} \\
SpargeAttn & No (sin entrenamiento) & dependiente del bloque disperso & dispersión de bloque predecida a nivel de bloque & Solo evaluación en este repositorio; genera error en modo entrenamiento. & \citep{zhang_spargeattn_2025} \\
FASA & No (sin entrenamiento) & dependiente del token seleccionado & fragmentos de frecuencia dominantes + tokens seleccionados & Solo evaluación en este repositorio; genera error en modo entrenamiento. & \citep{wang_fasa_2026} \\
\bottomrule
\end{longtable}
\normalsize

\subsection{Bloques de Atención con Compuerta Implementados (Detallado)}
Este código base incluye ocho bloques con compuerta \citep{yang_gla_2023,yang_deltanet_2024,yang_gated_deltanet_2024,hatamizadeh_gated_deltanet2_2026,sun_retentive_2023,qin_hgrn2_2024,lin_forgetting_transformer_2025,qiu_gated_attention_2025}.

La idea unificadora es que las compuertas controlan \emph{qué información sobrevive}. Algunas compuertas actúan sobre actualizaciones de estado recurrente (GLA, variantes de DeltaNet, HGRN2), mientras que otras modifican la ruta de atención completa (FoX y Gated Softmax). Esto hace que las compuertas sean especialmente útiles cuando el modelo debe equilibrar recuerdo, actualidad y memoria limitada.

\begin{figure}[t]
\centering
\fitfigure{\begin{tikzpicture}[node distance=1.2cm and 1.0cm]
\node[box, fill=green!8] (state) {Previous state\\$S_{t-1}$};
\node[box, fill=red!8, right=of state] (gate) {Gate(s)\\$\alpha_t,\beta_t,G_t,f_t$};
\node[box, fill=blue!8, right=of gate] (write) {New key/value\\or SDPA output};
\node[box, fill=yellow!12, below=of gate] (next) {Updated memory / output\\$S_t$ or $O_t$};
\draw[line] (state) -- (gate);
\draw[line] (gate) -- (write);
\draw[line] (state) |- (next);
\draw[line] (write) -- (next);
\end{tikzpicture}}
\caption{Plantilla de compuerta genérica. Una compuerta puede atenuar la memoria existente, regular la fuerza de escritura o modular las salidas de atención densa, dependiendo de la familia del bloque.}
\label{fig:gating-template}
\end{figure}

\paragraph{GLA (\texttt{gla\_attn}).}
La Atención Lineal con Compuerta aplica decaimiento multiplicativo dependiente de datos en las actualizaciones de estado recurrente:
\[
S_t=G_t\odot S_{t-1}+v_tk_t^\top,\qquad o_t=S_tq_t
\]
para controlar la acumulación de memoria.
\citep{yang_gla_2023}

\paragraph{DeltaNet (\texttt{deltanet\_attn}).}
Utiliza una escritura correctora de error tipo delta-regla con fuerza de escritura aprendida $\beta_t$:
\[
S_t=S_{t-1}(I-\beta_tk_tk_t^\top)+\beta_tv_tk_t^\top
\]
lo que mejora el reemplazo de memoria dirigido.
\citep{yang_deltanet_2024}

\paragraph{Gated DeltaNet (\texttt{gated\_deltanet\_attn}).}
Añade compuerta de decaimiento $\alpha_t$ sobre las escrituras de delta-regla:
\[
S_t=\alpha_t S_{t-1}(I-\beta_tk_tk_t^\top)+\beta_tv_tk_t^\top
\]
tanto para olvido global como para actualizaciones correctivas locales.
\citep{yang_gated_deltanet_2024}

\paragraph{Gated DeltaNet-2 (\texttt{gated\_deltanet2\_attn}).}
Desacopla la compuerta delta escalar $\beta_t$ en una compuerta de borrado canalizada $\mathbf{b}_t$ (eje de claves) y una compuerta de escritura canalizada $\mathbf{w}_t$ (eje de valores), con decaimiento canalizado $\boldsymbol{\alpha}_t$ (estilo KDA):
\[
S_t=(I-k_t(\mathbf{b}_t\odot k_t)^\top)\,\mathrm{Diag}(\boldsymbol{\alpha}_t)\,S_{t-1}+k_t(\mathbf{w}_t\odot v_t)^\top
\]
para control de memoria más fino; se reduce a KDA y Gated DeltaNet como casos especiales.
\citep{hatamizadeh_gated_deltanet2_2026}

\paragraph{Alias RetNet Attn (\texttt{retnet\_attn}).}
Proporciona un alias explícito con compuerta que envuelve el comportamiento de retención multiescala para consistencia de nomenclatura en los registros de capas.
\citep{sun_retentive_2023}

\paragraph{HGRN2 (\texttt{hgrn2\_attn}).}
Utiliza compuertas de olvido con cota inferior y expansión de estado con producto exterior:
\[
S_t=\text{diag}(g_t)S_{t-1}+v_tk_t^\top
\]
para aumentar la expresividad del estado recurrente mientras se mantiene la eficiencia.
\citep{qin_hgrn2_2024}

\paragraph{FoX (\texttt{fox\_attn}).}
Inyecta sesgo de olvido a nivel de token directamente en los logits de softmax:
\[
O=\text{softmax}(QK^\top + D)V
\]
donde $D$ se deriva de compuertas acumulativas de log-olvido.
\citep{lin_forgetting_transformer_2025}

\paragraph{Gated Softmax (\texttt{gated\_softmax\_attn}).}
Aplica una compuerta sigmoide post-SDPA:
\[
Y'=\text{SDPA}(Q,K,V)\odot \sigma(XW_g)
\]
lo que añade compuerta de canal multiplicativa sin reemplazar la atención softmax.
\citep{qiu_gated_attention_2025}

\small
\begin{longtable}{p{2.4cm}p{2.0cm}p{2.2cm}p{1.8cm}p{4.0cm}p{1.5cm}}
\toprule
\textbf{Bloque} & \textbf{Tipo de Estado} & \textbf{Mecanismo de Compuerta} & \textbf{Ruta Softmax} & \textbf{Notas de Integración Actual} & \textbf{Ref} \\
\midrule
\endhead
GLA & estado recurrente matricial & decaimiento multiplicativo dependiente de datos & No (recurrente lineal) & Actualización recurrente con proyección de compuerta de bajo rango. & \citep{yang_gla_2023} \\
DeltaNet & estado recurrente matricial & compuerta de escritura ($\beta$) & No (recurrente lineal) & Corrección delta-regla con Q/K normalizados. & \citep{yang_deltanet_2024} \\
Gated DeltaNet & estado recurrente matricial & compuertas de decaimiento + escritura ($\alpha,\beta$) & No (recurrente lineal) & Olvido combinado y escritura dirigida. & \citep{yang_gated_deltanet_2024} \\
Gated DeltaNet-2 & estado recurrente matricial & compuertas de borrado + escritura decoupladas (canalizadas) & No (recurrente lineal) & Borrado canalizado (claves) + escritura (valores); decaimiento estilo KDA. & \citep{hatamizadeh_gated_deltanet2_2026} \\
RetNet Attn & estado recurrente matricial & decaimiento multiescala fijo & No (retención) & Envoltorio alias sobre el mezclador RetNet existente. & \citep{sun_retentive_2023} \\
HGRN2 & estado recurrente matricial & compuerta de olvido con cota inferior & No (recurrente lineal) & Compuerta recurrente jerárquica con productos exteriores. & \citep{qin_hgrn2_2024} \\
FoX & matriz de atención completa & compuerta de olvido en espacio de logits & Sí & Sesgo de olvido añadido antes de softmax. & \citep{lin_forgetting_transformer_2025} \\
Gated Softmax & matriz de atención completa & compuerta sigmoide post-atención & Sí & Compuerta sigmoide aplicada después de la salida SDPA. & \citep{qiu_gated_attention_2025} \\
\bottomrule
\end{longtable}
\normalsize

\begin{algorithm}[t]
\caption{Paso Adelante del Mezclador Guiado por Patrón (Conceptual)}
\begin{algorithmic}[1]
\Require Estados ocultos $H$, patrón $P$, índice de capa $\ell$
\State $m \gets P[\ell \bmod |P|]$
\If{$m \in \{\texttt{fasa\_attn}, \texttt{sparge\_attn}\}$ y el modelo está en modo entrenamiento}
    \State lanzar error de configuración/tiempo de ejecución (bloque sin entrenamiento en modo entrenar)
\ElsIf{$m=\texttt{standard\_attn}$}
    \State $H \gets \text{softmax-attention}(H)$
\ElsIf{$m=\texttt{sigmoid\_attn}$}
    \State $H \gets \text{sigmoid-attention}(H)$
\ElsIf{$m\in\{\texttt{retnet},\texttt{retnet\_attn}\}$}
    \State $H \gets \text{retention}(H)$
\ElsIf{$m=\texttt{mamba}$}
    \State $H \gets \text{selective-ssm}(H)$
\ElsIf{$m=\texttt{ode}$}
    \State $H \gets \text{rk-step}(H)$
\ElsIf{$m$ es una clave de atención dispersa}
    \State $H \gets \text{sparse-attention-family}(H)$
\ElsIf{$m$ es una clave de atención con compuerta}
    \State $H \gets \text{gated-attention-family}(H)$
\Else
    \State $H \gets \text{memory-augmented-attn}(H)$
\EndIf
\State \Return $H$
\end{algorithmic}
\end{algorithm}

\section{Familias de Optimizadores y Dinámicas de Entrenamiento}
\label{sec:optimizer-families}

La optimización de arquitecturas transformer altamente parametrizadas presenta desafíos significativos debido a paisajes de pérdida no convexos, puntos de silla y heterogeneidad de bloques entre grupos de parámetros. Este sistema aborda estos desafíos mediante un marco unificado que soporta veintitrés familias de optimizadores que abarcan seis categorías algorítmicas: (1) líneas base clásicas (SGD, AdamW), (2) momento avanzado y reducción de varianza (Adan, ADOPT, AdEMAMix, MARS, Cautious), (3) variantes eficientes en memoria (Adafactor, GaLore, Lion, APOLLO, APOLLO-Mini, Q-APOLLO), (4) métodos sin programación y sin parámetros (Schedule-Free AdamW, Prodigy) más escalado de lotes grandes mediante LAMB, (5) conscientes de curvatura y segundo orden (Sophia, Shampoo, SOAP), (6) orientados a geometría (Muon, Turbo-Muon), y (7) optimizadores con adaptabilidad ajustable (Anon). Las descripciones algorítmicas detalladas, pseudocódigo y una tabla comparativa de memoria/complejidad para todos los optimizadores se proporcionan en el Apéndice~\ref{sec:annex-optimizers}.

\begin{figure}[t]
\centering
\fitfigure{\begin{tikzpicture}[node distance=1.1cm and 1.0cm]
\node[box, fill=blue!8] (start) {Elegir objetivo de optimizador};
\node[box, fill=green!8, below left=of start, xshift=-1.5cm] (baseline) {Línea base confiable\\AdamW / RAdam};
\node[box, fill=yellow!12, below=of start] (memory) {Menor memoria de optimizador\\Adafactor / GaLore / Lion};
\node[box, fill=orange!12, below right=of start, xshift=1.5cm] (advanced) {Agresivo o estructurado\\Adan, ADOPT, Sophia,\\Shampoo, SOAP, Muon};
\node[box, fill=red!8, below=of memory] (route) {Grupos con prefijo de esquema\\LR, decay de peso, betas, eps};
\draw[line] (start) -- (baseline);
\draw[line] (start) -- (memory);
\draw[line] (start) -- (advanced);
\draw[line] (baseline) -- (route);
\draw[line] (memory) -- (route);
\draw[line] (advanced) -- (route);
\end{tikzpicture}}
\caption{Selección de optimizador en la práctica: la clase determina la regla de actualización, mientras que el esquema controla cómo se aplican los hiperparámetros entre embeddings, normas, bloques recurrentes, bloques de atención y otros parámetros.}
\label{fig:optimizer-routing}
\end{figure}

\section{Cuantización y Despliegue}
\label{sec:quantization}

El stack de despliegue utiliza empaquetado ternario de pesos más cuantización de activaciones INT8 para producir artefactos eficientes.

\begin{figure}[t]
\centering
\fitfigure{\begin{tikzpicture}[node distance=1.3cm and 0.9cm]
\node[box, fill=blue!8] (ckpt) {Checkpoint entrenado};
\node[box, fill=yellow!10, right=of ckpt] (quant) {Empaquetado de pesos\\ternario / bajo bit};
\node[box, fill=orange!12, right=of quant] (act) {Escalado de activaciones\\ruta INT8};
\node[box, fill=green!8, right=of act] (artifact) {Artefacto desplegable\\huella de almacenamiento reducida};
\draw[line] (ckpt) -- (quant);
\draw[line] (quant) -- (act);
\draw[line] (act) -- (artifact);
\end{tikzpicture}}
\caption{Ruta de despliegue desde un checkpoint entrenado hasta un artefacto compacto. El código base trata la cuantización como una transformación de la etapa de despliegue en lugar de una familia de modelos separada.}
\label{fig:deploy-flow}
\end{figure}

\subsection{Cuantización Ternaria}
Dado el tensor de pesos $W$, un escalado práctico es:
\[
s=\text{mean}(|W|),\qquad
\tilde{W}=\text{clip}\left(\text{round}\left(\frac{W}{s}\right),-1,1\right)
\]
que aproxima las actualizaciones de bajo bit estilo BitNet \citep{wang_bitnet_2023}. El mapeo empaquetado utiliza dos bits por símbolo de peso para eficiencia de almacenamiento.

\subsection{Cuantización de Activaciones}
Para las activaciones $x$:
\[
q=\text{round}\left(x\cdot\frac{127}{\max(|x|)+\epsilon}\right),\qquad q\in[-128,127]
\]
con des-cuantización $x\approx q/\alpha$.

\subsection{Estimaciones de Tamaño}
Para $N$ parámetros:
\[
\text{Tamaño FP32}\approx 4N,\quad \text{Tamaño FP16}\approx 2N,\quad \text{Tamaño 1.58-bit}\approx \frac{1.58}{8}N
\]
antes de metadatos y sobrecarga de empaquetado. Esto se alinea con objetivos de despliegue ligero \citep{samson_lightweight_2026,bravin_embbert_2026}.

\section{Tareas Posteriores de SBERT}
\label{sec:sbert}

El embedding de oraciones se construye sobre entrenamiento estilo Siamese \citep{reimers_sentence-bert_2019}. Para el par de oraciones $(s_1,s_2)$ con embeddings $(e_1,e_2)$:
\[
\text{cos}(e_1,e_2)=\frac{e_1^\top e_2}{\|e_1\|\|e_2\|}
\]
y pérdida coseno estilo regresión:
\[
\mathcal{L}_{\text{cos}}=\left(\text{cos}(e_1,e_2)-y\right)^2
\]
con $y\in[-1,1]$ en este pipeline.

Modos posteriores soportados:
\begin{itemize}[leftmargin=*]
    \item \textbf{Similitud}: puntuación por pares entre dos oraciones.
    \item \textbf{Búsqueda}: $k$ vecinos más cercanos sobre un corpus.
    \item \textbf{Clustering}: agrupación de embeddings (ej., k-means).
    \item \textbf{Codificación}: exportación persistente de embeddings para recuperación posterior.
\end{itemize}

\begin{figure}[t]
\centering
\fitfigure{\begin{tikzpicture}[node distance=1.1cm and 1.0cm]
\node[box, fill=blue!8] (encoder) {Codificador compartido};
\node[box, fill=yellow!10, below left=of encoder, xshift=-1.2cm] (sim) {Similitud\\puntuación coseno};
\node[box, fill=green!8, below=of encoder] (search) {Búsqueda\\recuperación top-$k$};
\node[box, fill=orange!12, below right=of encoder, xshift=1.2cm] (cluster) {Cluster / codificar\\análisis offline};
\draw[line] (encoder) -- (sim);
\draw[line] (encoder) -- (search);
\draw[line] (encoder) -- (cluster);
\end{tikzpicture}}
\caption{Reutilización del flujo de trabajo SBERT. Un único codificador soporta puntuación de pares en línea, recuperación sobre corpus, clustering y exportación persistente de embeddings.}
\label{fig:sbert-modes}
\end{figure}

\begin{algorithm}[t]
\caption{Enrutador de Modo de Inferencia SBERT}
\begin{algorithmic}[1]
\Require modo $m$, modelo $E$, entradas $X$
\If{$m=\texttt{similarity}$}
    \State retornar $\cos(E(x_1),E(x_2))$
\ElsIf{$m=\texttt{search}$}
    \State retornar top-$k$ por producto-punto/coseno contra embeddings del corpus
\ElsIf{$m=\texttt{cluster}$}
    \State retornar etiquetas de clustering sobre $E(X)$
\Else
    \State retornar embeddings serializados $E(X)$
\EndIf
\end{algorithmic}
\end{algorithm}

\section{Tablas Resumen}
\label{sec:summary-tables}

\subsection{Resumen de Atención y Mezcladores de Secuencia}
\small
\begin{longtable}{p{2.4cm}p{3.2cm}p{2.0cm}p{2.0cm}p{4.8cm}}
\toprule
\textbf{Tipo} & \textbf{Ecuación Principal} & \textbf{Entrenar} & \textbf{Inferir} & \textbf{Notas} \\
\midrule
\endhead
Atención Estándar & $\text{softmax}(QK^\top/\sqrt{d_k})V$ & $\mathcal{O}(n^2d)$ & $\mathcal{O}(n)$/paso & Enrutamiento global expresivo de línea base \citep{vaswani_attention_2017}. \\
Atención Sigmoide & $\sigma(QK^\top/\sqrt{d_k}+b)V$ & $\mathcal{O}(n^2d)$ & $\mathcal{O}(n)$/paso & Compuerta elemento a elemento; a menudo necesita norma de estabilización \citep{ramapuram_theory_2024}. \\
RetNet & $(QK^\top\odot D)V$ & $\mathcal{O}(n^2d)$ o por fragmentos & $\mathcal{O}(1)$/paso & Forma dual paralela/recurrente con retención de decaimiento \citep{sun_retentive_2023}. \\
Mamba & $h_t=\bar{A}_t h_{t-1}+\bar{B}_t x_t$ & $\mathcal{O}(nd)$ & $\mathcal{O}(1)$/paso & Espacio de estados selectivo con barrido consciente de hardware \citep{gu_mamba_2023}. \\
Bloque estilo EDO & $\frac{dh}{dt}=f_\theta(h,t)$ & depende del solver & depende del solver & Interpretación de profundidad continua; integración RK \citep{zhang_continuous_2021}. \\
Memoria Titans & $M_t=(1-\alpha_t)M_{t-1}+S_t$ & aprox.\ $\mathcal{O}(nd)$ & centrado en recuperación & Actualizaciones de memoria en tiempo de prueba con dinámica impulsada por sorpresa \citep{behrouz_titans_2025}. \\
\bottomrule
\end{longtable}
\normalsize

\subsection{Resumen de Optimizadores}
\small
\begin{longtable}{p{2.2cm}p{2.8cm}p{2.0cm}p{5.4cm}p{1.8cm}}
\toprule
\textbf{Optimizador} & \textbf{Familia} & \textbf{Costo de Estado} & \textbf{Idea Clave} & \textbf{Ref} \\
\midrule
\endhead
SGD+Momentum & Primer orden clásico & Bajo & Línea base acelerada por momento con estado mínimo & \citep{polyak_1964} \\
AdamW & Adaptativo primer/segundo momento & Alto & Línea base con decaimiento de peso desacoplado & \citep{loshchilov_decoupled_2017} \\
RAdam & Adaptativo con corrección de varianza & Alto & Rectifica la varianza adaptativa temprana & \citep{liu_variance_2019} \\
Adan & Momento + reducción de varianza & Alto & Actualización adaptativa estilo Nesterov & \citep{xie_adan_2022} \\
ADOPT & Variante de Adam & Alto & Actualizaciones reordenadas con mejores garantías de convergencia & \citep{taniguchi_adopt_2024} \\
AdEMAMix & Adaptativo multi-EMA & Alto & Mezcla EMAs de horizonte corto y largo & \citep{pagliardini_ademamix_2024} \\
MARS & Precondicionado con reducción de varianza & Alto & Corrección de momento recursiva & \citep{yuan_mars_2024} \\
Cautious AdamW & Momento enmascarado & Alto & Aplica actualizaciones solo en direcciones consistentes con signo & \citep{liang_cautious_2024} \\
LAMB & Momentos adaptativos por capa & Alto & Escalado de razón de confianza para entrenamiento con lotes muy grandes & \citep{you_largebatch_2020} \\
Schedule-free AdamW & Adaptativo sin programador & Alto & Elimina la dependencia de programación LR explícita & \citep{defazio_road_2024} \\
Adafactor & Adaptativo eficiente en memoria & Medio & Segundos momentos factorizados para tensores matriciales & \citep{shazeer_adafactor_2018} \\
GaLore AdamW & Proyección de gradiente de bajo rango & Medio & Optimiza en espacio de gradiente de bajo rango proyectado & \citep{zhao_galore_2024} \\
APOLLO & Proyección adaptativa de bajo rango & Medio & Escalado estructurado a partir de momentos estilo Adam proyectados & \citep{zhu2025apollosgdlikememoryadamwlevel} \\
APOLLO-Mini & Proyección adaptativa rango-1 & Bajo & Variante APOLLO escalada por tensor para ahorro extremo de memoria & \citep{zhu2025apollosgdlikememoryadamwlevel} \\
Q-APOLLO & Proyección de bajo rango cuantizada & Bajo & Estado APOLLO cuantizado para entrenamiento de memoria ultra-baja & \citep{zhu2025apollosgdlikememoryadamwlevel} \\
Anon & Similar a Adam con adaptabilidad ajustable & Alto & Adaptabilidad $\gamma \in \mathbb{R}$ ajustable continuamente con actualización de retardo incremental (IDU) para convergencia & \citep{anon2026anon} \\
Prodigy & Adaptación sin parámetros & Medio & Calibración de paso adaptativa a distancia & \citep{mishchenko_prodigy_2023} \\
Lion & Momento de signo & Bajo & Actualización de signo de momento, estado reducido & \citep{chen_symbolic_2023} \\
Sophia & Segundo orden aproximado & Medio & Precondicionamiento diagonal de Hessiana con recorte & \citep{liu_sophia_2023} \\
Shampoo & Precondicionador matricial & Alto & Estadísticas de segundo orden estructuradas Kronecker & \citep{gupta_shampoo_2018} \\
SOAP & Shampoo + base Adam & Alto & Seguimiento similar a Adam en base propia del precondicionador & \citep{vyas_soap_2024} \\
Muon & Basado en ortogonalidad & Medio & Actualizaciones matriciales ortogonalizadas & \citep{shen_convergence_2025} \\
Turbo-Muon & Ortogonalización acelerada & Medio & Aceleración Newton-Schulz precondicionada & \citep{boissin_turbo-muon_2025} \\
\bottomrule
\end{longtable}
\normalsize

\section{Discusión}
\label{sec:discussion}

Las decisiones de diseño en Transformer Encoder Frankenstein reflejan varias tensiones de ingeniería e investigación en las herramientas modernas de aprendizaje profundo.

\subsection{Compromisos de Diseño Basado en Esquemas}

El enfoque basado en esquemas proporciona beneficios significativos de reproducibilidad al imponer contratos explícitos y fallar rápidamente ante configuraciones inválidas. Sin embargo, este enfoque también introduce rigidez: agregar nuevas arquitecturas u optimizadores requiere extensiones del esquema en lugar de argumentos de línea de comandos flexibles. El sistema de hiperparámetros prefijados permite un control detallado pero aumenta la complejidad de configuración para usuarios acostumbrados a interfaces más simples.

La decisión de imponer \texttt{additionalProperties: false} en todos los niveles del esquema elimina la absorción silenciosa de parámetros que ha afectado a sistemas de configuración anteriores, pero esta rigurosidad requiere un mantenimiento cuidadoso del esquema al extender las capacidades del sistema. Cada nuevo mecanismo de atención o variante de optimizador debe integrarse adecuadamente en el marco de validación, incluyendo definiciones de campos del esquema con tipos y restricciones apropiados, mapeo de hiperparámetros prefijados para grupos específicos de optimizadores, valores predeterminados alineados con las mejores prácticas de investigación y cadenas de documentación para la representación en la interfaz web.

\subsection{Cobertura Arquitectónica y Vacíos}

Las diecisiete arquitecturas de mezclador implementadas abarcan las principales direcciones de investigación en modelado de secuencias, pero ciertos vacíos permanecen. El sistema carece de arquitecturas híbridas recientes como Griffin \citep{soares_griffin_2024} y Jamba \citep{gu_jamba_2024}, que combinan compuertas con modelos de espacio de estados. El enrutamiento MoE (Mezcla de Expertos) está implementado para capas FFN pero no para el cómputo de atención, donde trabajos recientes han mostrado beneficios \citep{lewis_moe_2021}.

La cobertura de atención dispersa es completa, pero la implementación de métodos sin entrenamiento (FASA, SpargeAttn) genera errores en tiempo de ejecución durante el entrenamiento, reflejando restricciones arquitectónicas: estos métodos requieren checkpoints preentrenados de modelos de atención completa o procedimientos específicos de ajuste fino que actualmente no están automatizados.

La cobertura de mecanismos de compuerta es sólida en todas las categorías principales (GLA, DeltaNet, Gated DeltaNet, Gated DeltaNet-2, HGRN2, FoX, Gated Softmax).

\subsection{Fragmentación del Panorama de Optimizadores}

El soporte para veintidós familias de optimizadores en seis categorías algorítmicas demuestra exhaustividad pero también resalta el estado fragmentado de la investigación en optimización. Los usuarios enfrentan una complejidad significativa de decisión al elegir entre métodos de reducción de varianza (Adan, MARS), variantes eficientes en memoria (GaLore, Adafactor, APOLLO) y enfoques conscientes de curvatura (Sophia, Shampoo). El sistema de hiperparámetros prefijados, aunque poderoso, requiere comprender qué parámetros son relevantes para cada clase de optimizador.

La calidad de implementación varía entre optimizadores: los métodos clásicos (AdamW, SGD con momento) están altamente optimizados en PyTorch, mientras que los métodos más nuevos (Muon, Turbo-Muon, SOAP) pueden requerir implementaciones personalizadas que afectan la estabilidad numérica y las características de rendimiento.

\subsection{Consideraciones de Despliegue y Producción}

El pipeline de cuantización demuestra preocupaciones prácticas de despliegue pero realiza compromisos de ingeniería específicos. El empaquetado ternario de pesos reduce el almacenamiento a aproximadamente $1.58$ bits por parámetro, pero esta compresión agresiva puede degradar el rendimiento, especialmente para modelos más pequeños donde el error de cuantización es más significativo. La implementación actual aplica cuantización uniforme entre todos los tipos de parámetros.

Los flujos de trabajo SBERT proporcionan utilidad práctica para tareas de similitud semántica y recuperación, pero la implementación asume estrategias de pooling estándar (token CLS, pooling promedio). Avances recientes como los embeddings Matryoshka \citep{muennighoff_matryoshka_2021} y los refinamientos de aprendizaje contrastivo \citep{chen_supervised_2020} aún no están incorporados.

\subsection{Desafíos de Integración y Extensibilidad}

La estructura actual del código base, aunque funcional, presenta desafíos de mantenimiento a medida que se expanden las familias de arquitecturas y optimizadores. El patrón de despachador para la selección de mezcladores y el enrutamiento de optimizadores maneja la extensibilidad pero corre el riesgo de convertirse en un ``cajón de sastre'' de lógica condicional. Las versiones futuras se beneficiarían de arquitecturas basadas en plugins donde nuevos mezcladores y optimizadores puedan registrarse de forma declarativa en lugar de modificar la lógica central de despacho.

La interfaz de configuración web proporciona mejoras significativas de usabilidad pero introduce complejidad de despliegue: ejecutar Streamlit junto con trabajos de entrenamiento requiere recursos adicionales y consideraciones de infraestructura que pueden no ser apropiadas para todos los entornos, particularmente clústeres HPC sin acceso web.

\section{Conclusión}
\label{sec:conclusion}

Transformer Encoder Frankenstein presenta una plataforma de experimentación unificada, impulsada por configuración, que aborda desafíos críticos en la investigación moderna de aprendizaje profundo: fragmentación arquitectónica entre atención densa, modelos recurrentes, patrones dispersos y mecanismos de compuerta; complejidad del panorama de optimizadores que abarca líneas base clásicas, métodos de reducción de varianza, variantes eficientes en memoria, enfoques sin programación, algoritmos conscientes de curvatura y métodos orientados a geometría; y flujos de trabajo de despliegue de extremo a extremo que abarcan cuantización y aplicaciones de embeddings de oraciones.

Las contribuciones principales del sistema son:
\begin{enumerate}[leftmargin=*]
    \item \textbf{Diseño Basado en Esquemas}: Un contrato de configuración estricto basado en YAML con validación y enrutamiento de hiperparámetros prefijados que permite experimentos reproducibles en diecisiete arquitecturas de mezclador y veintitrés familias de optimizadores.
    \item \textbf{Soporte Arquitectónico Integral}: Implementación que abarca las principales categorías de investigación incluyendo líneas base densas (atención estándar, sigmoide), alternativas recurrentes (RetNet, Mamba, estilo EDO, Titans), atención dispersa (Sparse Transformer, Longformer, BigBird, SparseK, NSA, SpargeAttn, FASA) y mecanismos de compuerta (GLA, DeltaNet, Gated DeltaNet, HGRN2, FoX, Gated Softmax).
    \item \textbf{Marco Unificado de Optimizadores}: Grupos de hiperparámetros prefijados que permiten control detallado sobre embeddings, capas de normalización, bloques recurrentes, pesos de atención y parámetros FFN en líneas base clásicas (SGD+Momentum, AdamW), reducción de varianza (Adan, ADOPT, AdEMAMix, MARS, Cautious), eficientes en memoria (Adafactor, GaLore, Lion, APOLLO, APOLLO-Mini, Q-APOLLO), lotes grandes y simplificación de programación (LAMB, Schedule-Free AdamW, Prodigy), conscientes de curvatura (Sophia), segundo orden (Shampoo, SOAP) y orientados a geometría (Muon, Turbo-Muon).
    \item \textbf{Flujos de Trabajo de Extremo a Extremo}: Pipeline de despliegue integrado que soporta empaquetado ternario de pesos y cuantización de activaciones INT8; entrenamiento e inferencia inspirados en SBERT para similitud semántica, recuperación y tareas de clustering.
    \item \textbf{Configuración Interactiva}: Interfaz web basada en Streamlit que proporciona generación de formularios impulsada por esquemas, validación en tiempo real, documentación en línea y síntesis de comandos CLI.
\end{enumerate}

Este sistema permite iteración experimental rápida mientras mantiene reproducibilidad mediante contratos de configuración estrictos. Al consolidar diversas contribuciones de investigación en un conjunto de herramientas unificado, reduce las barreras para explorar arquitecturas y estrategias de optimización novedosas, particularmente para investigadores que pueden carecer de recursos para implementar y validar cada variante de forma independiente.

\subsection{Limitaciones y Direcciones Futuras}

Varias limitaciones y direcciones prometedoras para trabajo futuro surgen del diseño e implementación de este sistema:

\begin{enumerate}[leftmargin=*]
    \item \textbf{Integración Arquitectónica}: Arquitecturas híbridas recientes (Griffin, Jamba, Mamba-X) demuestran beneficios de combinar múltiples mecanismos en bloques unificados. Versiones futuras deberían integrar estas arquitecturas y explorar patrones de composición sistemática.
    
    \item \textbf{Cuantización Avanzada}: La implementación actual utiliza empaquetado ternario uniforme en todos los parámetros. La investigación sobre cuantización por capas, por canales y consciente de importancia sugiere que estrategias más sofisticadas podrían mejorar los compromisos calidad-eficiencia.
    
    \item \textbf{Extensibilidad Basada en Plugins}: El patrón de despacho actual se vuelve cada vez más complejo con cada nueva adición. Una arquitectura de plugins que permita el registro declarativo de nuevos mezcladores, optimizadores y métodos de normalización mejoraría el mantenimiento y reduciría el riesgo de errores en la lógica central de despacho.
    
    \item \textbf{Optimización Automatizada de Hiperparámetros}: El esquema soporta espacios extensos de hiperparámetros, pero los usuarios deben explorar estos espacios manualmente. La integración con optimización bayesiana, estrategias de bandidos multi-brazo o ajuste de hiperparámetros basado en gradientes podría automatizar el descubrimiento de configuraciones efectivas.
    
    \item \textbf{Despliegue en Producción}: La interfaz web mejora la usabilidad pero puede no ser apropiada para todos los entornos de despliegue. Modos de configuración sin interfaz gráfica, gestión de configuración basada en API o ergonomía CLI mejorada podrían servir a flujos de trabajo HPC y de producción.
    
    \item \textbf{Evaluación Comparativa}: Si bien el sistema permite entrenamiento con diversas arquitecturas, una evaluación comparativa exhaustiva que compare el rendimiento entre mezcladores y optimizadores en tareas estandarizadas proporcionaría orientación valiosa para la selección de configuraciones.
    
    \item \textbf{Garantías de Estabilidad de Entrenamiento}: La implementación actual incluye guardas contra NaN/Inf y recorte de gradiente, pero el análisis formal de condiciones de estabilidad para diferentes combinaciones mezclador-optimizador, particularmente con bloques en bucle y cuantización agresiva, sigue abierto.
    
    \item \textbf{Extensiones Multimodales y Específicas de Tarea}: El diseño actual se centra en modelado de secuencias. Extensiones para modelos de visión-lenguaje, arquitecturas multimodales y flujos de trabajo de ajuste fino específicos de tarea (ej., ajuste por instrucciones, RLHF) ampliarían la aplicabilidad.
\end{enumerate}

La trayectoria de investigación del modelado de secuencias continúa hacia enfoques híbridos que combinan fortalezas de múltiples paradigmas---compresión de la recurrencia, selectividad de la atención, compuertas para gestión de memoria y dispersión para eficiencia. Una plataforma de experimentación unificada como Transformer Encoder Frankenstein es cada vez más valiosa a medida que esta convergencia se acelera, permitiendo a los investigadores explorar sistemáticamente este creciente espacio de diseño con infraestructura reproducible y bien diseñada.

\bibliographystyle{plainnat}
\bibliography{bibliography/other,bibliography/optimizers,bibliography/attention_types}
\appendix

\section{Anexo A: Familias de Optimizadores}
\label{sec:annex-optimizers}

\subsection{Comparación de Memoria y Complejidad}

La Tabla~\ref{tab:optimizer-complexity} resume la sobrecarga de memoria (número de búferes de estado por parámetro), la complejidad computacional por paso y los hiperparámetros clave para todos los optimizadores soportados. La sobrecarga de memoria se expresa en términos del número de tensores de estado mantenidos por parámetro del modelo, donde cada tensor tiene la misma forma que el parámetro.

\begin{longtable}{llcl}
\toprule
\textbf{Optimizador} & \textbf{Búferes de Estado} & \textbf{Costo por Paso} & \textbf{Hiperparámetros Clave} \\
\midrule
\endfirsthead
\toprule
\textbf{Optimizador} & \textbf{Búferes de Estado} & \textbf{Costo por Paso} & \textbf{Hiperparámetros Clave} \\
\midrule
\endhead
\midrule
\multicolumn{4}{r}{\textit{continúa en la página siguiente}} \\
\endfoot
\bottomrule
\caption{Comparación de memoria y complejidad de todos los optimizadores soportados. $n$ = número de parámetros, $d$ = dimensionalidad del tensor, $r$ = dimensión de rango bajo, $k$ = número de iteraciones de Newton--Schulz. Búferes de estado indica el número de tensores de estado del optimizador mantenidos por grupo de parámetros. El costo por paso se enfoca en el costo adicional dominante más allá del propio cálculo del gradiente.}
\label{tab:optimizer-complexity}
\endlastfoot
SGD+Momentum & 1 ($m$) & $\mathcal{O}(n)$ & lr, momentum, wd \\
AdamW & 2 ($m$, $v$) & $\mathcal{O}(n)$ & lr, $\beta_1$, $\beta_2$, eps, wd \\
RAdam & 2 ($m$, $v$) & $\mathcal{O}(n)$ & lr, $\beta_1$, $\beta_2$, eps, wd \\
Adan & 3 ($m$, $v$, $s$) & $\mathcal{O}(n)$ & lr, $\beta_1$, $\beta_2$, $\beta_3$, eps, wd \\
ADOPT & 2 ($m$, $v$) & $\mathcal{O}(n)$ & lr, $\beta_1$, $\beta_2$, eps, wd \\
AdEMAMix & 3 ($m_1$, $m_2$, $v$) & $\mathcal{O}(n)$ & lr, $\beta_1$, $\beta_2$, $\beta_3$, eps, wd \\
MARS & 3 ($m$, $v$, $z$) & $\mathcal{O}(n)$ & lr, $\beta_1$, $\beta_2$, eps, wd, $\gamma$ \\
Cautious AdamW & 2 ($m$, $v$) & $\mathcal{O}(n)$ & lr, $\beta_1$, $\beta_2$, eps, wd \\
LAMB & 2 ($m$, $v$) & $\mathcal{O}(n)$ & lr, $\beta_1$, $\beta_2$, eps, wd \\
Schedule-Free & 3 ($z$, $x$, $n$) & $\mathcal{O}(n)$ & lr, $\beta_1$, $\beta_2$, wd \\
Adafactor & 1--2 (row/col) & $\mathcal{O}(n)$ & lr, $\beta_1$, $\beta_2$, eps, wd, factored \\
GaLore & 2 ($m$, $v$) + SVD/proj & $\mathcal{O}(nr)$ & lr, rank $r$, $\beta_1$, $\beta_2$, eps, wd \\
Prodigy & 3 ($m$, $v$, $d$) & $\mathcal{O}(n)$ & $\beta_1$, $\beta_2$, eps, wd, $d_0$ \\
Lion & 1 ($m$) & $\mathcal{O}(n)$ & lr, $\beta_1$, $\beta_2$, wd \\
Shampoo & $2d$ ($L_i$, $R_i$) & $\mathcal{O}(n^{1+2/d})$ & lr, eps, wd, matrix eps \\
SOAP & $2d$ + 2 ($m$, $v$) & $\mathcal{O}(n^{1+1/d})$ & lr, $\beta_1$, $\beta_2$, eps, wd, shampoo eps \\
Sophia & 3 ($m$, $v$, $h$) & $\mathcal{O}(n)$ & lr, $\beta_1$, $\beta_2$, eps, wd, $k$ \\
Muon & 2 ($m$, $v$) & $\mathcal{O}(n \cdot k)$ & lr, momentum, wd, NS steps $k$ \\
Turbo-Muon & 2 ($m$, $v$) & $\mathcal{O}(n \cdot k)$ & lr, momentum, wd, NS steps $k$ \\
APOLLO & 2 bajo rango ($m_R$, $v_R$) + proy & $\mathcal{O}(nr)$ & lr, rank $r$, update gap, scale, betas, eps, wd \\
APOLLO-Mini & 2 rango-1 ($m_R$, $v_R$) + proy & $\mathcal{O}(n)$ & lr, update gap, scale, betas, eps, wd \\
Q-APOLLO & 2 bajo rango cuantizados + proy & $\mathcal{O}(nr)$ & lr, rank $r$, update gap, scale, quant bits, betas, eps, wd \\
Anon & 3 ($m$, $v$, fixed\_lr) & $\mathcal{O}(n)$ & lr, $\beta_1$, $\beta_2$, eps, wd, $\gamma$ \\
\end{longtable}

\subsection{La Evolución de la Optimización en Redes Neuronales}
El estudio de optimizadores enmarca la optimización de transformadores como una respuesta a tres presiones estructurales: paisajes de pérdida no convexos, heterogeneidad severa de curvatura entre bloques de parámetros y el costo de memoria de almacenar el estado del optimizador para modelos muy grandes. El informe sostiene que el campo se ha diversificado en varias trayectorias: líneas base adaptativas de primer orden, métodos de reducción de varianza, métodos eficientes en memoria, precondicionadores estructurados de segundo orden, métodos sin programación de tasa de aprendizaje y actualizaciones orientadas a ortogonalidad.

\subsection{Línea Base Estándar y Optimizadores Adaptativos}
\paragraph{SGD con Momentum.}
La actualización clásica acumula un búfer de momentum y luego aplica una tasa de aprendizaje fija. En cada iteración $t$, el algoritmo mantiene un promedio móvil exponencial $m_t$ de los gradientes pasados:
\begin{align}
m_t &= \beta m_{t-1} + g_t \\
\theta_{t+1} &= \theta_t - \eta m_t
\end{align}
donde $g_t = \nabla f(\theta_t)$, $\beta \in [0.8, 0.99]$ controla la decadencia del momentum y $\eta$ es la tasa de aprendizaje fija. El término de momentum actúa como velocidad: acelera el movimiento en direcciones consistentes mientras amortigua oscilaciones en direcciones variables. Sus fortalezas son la baja sobrecarga de memoria (un único búfer de momentum por parámetro) y una fuerte generalización cuando se ajusta cuidadosamente. Su principal debilidad en cargas de trabajo de transformadores es la escasa robustez frente a la heterogeneidad de curvatura (diferentes espectros del Hessiano entre grupos de parámetros) y una fuerte dependencia de las programaciones de tasa de aprendizaje.

\begin{algorithm}[H]
\caption{SGD con Momentum}
\label{alg:sgd-momentum}
\begin{algorithmic}[1]
\Require Parámetros iniciales $\theta_0$, tasa de aprendizaje $\eta$, coeficiente de momentum $\beta$, decadencia de peso $\lambda$
\Ensure Parámetros actualizados $\theta_T$
\State Inicializar: búfer de momentum $m \gets 0$
\For{$t = 0, 1, 2, \ldots, T-1$}
    \State Calcular gradiente: $g_t \gets \nabla f(\theta_t)$
    \If{decadencia\_de\_peso $> 0$}
        \State $g_t \gets g_t + \lambda \theta_t$ \Comment{Regularización L2}
    \EndIf
    \State Actualizar momentum: $m \gets \beta \cdot m + g_t$
    \State Actualizar parámetros: $\theta_{t+1} \gets \theta_t - \eta \cdot m$
\EndFor
\Return $\theta_T$
\end{algorithmic}
\end{algorithm}

\paragraph{Adam y AdamW.}
Adam rastrea promedios móviles exponenciales tanto del primer momento (media) como del segundo momento (varianza no centrada) para lograr tasas de aprendizaje adaptativas elemento a elemento. La regla de actualización es:
\begin{align}
m_t &= \beta_1 m_{t-1} + (1-\beta_1) g_t \\
v_t &= \beta_2 v_{t-1} + (1-\beta_2) g_t^2 \\
\hat{m}_t &= \frac{m_t}{1-\beta_1^t}, \quad \hat{v}_t = \frac{v_t}{1-\beta_2^t} \quad \text{(corrección de sesgo)} \\
\theta_{t+1} &= \theta_t - \eta \left(\frac{\hat{m}_t}{\sqrt{\hat{v}_t}+\epsilon}\right)
\end{align}
AdamW introduce una modificación crucial: la decadencia de peso se aplica directamente a los parámetros (desacoplada de los gradientes): $\theta_t \gets \theta_t(1-\eta\lambda)$ en lugar de añadir $\lambda\theta_t$ al gradiente. Este desacoplamiento evita que el escalado adaptativo interfiera con la fuerza de la regularización. El informe trata a AdamW como la línea base práctica para el ajuste fino de transformadores porque converge rápidamente y es relativamente tolerante a la variación de hiperparámetros. La desventaja es el costo de memoria: ambos tensores de momento deben almacenarse para cada parámetro, duplicando la memoria del estado del optimizador en relación con SGD.

\begin{algorithm}[H]
\caption{AdamW (Adam con Decadencia de Peso Desacoplada)}
\label{alg:adamw}
\begin{algorithmic}[1]
\Require Parámetros iniciales $\theta_0$, tasa de aprendizaje $\eta$, tasas de decadencia exponencial $\beta_1, \beta_2 \in [0,1)$
\Require Constante de momentum $\epsilon$, decadencia de peso $\lambda$
\Ensure Parámetros actualizados $\theta_T$
\State Inicializar: primer momento $m \gets 0$, segundo momento $v \gets 0$, contador de pasos $t \gets 0$
\For{$t = 1, 2, \ldots, T$}
    \State Calcular gradiente: $g_t \gets \nabla f(\theta_{t-1})$
    \State Decadencia de peso (desacoplada): $\theta_{t-1} \gets \theta_{t-1}(1 - \eta\lambda)$
    \State Actualizar momentos: $m \gets \beta_1 m + (1-\beta_1) g_t$
    \State $v \gets \beta_2 v + (1-\beta_2) g_t^2$
    \State Corrección de sesgo: $\hat{m} \gets m / (1 - \beta_1^t)$, $\hat{v} \gets v / (1 - \beta_2^t)$
    \State Actualizar parámetros: $\theta_t \gets \theta_{t-1} - \eta (\hat{m} / (\sqrt{\hat{v}} + \epsilon))$
\EndFor
\Return $\theta_T$
\end{algorithmic}
\end{algorithm}

\paragraph{RAdam (Adam Rectificado).}
RAdam aborda la inestabilidad de Adam en el entrenamiento temprano rectificando dinámicamente la tasa de aprendizaje adaptativa. La observación clave es que el segundo momento $v_t$ tiene una varianza muy alta en los primeros pasos, causando un escalado adaptativo poco fiable. RAdam calcula la longitud efectiva de la ventana de promedio móvil simple (SMA):
\[\rho_t = \rho_{\infty} - \frac{2t\beta_2^t}{1-\beta_2^t}, \quad \text{donde} \quad \rho_{\infty} = \frac{2}{1-\beta_2}-1\]
Cuando $\rho_t > 4$ (suficientes muestras para la estimación de varianza), RAdam aplica el escalado adaptativo con un término de rectificación:
\[r_t = \sqrt{\frac{(\rho_t-4)(\rho_t-2)\rho_{\infty}}{(\rho_{\infty}-4)(\rho_{\infty}-2)\rho_t}}, \quad \theta_{t+1} = \theta_t - \eta r_t \frac{\hat{m}_t}{\sqrt{\hat{v}_t}+\epsilon}\]
Cuando $\rho_t \leq 4$, RAdam retrocede a SGD con momentum: $\theta_{t+1} = \theta_t - \eta\hat{m}_t$. Esta transición gradual elimina la necesidad de programaciones manuales de calentamiento de la tasa de aprendizaje y mejora la robustez.

\begin{algorithm}[H]
\caption{RAdam (Adam Rectificado)}
\label{alg:radam}
\begin{algorithmic}[1]
\Require Parámetros iniciales $\theta_0$, tasa de aprendizaje $\eta$, $\beta_1, \beta_2$, decadencia de peso $\lambda$
\Ensure Parámetros actualizados $\theta_T$
\State Inicializar: $m \gets 0$, $v \gets 0$, $t \gets 0$
\State Calcular: $\rho_\infty \gets 2/(1-\beta_2) - 1$
\For{$t = 1, 2, \ldots, T$}
    \State $g_t \gets \nabla f(\theta_{t-1})$
    \If{decadencia\_de\_peso $> 0$}
        \State $\theta_{t-1} \gets \theta_{t-1}(1 - \eta\lambda)$
    \EndIf
    \State $m \gets \beta_1 m + (1-\beta_1) g_t$
    \State $v \gets \beta_2 v + (1-\beta_2) g_t^2$
    \State $\hat{m} \gets m/(1-\beta_1^t)$, $\hat{v} \gets v/(1-\beta_2^t)$
    \State $\rho_t \gets \rho_\infty - 2t \beta_2^t / (1 - \beta_2^t)$
    \If{$\rho_t > 4$}
        \State $r_t \gets \sqrt{((\rho_t-4)(\rho_t-2)\rho_\infty) / ((\rho_\infty-4)(\rho_\infty-2)\rho_t)}$
        \State $\theta_t \gets \theta_{t-1} - \eta r_t \hat{m} / (\sqrt{\hat{v}} + \epsilon)$
    \Else
        \State $\theta_t \gets \theta_{t-1} - \eta \hat{m}$ \Comment{SGD con momentum}
    \EndIf
\EndFor
\Return $\theta_T$
\end{algorithmic}
\end{algorithm}

\subsection{Momentum Avanzado y Reducción de Varianza (2024--2025)}
\paragraph{Adan (Momentum Nesterov Adaptativo).}
Adan reformula la aceleración de Nesterov sin el cálculo de gradiente adicional requerido por el SGD Nesterov clásico. El algoritmo mantiene tres búferes de momentum:
\begin{align}
m_t &= (1-\beta_1)m_{t-1} + \beta_1 g_t \quad\text{(primer momento)} \\
v_t &= (1-\beta_2)v_{t-1} + \beta_2(g_t - g_{t-1}) \quad\text{(velocidad/diferencia de gradientes)} \\
n_t &= (1-\beta_3)n_{t-1} + \beta_3[g_t + (1-\beta_1)(g_t-g_{t-1})]^2 \quad\text{(segundo momento Nesterov)}
\end{align}
El término de \textbf{Estimación de Momentum Nesterov} (NME) $\bar{g}_t = g_t + (1-\beta_1)(g_t-g_{t-1})$ estima el gradiente en una posición futura sin evaluarlo. La actualización combina el momentum con la velocidad:
\[\bar{m}_t = m_t + (1-\beta_1)v_t, \quad \theta_{t+1} = \theta_t - \eta\frac{\bar{m}_t}{\sqrt{n_t}+\epsilon}\]
Adan logra una convergencia rápida en diversas arquitecturas (CNN, GAN, Transformadores) mediante esta aceleración. El costo es mantener tres búferes similares al momentum, aumentando la sobrecarga de memoria en relación con Adam.

\begin{algorithm}[H]
\caption{Adan (Momentum Nesterov Adaptativo)}
\label{alg:adan}
\begin{algorithmic}[1]
\Require Parámetros iniciales $\theta_0$, tasa de aprendizaje $\eta$, $\beta_1, \beta_2, \beta_3 \in (0,1)$
\Ensure Parámetros actualizados $\theta_T$
\State Inicializar: $m \gets 0$, $v \gets 0$, $n \gets 0$, $g_{-1} \gets 0$ (gradiente anterior)
\For{$t = 1, 2, \ldots, T$}
    \State $g_t \gets \nabla f(\theta_{t-1})$
    \State $m \gets (1-\beta_1) m + \beta_1 g_t$
    \State $\Delta g_t \gets g_t - g_{t-1}$
    \State $v \gets (1-\beta_2) v + \beta_2 \Delta g_t$
    \State Estimación Nesterov: $\bar{g}_t \gets g_t + (1-\beta_1) \Delta g_t$
    \State $n \gets (1-\beta_3) n + \beta_3 \bar{g}_t^2$
    \State Momentum combinado: $\bar{m} \gets m + (1-\beta_1) v$
    \State $\theta_t \gets \theta_{t-1} - \eta (\bar{m} / (\sqrt{n} + \epsilon))$
    \State $g_{t-1} \gets g_t$
\EndFor
\Return $\theta_T$
\end{algorithmic}
\end{algorithm}

\paragraph{ADOPT (Adam con Ajuste Óptimo Podado).}
ADOPT corrige un problema teórico fundamental en Adam: el gradiente aparece tanto en la estimación del primer momento como en la del segundo momento, creando circularidad. ADOPT \textbf{desacopla} utilizando el segundo momento del paso anterior para el denominador:
\begin{align}
v_t &= \beta_2 v_{t-1} + (1-\beta_2)g_t^2 \\
m_t &= \beta_1 m_{t-1} + (1-\beta_1)\frac{g_t}{\sqrt{v_{t-1}}+\epsilon} \quad\text{(usa $v_{t-1}$, no $v_t$)} \\
\theta_{t+1} &= \theta_t - \eta m_t
\end{align}
Este simple reordenamiento logra la tasa de convergencia óptima $O(1/\sqrt{T})$ con \textbf{cualquier} elección de $\beta_2 \in (0,1)$, sin suposiciones de ruido acotado. La consecuencia práctica es que ADOPT es un reemplazo directo de Adam con garantías teóricas más sólidas y un rendimiento empírico comparable o superior en los dominios de visión, PLN, RL y modelado generativo.

\begin{algorithm}[H]
\caption{ADOPT (Adam con Ajuste Óptimo Podado)}
\label{alg:adopt}
\begin{algorithmic}[1]
\Require Parámetros iniciales $\theta_0$, tasa de aprendizaje $\eta$, $\beta_1, \beta_2$, tolerancia $\epsilon$
\Ensure Parámetros actualizados $\theta_T$
\State Inicializar: $m \gets 0$, $v \gets 0$
\For{$t = 1, 2, \ldots, T$}
    \State $g_t \gets \nabla f(\theta_{t-1})$
    \If{decadencia\_de\_peso $> 0$}
        \State $\theta_{t-1} \gets \theta_{t-1}(1 - \eta\lambda)$
    \EndIf
    \State Usar segundo momento \textbf{anterior}: $\text{denom} \gets \sqrt{\max(v, \epsilon)} + \epsilon$
    \State Actualizar primer momento usando varianza anterior: $m \gets \beta_1 m + (1-\beta_1)(g_t / \text{denom})$
    \State Actualizar segundo momento con gradiente \textbf{actual}: $v \gets \beta_2 v + (1-\beta_2) g_t^2$
    \State $\theta_t \gets \theta_{t-1} - \eta m$
\EndFor
\Return $\theta_T$
\end{algorithmic}
\end{algorithm}

\paragraph{AdEMAMix (Mezcla de Promedios Móviles Exponenciales).}
AdEMAMix reemplaza el EMA único de gradientes de Adam por una \textbf{mezcla de dos EMA}: uno de decadencia rápida y otro de decadencia lenta. Esto aborda la observación de que los gradientes siguen siendo informativos durante decenas de miles de pasos, no solo cientos:
\begin{align}
m_{1,t} &= \beta_1 m_{1,t-1} + (1-\beta_1)g_t \quad\text{(EMA rápido, $\beta_1 \approx 0.9$)} \\
m_{2,t} &= \beta_3 m_{2,t-1} + (1-\beta_3)g_t \quad\text{(EMA lento, $\beta_3 \approx 0.9999$)} \\
v_t &= \beta_2 v_{t-1} + (1-\beta_2)g_t^2 \\
m_t &= m_{1,t} + \alpha_t m_{2,t} \quad\text{(mezcla con peso programado $\alpha_t$)} \\
\theta_{t+1} &= \theta_t - \eta\frac{m_t}{\sqrt{v_t}+\epsilon}
\end{align}
El peso de la mezcla $\alpha_t$ típicamente aumenta durante el entrenamiento, permitiendo que el EMA rápido proporcione adaptación inmediata mientras que el EMA lento acumula correlaciones de gradiente a largo plazo. Empíricamente, un modelo de 1.3B parámetros en 101B tokens alcanza una pérdida similar a AdamW en 197B tokens (ganancia del 95\% en eficiencia de datos), lo que sugiere que el EMA lento reduce significativamente el olvido del modelo.

\begin{algorithm}[H]
\caption{AdEMAMix (Mezcla de Promedios Móviles Exponenciales)}
\label{alg:ademamix}
\begin{algorithmic}[1]
\Require Parámetros iniciales $\theta_0$, tasa de aprendizaje $\eta$, $\beta_1$ (rápido), $\beta_3$ (lento), $\beta_2$ (segundo momento)
\Require Peso de mezcla $\alpha_t$ (típicamente creciente), tolerancia $\epsilon$
\Ensure Parámetros actualizados $\theta_T$
\State Inicializar: $m_1 \gets 0$ (EMA rápido), $m_2 \gets 0$ (EMA lento), $v \gets 0$
\For{$t = 1, 2, \ldots, T$}
    \State $g_t \gets \nabla f(\theta_{t-1})$
    \If{decadencia\_de\_peso $> 0$}
        \State $\theta_{t-1} \gets \theta_{t-1}(1 - \eta\lambda)$
    \EndIf
    \State $m_1 \gets \beta_1 m_1 + (1-\beta_1) g_t$ \Comment{EMA rápido}
    \State $m_2 \gets \beta_3 m_2 + (1-\beta_3) g_t$ \Comment{EMA lento}
    \State $v \gets \beta_2 v + (1-\beta_2) g_t^2$
    \State $m_t \gets m_1 + \alpha_t m_2$ \Comment{Mezcla}
    \State $\theta_t \gets \theta_{t-1} - \eta (m_t / (\sqrt{v} + \epsilon))$
\EndFor
\Return $\theta_T$
\end{algorithmic}
\end{algorithm}

\paragraph{MARS (Haciendo Brillar las Tasas de Aprendizaje Adaptativas).}
MARS combina métodos de gradiente precondicionado (ej., AdamW) con \textbf{reducción de varianza} mediante momentum recursivo estocástico escalado. La innovación central es una estimación del gradiente con varianza reducida:
\begin{align}
c_t &= g_t + \gamma(c_{t-1} - g_{t-1}) \quad\text{(estimación recursiva tipo SVRG)} \\
m_t &= \beta_1 m_{t-1} + (1-\beta_1)c_t \\
v_t &= \beta_2 v_{t-1} + (1-\beta_2)c_t^2 \\
\theta_{t+1} &= \theta_t - \eta\frac{m_t}{\sqrt{v_t}+\epsilon}
\end{align}
donde $\gamma \in [0.01, 0.1]$ controla cuánta información de gradiente histórico se retiene. El gradiente con varianza reducida $c_t$ actúa como un filtro de ruido implícito, amplificando las direcciones de señal consistentes y amortiguando el ruido contradictorio. MARS-AdamW supera consistentemente a AdamW por márgenes significativos en el preentrenamiento de GPT-2, lo que sugiere que la reducción de varianza mejora sustancialmente la convergencia en el entrenamiento estocástico por mini-lotes.

\begin{algorithm}[H]
\caption{MARS (Haciendo Brillar las Tasas de Aprendizaje Adaptativas)}
\label{alg:mars}
\begin{algorithmic}[1]
\Require Parámetros iniciales $\theta_0$, tasa de aprendizaje $\eta$, $\beta_1, \beta_2$, decadencia de peso $\lambda$
\Require Coeficiente de reducción de varianza $\gamma \in [0.01, 0.1]$
\Ensure Parámetros actualizados $\theta_T$
\State Inicializar: $m \gets 0$, $v \gets 0$, $c \gets 0$ (gradiente con varianza reducida), $g_{-1} \gets 0$
\For{$t = 1, 2, \ldots, T$}
    \State $g_t \gets \nabla f(\theta_{t-1})$
    \State Reducción de varianza: $c \gets g_t + \gamma(c - g_{t-1})$
    \State $m \gets \beta_1 m + (1-\beta_1) c$
    \State $v \gets \beta_2 v + (1-\beta_2) c^2$
    \If{decadencia\_de\_peso $> 0$}
        \State $\theta_{t-1} \gets \theta_{t-1}(1 - \eta\lambda)$
    \EndIf
    \State $\theta_t \gets \theta_{t-1} - \eta (m / (\sqrt{v} + \epsilon))$
    \State $g_{t-1} \gets g_t$
\EndFor
\Return $\theta_T$
\end{algorithmic}
\end{algorithm}

\paragraph{Optimizadores Cautelosos (Cautious AdamW, Cautious Lion).}
El marco Cauteloso aplica una \textbf{modificación de una línea} a cualquier optimizador basado en momentum: enmascarar la actualización de modo que solo se apliquen las dimensiones donde el momentum y la dirección del gradiente coinciden:
\begin{align}
\text{mask}_i &= \begin{cases} 1 & \text{si } m_t[i] \cdot g_t[i] > 0 \quad\text{(acuerdo)} \\ 0 & \text{en caso contrario} \end{cases} \\
u_t &= \left(\frac{m_t}{\sqrt{v_t}+\epsilon}\right) \odot \text{mask} \quad\text{(enmascaramiento elemento a elemento)} \\
\theta_{t+1} &= \theta_t - \eta u_t
\end{align}
La intuición: el momentum $m_t$ estima la dirección del gradiente a partir del historial; el gradiente actual $g_t$ es la señal instantánea. Cuando ambos coinciden, el optimizador está seguro y debe actualizar agresivamente. Cuando discrepan, el optimizador está en conflicto---la tendencia histórica apunta hacia una región que pudo haber sido buena antes, pero la evidencia actual la contradice. Al enmascarar las dimensiones en conflicto, Cauteloso se vuelve más conservador y evita pasos corruptos. Empíricamente, este simple enmascaramiento logra hasta \textbf{1.47$\times$ de aceleración} en el preentrenamiento de Llama y MAE mientras preserva las garantías de convergencia.

\begin{algorithm}[H]
\caption{Cautious AdamW (Enmascaramiento de Actualización por Consenso)}
\label{alg:cautious}
\begin{algorithmic}[1]
\Require Parámetros iniciales $\theta_0$, tasa de aprendizaje $\eta$, $\beta_1, \beta_2$, decadencia de peso $\lambda$
\Ensure Parámetros actualizados $\theta_T$
\State Inicializar: $m \gets 0$, $v \gets 0$
\For{$t = 1, 2, \ldots, T$}
    \State $g_t \gets \nabla f(\theta_{t-1})$
    \State $m \gets \beta_1 m + (1-\beta_1) g_t$
    \State $v \gets \beta_2 v + (1-\beta_2) g_t^2$
    \State Máscara de consenso: $\text{mask}_i \gets 1$ si $m[i] \cdot g_t[i] > 0$, si no $0$ \Comment{Acuerdo}
    \State Actualización base Adam: $u \gets m / (\sqrt{v} + \epsilon)$
    \State Aplicar máscara: $u_{\text{masked}} \gets u \odot \text{mask}$ \Comment{Elemento a elemento}
    \If{decadencia\_de\_peso $> 0$}
        \State $\theta_{t-1} \gets \theta_{t-1}(1 - \eta\lambda)$
    \EndIf
    \State $\theta_t \gets \theta_{t-1} - \eta u_{\text{masked}}$
\EndFor
\Return $\theta_T$
\end{algorithmic}
\end{algorithm}

\subsection{Optimizadores de Lote Grande, Eficientes en Memoria y Sin Tasa de Aprendizaje}
\paragraph{LAMB (Layer-wise Adaptive Moments optimizer for Batch training).}
LAMB extiende Adam con \textbf{escalado de tasa adaptativo por capas} (inspirado en LARS), permitiendo entrenamiento estable con tamaños de lote extremos (e.g., 64K en BERT):
\begin{align}
m_t^L &= \beta_1 m_{t-1}^L + (1-\beta_1)g_t^L \quad\text{(momento de primer orden por capas)} \\
v_t^L &= \beta_2 v_{t-1}^L + (1-\beta_2)(g_t^L)^2 \\
u_{\text{adam}}^L &= \frac{m_t^L}{\sqrt{v_t^L}+\epsilon} + \lambda\theta_{t-1}^L \quad\text{(paso adaptativo base con decaimiento)} \\
\phi^L &= \frac{\|\theta_{t-1}^L\|_2}{\|u_{\text{adam}}^L\|_2} \quad\text{(razón de confianza: normalización por capas)} \\
\theta_t^L &= \theta_{t-1}^L - \eta \cdot \phi^L \cdot u_{\text{adam}}^L
\end{align}
La \textbf{razón de confianza} $\phi^L = \|\theta^L\|_2 / \|u_{\text{adam}}^L\|_2$ normaliza el paso de actualización efectivo en relación con la magnitud de los pesos. En lotes muy grandes, el ruido estocástico del gradiente puede superar a la señal, desestabilizando las tasas de aprendizaje por capas. Al escalar las actualizaciones proporcionalmente a las normas de los pesos, LAMB asegura cambios relativos en lugar de absolutos, evitando que actualizaciones pequeñas y seguras en vectores grandes dominen. LAMB preserva los beneficios de generalización de lotes pequeños mientras permite entrenamiento eficiente con lotes grandes.

\begin{algorithm}[H]
\caption{LAMB (Layer-wise Adaptive Moments optimizer for Batch training)}
\label{alg:lamb}
\begin{algorithmic}[1]
\Require Parámetros iniciales $\theta_0$, tasa de aprendizaje $\eta$, $\beta_1, \beta_2$, decaimiento de pesos $\lambda$
\Ensure Parámetros actualizados $\theta_T$
\State Inicializar: Para cada capa $L$: $m^L \gets 0$, $v^L \gets 0$
\For{$t = 1, 2, \ldots, T$}
    \For{cada capa $L$}
        \State $g_t^L \gets \nabla f(\theta_t^L)$
        \State $m^L \gets \beta_1 m^L + (1-\beta_1) g_t^L$
        \State $v^L \gets \beta_2 v^L + (1-\beta_2) (g_t^L)^2$
        \State Adam base: $u_{\text{adam}}^L \gets m^L / (\sqrt{v^L} + \epsilon)$
        \If{weight\_decay $> 0$}
            \State $u_{\text{adam}}^L \gets u_{\text{adam}}^L + \lambda \theta_{t-1}^L$
        \EndIf
        \State Razón de confianza: $\phi^L \gets \|\theta_{t-1}^L\|_2 / \|u_{\text{adam}}^L\|_2$ si ambos son no nulos, sino $1$
        \State $\theta_t^L \gets \theta_{t-1}^L - \eta \phi^L u_{\text{adam}}^L$
    \EndFor
\EndFor
\Return $\theta_T$
\end{algorithmic}
\end{algorithm}

\paragraph{Schedule-Free AdamW.}
Los métodos Schedule-Free eliminan el diseño explícito de programación de la receta de optimización. El algoritmo mantiene dos flujos de parámetros: $z_t$ (exploración) y $x_t$ (promedio suavizado):
\begin{align}
v_t &= \beta_2 v_{t-1} + (1-\beta_2)g_t^2 \quad\text{(varianza)} \\
z_t &= z_{t-1}(1-\eta\lambda) - \eta\frac{g_t}{\sqrt{v_t}+\epsilon} \quad\text{(paso adaptativo con decaimiento)} \\
c_t &= \frac{1}{t+1} \quad\text{(coeficiente de promediado, decrece con el tiempo)} \\
x_t &= (1-c_t)x_{t-1} + c_t z_t \quad\text{(promediado de iteraciones)} \\
y_t &= (1-\beta)z_t + \beta x_t \quad\text{(interpolación para el punto de evaluación)}
\end{align}
El coeficiente de promediado $c_t = 1/(t+1)$ implementa una programación de tasa de aprendizaje implícita sin especificar explícitamente el número total de pasos $T$. Este marco unifica la programación y el promediado de iteraciones: el algoritmo explora mediante $z_t$ mientras acumula dirección estable mediante $x_t$. El punto de evaluación $y_t$ (usado para el cómputo del gradiente) interpola entre exploración y estabilidad. Schedule-Free logra convergencia de última generación en optimización convexa, aprendizaje profundo a gran escala y aprendizaje por refuerzo, eliminando un hiperparámetro importante.

\begin{algorithm}[H]
\caption{Schedule-Free AdamW}
\label{alg:schedulefree}
\begin{algorithmic}[1]
\Require Parámetros iniciales $\theta_0$, tasa de aprendizaje $\eta$ (fija), $\beta_1, \beta_2$, decaimiento de pesos $\lambda$
\Ensure Parámetros actualizados $\theta_T$ o promediado $x_T$
\State Inicializar: $z \gets \theta_0$ (exploración), $x \gets \theta_0$ (promedio), $v \gets 0$, $t \gets 0$
\For{$t = 1, 2, \ldots, T$}
    \State $g_t \gets \nabla f(y_{t-1})$ (gradiente en el punto de interpolación)
    \State $v \gets \beta_2 v + (1-\beta_2) g_t^2$ \Comment{Varianza}
    \State Decaimiento de pesos sobre $z$: $z \gets z(1 - \eta\lambda)$
    \State $z \gets z - \eta g_t / (\sqrt{v} + \epsilon)$ \Comment{Paso adaptativo sobre el gradiente crudo}
    \State $c_t \gets 1 / (t+1)$ \Comment{Coeficiente de promediado}
    \State $x \gets (1-c_t) x + c_t z$ \Comment{Promediado de iteraciones}
    \State $y_t \gets (1-\beta_1) z + \beta_1 x$ \Comment{Interpolación para la siguiente evaluación}
\EndFor
\Return $x_T$ o $y_T$
\end{algorithmic}
\end{algorithm}

\paragraph{Adafactor.}
Adafactor reduce la memoria del optimizador \textbf{factorizando} estadísticas de segundo momento para parámetros con forma de matriz, almacenando solo acumuladores de filas y columnas en lugar de estados de varianza densos. Para una matriz de gradiente $G_t \in \mathbb{R}^{m \times n}$:
\begin{align}
R_t &= \beta_2 R_{t-1} + (1-\beta_2)(G_t^2) \mathbf{1}_n^T \quad\text{(varianza por filas)} \\
C_t &= \beta_2 C_{t-1} + (1-\beta_2) \mathbf{1}_m^\top (G_t^2) \quad\text{(varianza por columnas)} \\
\hat{V}_t &= \frac{R_t C_t}{\mathbf{1}_n^\top R_t} \quad\text{(varianza reconstruida mediante producto exterior)} \\
U_t &= \frac{G_t}{\sqrt{\hat{V}_t}+\epsilon} \quad\text{(paso adaptativo normalizado)} \\
\hat{U}_t &= \frac{U_t}{\max(1, \text{RMS}(U_t))} \quad\text{(recorte de estabilidad)}
\end{align}
En lugar de almacenar $m \times n$ valores de segundo momento, Adafactor almacena solo $m + n$ acumuladores, reduciendo la memoria de $O(n_{\text{params}})$ a $O(\sqrt{n_{\text{params}}})$. Esto es más atractivo cuando la VRAM está dominada por el estado del optimizador en lugar de las activaciones. La contrapartida es una expresividad de optimización reducida y posible inestabilidad en algunas tareas.

\begin{algorithm}[H]
\caption{Adafactor (Estadísticas de Segundo Momento Factorizadas)}
\label{alg:adafactor}
\begin{algorithmic}[1]
\Require Matriz de gradiente $G_t \in \mathbb{R}^{m \times n}$, tasa de aprendizaje $\eta$, $\beta_2 \approx 0.999$
\Ensure Parámetros actualizados
\State Inicializar: Acumuladores de fila $R \gets \epsilon \mathbf{1}^T_m$, Columna $C \gets \epsilon \mathbf{1}_n$
\For{$t = 1, 2, \ldots, T$}
    \State $G_t \gets \nabla f(\theta_t)$ (gradiente matricial)
    \State $R \gets \beta_2 R + (1-\beta_2) (G_t^2) \mathbf{1}_n^T$ \Comment{Productos internos por filas}
    \State $C \gets \beta_2 C + (1-\beta_2) \mathbf{1}_m^\top (G_t^2)$ \Comment{Productos internos por columnas}
    \State Varianza reconstruida: $\hat{V} \gets (R \cdot C) / (\mathbf{1}_n^\top R)$ \Comment{Mediante producto exterior}
    \State Paso adaptativo normalizado: $U_t \gets G_t / (\sqrt{\hat{V}} + \epsilon)$
    \State Recorte RMS por fila: $\hat{U}_t \gets U_t / \max(1, \text{RMS}(U_t))$
    \State $\theta_{t+1} \gets \theta_t - \eta \hat{U}_t$
\EndFor
\end{algorithmic}
\end{algorithm}

\paragraph{GaLore (Gradient Low-Rank Projection).}
GaLore proyecta gradientes 2D en un subespacio de bajo rango antes de la optimización, reduciendo la memoria del estado del optimizador mientras mantiene el escalado adaptativo del paso. Para una matriz de parámetros 2D $W \in \mathbb{R}^{m \times n}$:
\begin{align}
U, S, V &= \text{SVD}(G_t) \quad\text{(calcular descomposición singular)} \\
P &\in \mathbb{R}^{n \times r} \quad\text{o} \quad P^\top \in \mathbb{R}^{r \times m} \quad\text{(seleccionar los $r$ vectores singulares superiores)} \\
G_{\text{low}} &= P^\top G_t \quad\text{(proyectar al espacio de bajo rango)} \\
\Delta_{\text{low}} &= \text{Adam}(G_{\text{low}}) \quad\text{(optimizar en el espacio comprimido)} \\
\Delta &= P \Delta_{\text{low}} \quad\text{(reconstruir en el espacio original)}
\end{align}
Al proyectar en un subespacio de rango $r$ (típicamente $r \ll \min(m,n)$), el estado del optimizador se reduce de $O(mn)$ a $O(r(m+n))$ para parámetros 2D. Este enfoque complementario de ahorro de memoria es especialmente relevante cuando el modelo es demasiado grande para el estado completo del optimizador. GaLore funciona sinérgicamente con otras técnicas y ha mostrado resultados empíricos sólidos en modelos de miles de millones de parámetros.

\begin{algorithm}[H]
\caption{GaLore (Gradient Low-Rank Projection)}
\label{alg:galore}
\begin{algorithmic}[1]
\Require Matriz de parámetros 2D $W \in \mathbb{R}^{m \times n}$, tasa de aprendizaje $\eta$, rango $r \ll \min(m,n)$
\Ensure Parámetros actualizados
\State Inicializar: Estado de Adam en el espacio de bajo rango (si es variante de rango 2)
\For{$t = 1, 2, \ldots, T$}
    \State $G_t \gets \nabla f(W_t)$
    \If{$t \mod K = 1$}  \Comment{SVD periódico}
        \State $U_t, S_t, V_t^\top \gets \text{SVD}(G_t)$ (completa o delgada)
        \State Seleccionar proyección: $P \gets U_t[:, :r]$ o $P \gets V_t[:, :r]$
    \EndIf
    \State Proyectar gradiente: $G_{\text{low}} \gets P^\top G_t$ (o $G_t P^\top$ para proyección derecha)
    \State Ejecutar Adam en bajo rango: $\Delta_{\text{low}} \gets \text{Adam}(G_{\text{low}})$
    \State Reconstruir en espacio original: $\Delta \gets P \Delta_{\text{low}}$ (o $\Delta_{\text{low}} P^\top$)
    \State $W_t \gets W_t - \eta \Delta$
\EndFor
\end{algorithmic}
\end{algorithm}

\paragraph{APOLLO, APOLLO-Mini, y Q-APOLLO.}
La familia APOLLO \citep{zhu2025apollosgdlikememoryadamwlevel} parte de la observación de que el denominador elemento a elemento de AdamW puede \textbf{transformarse en una actualización estructurada de la tasa de aprendizaje}. En lugar de almacenar momentos densos para cada entrada de parámetro, APOLLO proyecta un gradiente matricial $G_t \in \mathbb{R}^{m \times n}$ en un subespacio aleatorio compacto y rastrea momentos estilo Adam allí:
\begin{align}
R_t &= P_t G_t \quad\text{o}\quad R_t = G_t P_t^\top \\
M_t^R &= \beta_1 M_{t-1}^R + (1-\beta_1)R_t \\
V_t^R &= \beta_2 V_{t-1}^R + (1-\beta_2)R_t^2 \\
\widetilde{R}_t &= \frac{M_t^R}{\sqrt{V_t^R}+\epsilon}
\end{align}
El estado proyectado \emph{no} se expande de vuelta a una actualización densa de bajo rango como en los métodos basados en SVD. En su lugar, APOLLO estima un tensor de escalado estructurado $S_t$ en el espacio original. En la variante estándar de APOLLO, ese escalado es \textbf{por canales}: cada fila o canal recibe su propia razón de norma. En APOLLO-Mini, el escalado se reduce a un único escalar \textbf{por tensor}, correspondiente al caso extremo de rango 1 descrito en el artículo. La actualización de parámetros resultante es, por tanto, similar a Adam en adaptación pero mucho más cercana a SGD en costo de estado:
\[
W_t \gets (1-\eta\lambda)W_{t-1} - \eta\,\alpha\,(G_t \odot S_t)
\]
donde $\alpha$ es el factor de escala adicional usado para estabilizar variantes altamente comprimidas.

La contribución práctica del artículo es doble. Primero, APOLLO reemplaza el costoso SVD repetido con \textbf{proyección aleatoria gaussiana} actualizada periódicamente, de modo que la carga computacional es multiplicación de matrices ordinaria en lugar de descomposición espectral. Segundo, el optimizador es inusualmente tolerante a la compresión extrema: incluso \textbf{APOLLO-Mini}, que mantiene solo estado auxiliar de rango 1, sigue siendo competitivo o mejor que AdamW en los experimentos de preentrenamiento reportados, mientras se acerca al costo de memoria de SGD.

\begin{algorithm}[H]
\caption{APOLLO / APOLLO-Mini Escalado de Gradiente Estructurado}
\label{alg:apollo-family}
\begin{algorithmic}[1]
\Require Matriz de pesos $W \in \mathbb{R}^{m \times n}$ con $m \le n$, tasa de aprendizaje $\eta$, factor de escala $\alpha$, tasas de decaimiento $(\beta_1, \beta_2)$, decaimiento de pesos $\lambda$, rango $r$, intervalo de actualización de proyección $T$
\Ensure Parámetros actualizados $W_T$
\State Inicializar momentos proyectados: $M^R \gets 0$, $V^R \gets 0$, paso $t \gets 0$
\Repeat
    \State Calcular gradiente: $G_t \gets \nabla_W \phi(W_t)$
    \If{$t \bmod T = 0$}
        \State Muestrear proyector gaussiano $P_t \sim \mathcal{N}(0, 1/r)$ con una semilla nueva
    \EndIf
    \State Proyectar gradiente: $R_t \gets P_t G_t$ \Comment{o $G_t P_t^\top$ según la disposición}
    \State Actualizar momentos proyectados AdamW: $M_t^R, V_t^R \gets \text{AdamWState}(R_t; \beta_1, \beta_2)$
    \State Normalizar estado proyectado: $\widetilde{R}_t \gets M_t^R / (\sqrt{V_t^R}+\epsilon)$
    \If{APOLLO}
        \State $S_t \gets \operatorname{diag}(s_1^R, \ldots, s_m^R)$, donde $s_i^R = \|\widetilde{R}_t[i,:]\|_2 / \|R_t[i,:]\|_2$
    \Else
        \State $S_t \gets s_t^R$, donde $s_t^R = \|\widetilde{R}_t\|_2 / \|R_t\|_2$
    \EndIf
    \State Actualizar pesos: $W_t \gets (1-\eta\lambda)W_{t-1} - \eta\,\alpha\,(G_t \odot S_t)$
    \State $t \gets t+1$
\Until{convergencia}
\end{algorithmic}
\end{algorithm}

\paragraph{APOLLO-Mini.}
APOLLO-Mini es el miembro de la familia optimizado para \textbf{eficiencia de memoria extrema}. El artículo lo motiva argumentando que, en un espacio compacto de rango 1, el escalado por canales se vuelve demasiado ruidoso, por lo que la actualización se simplifica aún más a un único escalar por tensor. La pérdida de granularidad se compensa parcialmente con el factor de escala explícito $\alpha$ (el artículo discute valores como $128$ en este régimen), proporcionando un punto útil en la frontera de Pareto de optimización: estado muy pequeño, sin sobrecarga de SVD, y aún un comportamiento sólido en preentrenamiento.

\paragraph{Q-APOLLO.}
El artículo de APOLLO también enfatiza que la familia se combina naturalmente con \textbf{cuantización} para entrenamiento con memoria ultra-baja. Este repositorio convierte esa observación de sistemas en una variante concreta del optimizador, \texttt{q\_apollo}. En la implementación, los momentos de primer y segundo orden de bajo rango de APOLLO se almacenan en forma cuantizada junto con escalas y desplazamientos por tensor, y se descuantizan solo cuando es necesario para el siguiente paso. Q-APOLLO preserva, por tanto, la lógica de gradiente proyectado de APOLLO mientras reduce la precisión del estado restante del optimizador, convirtiéndolo en el miembro más agresivo en ahorro de memoria de la pila local de optimizadores.

\paragraph{Anon (Adaptivity Non-restricted Optimizer with Novel convergence technique).}
Anon \citep{anon2026anon} introduce adaptividad ajustable continuamente mediante un único parámetro $\gamma \in \mathbb{R}$, cerrando la brecha entre comportamientos tipo SGD ($\gamma=0$) y tipo Adam ($\gamma=1$), y extrapolando más allá de ambos. La innovación central es el mecanismo Incremental Delay Update (IDU), que permite convergencia estable a través de todo el espectro de adaptividad de valores reales:

\begin{align}
m_t &= \beta_1 m_{t-1} + (1-\beta_1) g_t \quad\text{(momento de primer orden)}\\
v_t &= \beta_2 v_{t-1} + (1-\beta_2) g_t^2 \quad\text{(momento de segundo orden)}\\
\bar{v}_t &= \frac{v_t}{1-\beta_2^{\max(t,1)}} + \epsilon \quad\text{(momento de segundo orden con corrección de sesgo)}\\
\text{fixed\_lr}_k &= \begin{cases}
\bar{v}_k^{-\gamma/2} & \text{si } k = 0 \\
\sqrt{2 / \left(\frac{1}{\text{fixed\_lr}_{k-1}^2} + \bar{v}_k^\gamma\right)} & \text{si } k > 0
\end{cases}\\
\theta_{t+1} &= \theta_t - \frac{\eta}{1-\beta_1^t} \cdot m_t \odot \text{fixed\_lr}_{\lfloor\log_2 t\rfloor}
\end{align}

A diferencia de los optimizadores adaptativos estándar, IDU actualiza el precondicionador solo en potencias de dos ($t = 2^k$), acumulando estadísticas de gradiente entre actualizaciones y agregándolas recursivamente. Este acumulador suave y multiescala evita el seguimiento estricto del máximo de AMSGrad mientras estabiliza demostrablemente la convergencia para cualquier $\gamma \in \mathbb{R}$. El parámetro $\gamma$ controla directamente la adaptividad: $\gamma > 1$ acelera la salida de puntos de silla (favorable para arquitecturas complejas), $\gamma = 1$ recupera el comportamiento tipo Adam, $\gamma = 0$ produce escalado tipo SGD, y $\gamma < 0$ se sesga hacia mínimos más planos (beneficioso para arquitecturas clásicas como CNNs).

Anon mantiene tres buffers de estado ($m$, $v$, fixed\_lr$) con costo $O(n)$ por paso, situándolo en la misma clase de complejidad que AdamW pero con la flexibilidad adicional de adaptividad ajustable. La técnica IDU converge demostrablemente con regret $O(\sqrt{T})$ para problemas convexos y $O(\ln T / \sqrt{T})$ para entornos no convexos, igualando las garantías de los optimizadores adaptativos convencionales mientras soporta todo el rango de adaptividad de valores reales.

\begin{algorithm}[H]
\caption{Anon (Adaptivity Non-restricted Optimizer con IDU)}
\label{alg:anon}
\begin{algorithmic}[1]
\Require Parámetros iniciales $\theta_0$, tasa de aprendizaje $\eta$, $\beta_1, \beta_2$, tolerancia $\epsilon$, adaptividad $\gamma \in \mathbb{R}$
\Ensure Parámetros actualizados $\theta_T$
\State Inicializar: $m \gets 0$, $v \gets 0$, fixed\_lr $\gets 0$, paso $t \gets 0$, contador IDU $k \gets -1$
\For{$t = 1, 2, \ldots, T$}
    \State $g_t \gets \nabla f(\theta_{t-1})$
    \State $m \gets \beta_1 m + (1-\beta_1) g_t$
    \State $v \gets \beta_2 v + (1-\beta_2) g_t^2$
    \If{$t = 2^k$}
        \State $k \gets k+1$
        \State $\bar{v} \gets v/(1 - \beta_2^{\max(t/2,1)}) + \epsilon$
        \If{$k = 0$}
            \State fixed\_lr $\gets \bar{v}^{-\gamma/2}$
        \Else
            \State fixed\_lr $\gets \sqrt{2 / (1/\text{fixed\_lr}^2 + \bar{v}^\gamma)}$
        \EndIf
        \State $v \gets 0$
    \EndIf
    \State $\theta_t \gets \theta_{t-1} - \frac{\eta}{1-\beta_1^t} \cdot m \odot \text{fixed\_lr}$
\EndFor
\Return $\theta_T$
\end{algorithmic}
\end{algorithm}

\paragraph{Prodigy (Approximating the Distance Estimate).}
Prodigy adapta la escala efectiva del paso mediante una estadística tipo distancia acumulativa, eliminando la necesidad de ajuste explícito de la tasa de aprendizaje. El algoritmo mantiene una estimación de distancia acumulativa:
\begin{align}
u_t &= \frac{g_t}{\sqrt{v_t}+\epsilon} \quad\text{(paso adaptativo no normalizado)} \\
s_t &= s_{t-1} + \langle u_t, \theta_t - \theta_0 \rangle \quad\text{(distancia signada acumulativa)} \\
d_t &= \max(d_{t-1}, d_0 + d_{\text{coef}} \cdot s_t) \quad\text{(estimación de distancia con cota inferior)} \\
\theta_{t+1} &= \theta_t - \eta \cdot d_t \cdot u_t
\end{align}
donde $d_0$ es una cota inicial y $d_{\text{coef}}$ controla qué tan agresivamente se adapta la estimación. La estimación de distancia $d_t$ captura la magnitud combinada de gradientes pasados ponderada por el desplazamiento real de los parámetros, implementando una tasa de aprendizaje efectiva implícita. Este escalado consciente de la distancia logra convergencia robusta en distintas escalas de problemas y tamaños de lote sin requerir una programación de la tasa de aprendizaje. Prodigy reduce la sensibilidad a hiperparámetros estimando los tamaños de paso a partir de la geometría de optimización en lugar del conocimiento previo específico del problema.

\begin{algorithm}[H]
\caption{Prodigy (Approximating the Distance Estimate)}
\label{alg:prodigy}
\begin{algorithmic}[1]
\Require Parámetros iniciales $\theta_0$, tasa de aprendizaje $\eta$, coeficiente $d_{\text{coef}}$, distancia inicial $d_0 > 0$
\Ensure Parámetros actualizados $\theta_T$
\State Inicializar: $s \gets 0$ (distancia signada acumulativa), $d \gets d_0$
\For{$t = 1, 2, \ldots, T$}
    \State $g_t \gets \nabla f(\theta_{t-1})$
    \State $v_t \gets \beta_2 v_{t-1} + (1-\beta_2) g_t^2$ \Comment{Segundo momento}
    \State $u_t \gets g_t / (\sqrt{v_t} + \epsilon)$ \Comment{Adaptativo no normalizado}
    \State Actualizar distancia: $s \gets s + \langle u_t, \theta_t - \theta_0 \rangle$ \Comment{Distancia signada}
    \State $d \gets \max(d, d_0 + d_{\text{coef}} \cdot s)$ \Comment{Distancia con cota inferior}
    \State $\theta_t \gets \theta_{t-1} - \eta \cdot d \cdot u_t$ \Comment{Actualización escalada por distancia}
\EndFor
\Return $\theta_T$
\end{algorithmic}
\end{algorithm}

\subsection{Optimizadores de Segundo Orden, Geométricos y de Ortogonalidad}

\begin{algorithm}[H]
\caption{Shampoo (Precondicionamiento Matricial mediante Productos de Kronecker)}
\label{alg:shampoo}
\begin{algorithmic}[1]
\Require Parámetros iniciales $\theta_0$, tasa de aprendizaje $\eta$, intervalo de eigendescomposición $K$
\Ensure Parámetros actualizados $\theta_T$
\State Inicializar: $L \gets \epsilon I_m$, $R \gets \epsilon I_n$ (matrices de Gram izquierda y derecha)
\For{$t = 1, 2, \ldots, T$}
    \State $G \gets \nabla f(\theta_{t-1})$ \Comment{Gradiente}
    \State Actualizar matrices de Gram: $L \gets L + G G^\top$, $R \gets R + G^\top G$
    \If{$t \mod K == 0$}
        \State $Q_L, \Lambda_L \gets \text{eigh}(L)$ \Comment{Eigendescomposición de $L$}
        \State $Q_R, \Lambda_R \gets \text{eigh}(R)$ \Comment{Eigendescomposición de $R$}
        \State $L^{-1/4} \gets Q_L (\Lambda_L + \epsilon)^{-1/4} Q_L^\top$
        \State $R^{-1/4} \gets Q_R (\Lambda_R + \epsilon)^{-1/4} Q_R^\top$
    \EndIf
    \State $\Delta \theta \gets L^{-1/4} G R^{-1/4}$ \Comment{Gradiente precondicionado}
    \State $\theta_t \gets \theta_{t-1} - \eta \cdot \Delta \theta$
\EndFor
\Return $\theta_T$
\end{algorithmic}
\end{algorithm}

\paragraph{Shampoo (Precondicionamiento Matricial mediante Productos de Kronecker).}
Shampoo es un \textbf{método estructurado de segundo orden} que calcula precondicionadores matriciales a partir de estadísticas de producto exterior con estructura de Kronecker. Para una matriz de parámetros 2D $W \in \mathbb{R}^{m \times n}$:
\begin{align}
L_t &= L_{t-1} + G_t G_t^\top \quad\text{(matriz de Gram izquierda/filas, $m \times m$)} \\
R_t &= R_{t-1} + G_t^\top G_t \quad\text{(matriz de Gram derecha/columnas, $n \times n$)} \\
L_t^{-1/4} &= Q_L (\Lambda_L)^{-1/4} Q_L^\top \quad\text{(mediante eigendescomposición)} \\
R_t^{-1/4} &= Q_R (\Lambda_R)^{-1/4} Q_R^\top \\
\Delta W_t &= L_t^{-1/4} G_t R_t^{-1/4} \quad\text{(gradiente precondicionado)}
\end{align}
Shampoo aproxima el precondicionador de matriz completa de Adagrad $H^{-1/2}$ (donde $H$ es el Hessiano) utilizando estructura factorizada de Kronecker. En lugar de almacenar un precondicionador de $(mn) \times (mn)$, Shampoo mantiene dos matrices más pequeñas ($m \times m$ y $n \times n$), capturando correlaciones entre parámetros dentro y entre grupos. El método ha demostrado ser efectivo a escala (sistemas de producción de Google) y es adecuado para arquitecturas transformer con estructura de alto rango. Las eigendescomposiciones se realizan periódicamente (e.g., cada $K=10$ pasos) para amortizar el costo. Contrapartida: mayor cómputo por paso y eigendescomposición $O(m^3 + n^3)$ periódica frente a mejor condicionamiento y convergencia acelerada.

\paragraph{SOAP (Shampoo with Adam in eigenbasis).}
SOAP es una variante simplificada de Shampoo que desacopla el precondicionamiento del seguimiento de momento. En lugar de álgebra matricial compleja sobre gradientes precondicionados, SOAP ejecuta Adam estándar en la base de eigenvectores de los precondicionadores de Shampoo:
\begin{align}
L_t &= L_{t-1} + G_t G_t^\top, \quad R_t = R_{t-1} + G_t^\top G_t \quad\text{(acumulación de Gram)} \\
Q_L, \Lambda_L &= \text{eigh}(L_t), \quad Q_R, \Lambda_R = \text{eigh}(R_t) \quad\text{(eigendescomposición periódica)} \\
G_t^{\text{rot}} &= Q_L^\top G_t Q_R \quad\text{(rotar gradiente a la base de eigenvectores)} \\
m_t^{\text{rot}} &= \beta_1 m_{t-1}^{\text{rot}} + (1-\beta_1) G_t^{\text{rot}} \\
v_t^{\text{rot}} &= \beta_2 v_{t-1}^{\text{rot}} + (1-\beta_2)(G_t^{\text{rot}})^2 \quad\text{(Adam en el espacio rotado)} \\
U_t^{\text{rot}} &= \frac{m_t^{\text{rot}}}{\sqrt{v_t^{\text{rot}}}+\epsilon}, \quad U_t = Q_L U_t^{\text{rot}} Q_R^\top \quad\text{(rotar de vuelta)}
\end{align}
La idea clave: todo el seguimiento de momento ocurre en la base de eigenvectores bien condicionada, simplificando la estabilidad numérica y el análisis teórico. SOAP combina los beneficios de Shampoo (estructura de curvatura explícita) y Adam (mecánica de momento probada), siendo más limpio y a menudo más estable que Shampoo completo. Las eigendescomposiciones se recalculan cada $K$ pasos, amortizando el costo.

\begin{algorithm}[H]
\caption{SOAP (Shampoo with Adam in Eigenbasis)}
\label{alg:soap}
\begin{algorithmic}[1]
\Require Parámetros iniciales $\theta_0$, tasa de aprendizaje $\eta$, $\beta_1, \beta_2$, intervalo de eigendescomposición $K$
\Ensure Parámetros actualizados $\theta_T$
\State Inicializar: $L \gets \epsilon I_m$, $R \gets \epsilon I_n$, $m \gets 0$, $v \gets 0$
\For{$t = 1, 2, \ldots, T$}
    \State $G \gets \nabla f(\theta_{t-1})$ \Comment{Gradiente}
    \State Actualizar Gram: $L \gets L + G G^\top$, $R \gets R + G^\top G$
    \If{$t \mod K == 0$}
        \State $Q_L, \Lambda_L \gets \text{eigh}(L)$, $Q_R, \Lambda_R \gets \text{eigh}(R)$
    \EndIf
    \State $G^{\text{rot}} \gets Q_L^\top G Q_R$ \Comment{Rotar gradiente a la base de eigenvectores}
    \State $m \gets \beta_1 m + (1-\beta_1) G^{\text{rot}}$ \Comment{Media móvil exponencial (corregir sesgo externamente)}
    \State $v \gets \beta_2 v + (1-\beta_2) (G^{\text{rot}})^2$ \Comment{Segundo momento (corregir sesgo externamente)}
    \State $u \gets m / (\sqrt{v} + \epsilon)$ \Comment{Paso adaptativo en la base de eigenvectores}
    \State $\Delta \theta \gets Q_L u Q_R^\top$ \Comment{Rotar de vuelta al espacio de parámetros}
    \State $\theta_t \gets \theta_{t-1} - \eta \cdot \Delta \theta$
\EndFor
\Return $\theta_T$
\end{algorithmic}
\end{algorithm}

\paragraph{Lion (EvoLved Sign Momentum).}
Lion logra un \textbf{sobrecosto de memoria mínimo} al usar actualizaciones basadas en signo en lugar de escalado adaptativo completo:\begin{align}
c_t &= \beta_1 m_t + (1-\beta_1) g_t \quad\text{(entrada de momento)} \\
\theta_{t+1} &= \theta_t - \eta \left(\text{sign}(c_t) + \lambda\theta_t\right) \quad\text{(actualización con operador signo)} \\
m_{t+1} &= \beta_2 m_t + (1-\beta_2) g_t \quad\text{(acumulación de momento)}
\end{align}
Todas las operaciones de escalado elemento a elemento se reemplazan con \texttt{sign()}, que produce $\{-1, 0, +1\}$. Esto reduce drásticamente la memoria en comparación con métodos tipo Adam: Lion almacena solo el buffer de momento, sin varianza de segundo momento. La contrapartida: las actualizaciones basadas en signo sacrifican la adaptación de tasa de aprendizaje elemento a elemento que hace efectivo a Adam. Lion se posiciona no como un optimizador universalmente superior sino como una alternativa especializada de baja memoria y alto rendimiento para escenarios donde la memoria de activaciones domina y se puede sacrificar algo de sofisticación del optimizador. Funciona bien en regímenes de lotes grandes y cuando el rendimiento del hardware es la restricción principal.

\begin{algorithm}[H]
\caption{Lion (EvoLved Sign Momentum)}
\label{alg:lion}
\begin{algorithmic}[1]
\Require Parámetros iniciales $\theta_0$, tasa de aprendizaje $\eta$, $\beta_1, \beta_2$, decaimiento de pesos $\lambda$
\Ensure Parámetros actualizados $\theta_T$
\State Inicializar: $m \gets 0$ (buffer de momento)
\For{$t = 1, 2, \ldots, T$}
    \State $g \gets \nabla f(\theta_{t-1})$ \Comment{Gradiente}
    \State $c \gets \beta_1 m + (1-\beta_1) g$ \Comment{Entrada de momento}
    \State $\theta_t \gets \theta_{t-1} - \eta (\text{sign}(c) + \lambda \theta_{t-1})$ \Comment{Actualización basada en signo con decaimiento de pesos}
    \State $m \gets \beta_2 m + (1-\beta_2) g$ \Comment{Acumulación de momento (desacoplada)}
\EndFor
\Return $\theta_T$
\end{algorithmic}
\end{algorithm}

\paragraph{Sophia (Second-Order Hessian Information with Optimized Approximation).}
Sophia utiliza \textbf{estimaciones diagonales del Hessiano} para escalado consciente de la curvatura sin el costo de memoria y cómputo de los métodos densos de segundo orden.\begin{align}
m_t &= \beta_1 m_{t-1} + (1-\beta_1) g_t \quad\text{(momento de primer orden)} \\
h_t &= \beta_2 h_{t-(k-1)} + (1-\beta_2) \hat{h}_t \quad\text{(Hessiano diagonal, actualizado cada $k$ pasos)} \\
\text{clip}(x, C) &= \min(\max(x, -C), C) \quad\text{(recorte elemento a elemento)} \\
\theta_{t+1} &= \theta_t - \eta \cdot \text{clip}\left(\frac{m_t}{\max(\gamma h_t, \epsilon)}, 1\right)
\end{align}
donde $\hat{h}_t$ es una estimación diagonal (e.g., estimador de traza de Hutchinson) y $\gamma$ es un factor de escala. El Hessiano diagonal $h_t$ captura la curvatura local, permitiendo que el optimizador tome pasos más pequeños en direcciones pronunciadas y pasos más grandes en direcciones planas. La operación de recorte $\text{clip}(\cdot, 1)$ evita que los pasos adaptativos exploten. Sophia pertenece a la familia de métodos conscientes de la curvatura que buscan mejor condicionamiento sin el costo $O(n^2)$ o $O(n^3)$ de los métodos completos de segundo orden. Funciona particularmente bien en escenarios de ajuste fino de segunda pasada.

\begin{algorithm}[H]
\caption{Sophia (Hessiano Diagonal con Actualizaciones Recortadas)}
\label{alg:sophia}
\begin{algorithmic}[1]
\Require Parámetros iniciales $\theta_0$, tasa de aprendizaje $\eta$, $\beta_1, \beta_2$, frecuencia de actualización del Hessiano $k$, umbral de recorte $C$ 
\Ensure Parámetros actualizados $\theta_T$
\State Inicializar: $m \gets 0$, $h \gets \epsilon$ (estimación diagonal del Hessiano)
\For{$t = 1, 2, \ldots, T$}
    \State $g \gets \nabla f(\theta_{t-1})$
    \State $m \gets \beta_1 m + (1-\beta_1) g$ \Comment{Momento de primer orden (corregir sesgo externamente)}
    \If{$t \mod k == 0$}
        \State $\hat{h} \gets \text{HutchinsonEstimate}(g)$ \Comment{Hessiano diagonal mediante Hutchinson}
        \State $h \gets \beta_2 h + (1-\beta_2) \hat{h}$ \Comment{Actualizar estimación del Hessiano}
    \EndIf
    \State $u \gets \frac{m}{\max(\gamma h, \epsilon)}$ \Comment{Paso adaptativo (elemento a elemento)}
    \State $u \gets \text{clip}(u, C)$ \Comment{Recorte elemento a elemento a $[-C, C]$}
    \State $\theta_t \gets \theta_{t-1} - \eta \cdot u$
\EndFor
\Return $\theta_T$
\end{algorithmic}
\end{algorithm}

\paragraph{Muon y Turbo-Muon (Optimizadores Basados en Ortogonalidad).}
Muon y Turbo-Muon son \textbf{optimizadores orientados a la ortogonalidad} que reformulan la geometría de actualización utilizando iteraciones polinomiales de Newton-Schulz para ortogonalización. Para una matriz de parámetros 2D $W$ con gradiente $G_t$:
\begin{align}
X_0 &= \frac{G_t}{\|G_t\|_F + \epsilon} \quad\text{(gradiente normalizado)} \\
A_k &= X_k X_k^\top \quad\text{(Gramiana)} \\
B_k &= b A_k + c A_k^2 \quad\text{(paso polinomial, coeficientes $b, c$ de Newton-Schulz)} \\
X_{k+1} &= a X_k + B_k X_k \quad\text{(iteración $k=0,1,\ldots,4$)} \\
W_{t+1} &= W_t - \eta X_K \quad\text{(actualización con dirección ortogonalizada)}
\end{align}
Muon realiza 5 iteraciones de Newton-Schulz para producir una dirección aproximadamente ortogonal, asegurando que las actualizaciones respeten restricciones geométricas. Turbo-Muon añade un paso de \textbf{precondicionamiento casi-ortogonal} (AOL) antes de las iteraciones, reduciendo el número de pasos de ortogonalización requeridos a 4. La justificación: las actualizaciones ortogonales preservan las normas durante el entrenamiento, evitando la deriva de norma observada en métodos basados en momento. Estos optimizadores muestran potencial en entrenamiento a gran escala pero requieren kernels CUDA personalizados para eficiencia. Contrapartida: alto cómputo por paso (multiplicaciones de matrices e iteraciones polinomiales) frente a una geometría de actualización fundamentalmente mejor condicionada.

\begin{algorithm}[H]
\caption{Muon (Ortogonalización de Newton-Schulz)}
\label{alg:muon}
\begin{algorithmic}[1]
\Require Parámetros iniciales $\theta$ (matrices), tasa de aprendizaje $\eta$, iteraciones Newton-Schulz $K$
\Ensure Parámetros actualizados $\theta$
\For{cada matriz de parámetros 2D $W$}
    \State $G \gets \nabla f(W)$ \Comment{Gradiente}
    \State $X_0 \gets \frac{G}{\|G\|_F + \epsilon}$ \Comment{Normalizar gradiente por norma de Frobenius}
    \For{$k = 0, 1, \ldots, K-1$}
        \State $A_k \gets X_k X_k^\top$ \Comment{Gramiana (costo: $O(mn^2)$ o $O(m^2n)$ para matrices altas/anchas)}
        \State $B_k \gets (3/2) A_k - (1/2) A_k^2$ \Comment{Coeficientes de Newton-Schulz}
        \State $X_{k+1} \gets B_k X_k$ \Comment{Paso polinomial}
    \EndFor
    \State $W \gets W - \eta X_K$ \Comment{Actualizar con dirección ortogonalizada}
\EndFor
\Return $\theta$ actualizado
\end{algorithmic}
\end{algorithm}

\begin{algorithm}[H]
\caption{Turbo-Muon (Ortogonalización Precondicionada con AOL)}
\label{alg:turbo-muon}
\begin{algorithmic}[1]
\Require Parámetros iniciales $\theta$ (matrices), tasa de aprendizaje $\eta$, iteraciones Newton-Schulz $K'$ (típicamente 4)
\Ensure Parámetros actualizados $\theta$
\For{cada matriz de parámetros 2D $W$}
    \State $G \gets \nabla f(W)$ \Comment{Gradiente}
    \State $X_0 \gets \frac{G}{\|G\|_F + \epsilon}$ \Comment{Normalizar gradiente}
    \State \bf{AOL (Almost-Orthogonal preconditioner):} \Comment{Paso de precondicionamiento}
    \State $P \gets (1.5 I - 0.5 X_0 X_0^\top)$ \Comment{Matriz de precondicionamiento}
    \State $X_0 \gets P X_0$ \Comment{Inicio precondicionado}
    \For{$k = 0, 1, \ldots, K'-1$}
        \State $A_k \gets X_k X_k^\top$
        \State $B_k \gets (3/2) A_k - (1/2) A_k^2$
        \State $X_{k+1} \gets B_k X_k$ \Comment{Iteraciones reducidas gracias al precondicionamiento AOL}
    \EndFor
    \State $W \gets W - \eta X_{K'}$ \Comment{Actualizar con dirección ortogonalizada precondicionada}
\EndFor
\Return $\theta$ actualizado
\end{algorithmic}
\end{algorithm}

\small
\begin{longtable}{p{2.5cm}p{3.5cm}p{2.0cm}p{5.3cm}}
\toprule
\textbf{Grupo} & \textbf{Métodos} & \textbf{Objetivo Principal} & \textbf{Interpretación desde el Estudio} \\
\midrule
\endhead
Línea base clásica & SGD, AdamW, RAdam & estabilidad y referencias base & Definen el piso de comparación para las afirmaciones de optimizadores más nuevos. \\
Rediseño de momento & Adan, AdEMAMix, MARS, Cautious AdamW & adaptación de primer orden más rápida o segura & Mejor cuando la velocidad de convergencia o la estabilidad ante gradiente ruidoso es la preocupación principal. \\
Simplificación de lotes grandes y programación & LAMB, Schedule-Free AdamW & robustez operacional a escala & Reducen la fragilidad debida al crecimiento del tamaño de lote o la ingeniería de programación. \\
Eficiencia de memoria & Adafactor, GaLore, APOLLO, APOLLO-Mini, Q-APOLLO, Lion & reducción del estado del optimizador & Más útil cuando la VRAM está dominada por el estado del optimizador en lugar de las activaciones; los métodos de la familia APOLLO reemplazan los momentos densos de AdamW con escalado estructurado proyectado o cuantizado. \\
Consciente de curvatura & Shampoo, SOAP, Sophia & mejor condicionamiento & Preferir cuando una geometría más rica justifica la sobrecarga de implementación y cómputo. \\
Orientado a geometría & Muon, Turbo-Muon & estructura de actualización ortogonalizada & Opciones especializadas para geometría matricial y modelado de representaciones. \\
\bottomrule
\end{longtable}
\normalsize

\section{Anexo B: Transformadores Densos, Recurrentes y Aumentados con Memoria}
\label{sec:annex-transformer}

Este anexo exhaustivo sintetiza las arquitecturas modernas de transformers más allá de la línea base estándar de softmax. El campo ha evolucionado hacia seis paradigmas principales: atención densa global con variantes, compresión de estado recurrente con decaimiento, modelos de espacio de estado selectivos, integración numérica de profundidad continua, arquitecturas aumentadas con memoria y enfoques híbridos. Comprender esta taxonomía ilumina las compensaciones fundamentales entre expresividad, costo computacional, huella de memoria y simplicidad de despliegue.

\subsection{Atención Densa de Referencia: Estándar y Sigmoide}

\subsubsection{Atención Softmax Estándar}
La atención multi-cabeza estándar \citep{vaswani_attention_2017} calcula la similitud de producto punto escalado entre embeddings de tokens:

\begin{algorithmic}[1]
\State \textbf{Entrada:} Matriz de consulta $\mathbf{Q} \in \mathbb{R}^{n \times d}$, Matriz de clave $\mathbf{K} \in \mathbb{R}^{n \times d}$, Matriz de valor $\mathbf{V} \in \mathbb{R}^{n \times d}$
\State $\text{puntajes} \gets \frac{\mathbf{Q}\mathbf{K}^\top}{\sqrt{d}} \in \mathbb{R}^{n \times n}$ \Comment{calcular similitudes por pares}
\State $\text{pesos\_atencion} \gets \texttt{softmax}(\text{puntajes}, \text{eje}=1) \in \mathbb{R}^{n \times n}$ \Comment{normalizar sobre las claves}
\State \textbf{Salida:} $\mathbf{Y} = \text{pesos\_atencion} \cdot \mathbf{V} \in \mathbb{R}^{n \times d}$ \Comment{combinación ponderada de valores}
\end{algorithmic}

Para generación autorregresiva, el almacenamiento en caché KV guarda claves y valores pasados para evitar el recálculo $\mathcal{O}(n^2)$. Sin embargo, este caché de crecimiento lineal ocupa $\sim n \cdot d_{\text{oculto}}$ bytes, lo que puede consumir cientos de gigabytes para modelos de miles de millones de parámetros. La atención estándar logra una \textbf{expresividad perfecta} dentro de la ventana de contexto---cualquier token puede atender a cualquier otro token con pesos aprendidos---pero paga el precio del cómputo denso.

\begin{itemize}[leftmargin=*]
    \item \textbf{Complejidad de entrenamiento}: $\mathcal{O}(n^2 \cdot d)$ tiempo, $\mathcal{O}(n^2)$ espacio (materialización de la matriz de atención).
    \item \textbf{Complejidad de inferencia}: $\mathcal{O}(n)$ tiempo por token, $\mathcal{O}(n)$ espacio (caché KV).
    \item \textbf{Fortalezas}: Expresividad incomparable; recuperación perfecta del historial; altamente paralelizable.
    \item \textbf{Debilidades}: El cuello de botella cuadrático prohíbe contextos muy largos; el caché KV domina la memoria durante la generación.
\end{itemize}

\subsubsection{Atención Sigmoide}
La atención sigmoide reemplaza el softmax por filas con una activación sigmoide elemento a elemento \citep{ramapuram_theory_2024}:

\begin{algorithmic}[1]
\State \textbf{Entrada:} $\mathbf{Q}, \mathbf{K}, \mathbf{V}$ como arriba, sesgo aprendible $\mathbf{b} \in \mathbb{R}^{n \times n}$
\State $\text{logits} \gets \frac{\mathbf{Q}\mathbf{K}^\top}{\sqrt{d}} + \mathbf{b}$
\State $\text{pesos\_atencion} \gets \sigma(\text{logits})$ \Comment{sigmoide elemento a elemento, no softmax}
\State \textbf{Salida:} $\mathbf{Y} = \text{pesos\_atencion} \odot \mathbf{V}$
\end{algorithmic}

A diferencia de softmax, la sigmoide no impone una distribución de probabilidad (los pesos no necesitan sumar 1), lo que permite una independencia de tokens más fuerte. El análisis teórico mediante mezcla de expertos muestra que la sigmoide alcanza una complejidad muestral superior: convergencia $\mathcal{O}(n^{-0.51})$ para expertos ReLU frente a $\mathcal{O}(n^{-0.24})$ de softmax. Sin embargo, el entrenamiento empírico reveló inestabilidades de gradiente a escala. El remedio es la \textbf{norma híbrida}---agregar normalización después de la salida de atención---que estabiliza los gradientes sin sacrificar los beneficios teóricos del enrutamiento elemento a elemento.

\begin{itemize}[leftmargin=*]
    \item \textbf{Complejidad de entrenamiento}: $\mathcal{O}(n^2 \cdot d)$ (idéntico al estándar), pero las operaciones elemento a elemento permiten una aceleración de inferencia del 17\% mediante FlashSigmoid.
    \item \textbf{Complejidad de inferencia}: $\mathcal{O}(n)$ por token con caché KV (asintóticamente igual, pero con factores constantes más bajos).
    \item \textbf{Fortalezas}: Supera la competencia de suma cero; evita la sincronización por filas; implementación eficiente en hardware.
    \item \textbf{Debilidades}: Requiere estabilización cuidadosa (norma híbrida); inestabilidad de entrenamiento a grandes escalas sin pérdida auxiliar.
\end{itemize}

\subsection{Arquitecturas Recurrentes y Retentivas}

\subsubsection{Redes Retentivas (RetNet)}
RetNet \citep{sun_retentive_2023} unifica tres paradigmas de cómputo: entrenamiento paralelo, inferencia recurrente y despliegue por fragmentos. Su innovación central es el \textbf{mecanismo de retención}, que utiliza una matriz de decaimiento exponencial fija para modelar la importancia temporal:

\begin{algorithmic}[1]
\State \textbf{Forma paralela (entrenamiento):}
\State $\text{matriz\_decaimiento}[i,j] \gets \gamma^{i-j}$ para $i \geq j$, si no $0$ \Comment{decaimiento exponencial causal}
\State $\text{matriz\_decaimiento}[i,j] \gets 0$ para $i < j$ \Comment{enmascaramiento causal}
\State $\mathbf{Y}_{\text{paralelo}} \gets (\mathbf{Q}\mathbf{K}^\top \odot \text{matriz\_decaimiento}) \mathbf{V}$ \Comment{multiplicación elemento a elemento con decaimiento}
\end{algorithmic}

\begin{algorithmic}[1]
\State \textbf{Forma recurrente (inferencia):}
\For{$t = 1, \ldots, n$}
    \State $\mathbf{s}_t \gets \gamma \mathbf{s}_{t-1} + \mathbf{k}_t \mathbf{v}_t^\top$ \Comment{actualización de estado con decaimiento}
    \State $\mathbf{y}_t \gets \mathbf{q}_t \mathbf{s}_t$ \Comment{salida mediante interacción consulta-estado}
\EndFor
\end{algorithmic}

El escalar de decaimiento $\gamma \in (0,1)$ controla la ventana temporal. RetNet utiliza \textbf{retención multi-escala} con diferentes valores de $\gamma$ por cabeza (ej., $\gamma = 1 - 2^{-5}, 1 - 2^{-6}, \ldots$), permitiendo dependencias a corto y largo plazo simultáneamente. El modo recurrente por fragmentos divide las secuencias en fragmentos, procesa cada fragmento en paralelo y enhebra un estado recurrente entre fragmentos.

\begin{itemize}[leftmargin=*]
    \item \textbf{Complejidad de entrenamiento}: $\mathcal{O}(n^2 \cdot d)$ (paralela), o $\mathcal{O}(n \cdot c \cdot d)$ (recurrente por fragmentos con tamaño de fragmento $c$).
    \item \textbf{Complejidad de inferencia}: $\mathcal{O}(1)$ por token, $\mathcal{O}(d^2)$ espacio de estado (matriz fija).
    \item \textbf{Fortalezas}: Inferencia en tiempo constante; triple paradigma de cómputo; consolidación multi-escala.
    \item \textbf{Debilidades}: El decaimiento fijo impone un sesgo inductivo rígido; puede truncar patrones de largo alcance aprendidos.
\end{itemize}

\subsubsection{Mamba: Modelos de Espacio de Estado Selectivos}
Mamba \citep{gu_mamba_2023} enmarca la recurrencia como un sistema dinámico de tiempo continuo con parámetros dependientes de la entrada, logrando tanto entrenamiento lineal como inferencia en tiempo constante:

\begin{algorithmic}[1]
\State \textbf{Dinámica continua:} $h'(t) = \mathbf{A} h(t) + \mathbf{B} x(t)$, \quad $y(t) = \mathbf{C} h(t)$
\State \textbf{Discretización con tamaño de paso $\Delta_t$:}
\State $\bar{\mathbf{A}}_t \gets \exp(\Delta_t \mathbf{A})$
\State $\bar{\mathbf{B}}_t \gets (\Delta_t \mathbf{A})^{-1} (\exp(\Delta_t \mathbf{A}) - \mathbf{I}) \Delta_t \mathbf{B}_t$
\State \textbf{Actualización recurrente:}
\For{$t = 1, \ldots, n$}
    \State $\Delta_t \gets \text{softplus}(\text{Lineal}(x_t))$ \Comment{tamaño de paso dependiente de la entrada}
    \State $\mathbf{B}_t \gets \text{Lineal}(x_t)$, \quad $\mathbf{C}_t \gets \text{Lineal}(x_t)$ \Comment{matrices de proyección dependientes de la entrada}
    \State $h_t \gets \bar{\mathbf{A}}_t h_{t-1} + \bar{\mathbf{B}}_t x_t$ \Comment{transición de estado}
    \State $y_t \gets \mathbf{C}_t h_t$ \Comment{proyección de salida}
\EndFor
\end{algorithmic}

La innovación crítica es que $\mathbf{A}$ (la matriz del sistema) \textit{no} depende de la entrada, pero $\Delta_t$, $\mathbf{B}_t$ y $\mathbf{C}_t$ sí, haciendo que el sistema sea \textbf{variante en el tiempo}. Esta selectividad permite al modelo \textit{ignorar} información irrelevante estableciendo $\Delta_t$ cerca de cero, creando efectivamente una compuerta. Un algoritmo de escaneo paralelo consciente del hardware implementa la recurrencia eficientemente en GPUs fusionando el cómputo dentro de SRAM, evitando el costoso ancho de banda de HBM.

\begin{itemize}[leftmargin=*]
    \item \textbf{Complejidad de entrenamiento}: $\mathcal{O}(n \cdot d)$ mediante escaneo consciente del hardware (lineal en la longitud de la secuencia).
    \item \textbf{Complejidad de inferencia}: $\mathcal{O}(1)$ por token, $\mathcal{O}(d)$ estado (vector oculto).
    \item \textbf{Fortalezas}: Logra entrenamiento lineal e inferencia constante simultáneamente; práctico en secuencias largas (millones de tokens).
    \item \textbf{Debilidades}: La compresión del vector de estado puede debilitar la copia exacta y la recuperación asociativa densa en comparación con la atención completa.
\end{itemize}

\subsection{Transformers de Profundidad Continua: Integración EDO}

\subsubsection{Transformer EDO}
El Transformer EDO \citep{zhang_continuous_2021} interpreta la profundidad de la red como integración numérica de un sistema dinámico continuo, utilizando solucionadores Runge-Kutta de orden superior para reducir el error de truncamiento:

\begin{algorithmic}[1]
\State \textbf{Formulación continua:} $\frac{dh(t)}{dt} = f_\theta(h(t), t)$ donde $f_\theta$ es la subred del transformer
\State \textbf{Aproximación discreta Runge-Kutta-4:}
\State $k_1 \gets f_\theta(h_t, t)$
\State $k_2 \gets f_\theta(h_t + \tfrac{1}{2}k_1, t + \tfrac{1}{2}\Delta t)$
\State $k_3 \gets f_\theta(h_t + \tfrac{1}{2}k_2, t + \tfrac{1}{2}\Delta t)$
\State $k_4 \gets f_\theta(h_t + k_3, t + \Delta t)$
\State $h_{t+1} \gets h_t + \frac{\Delta t}{6}(k_1 + 2k_2 + 2k_3 + k_4)$ \Comment{combinación ponderada}
\end{algorithmic}

En lugar de las conexiones residuales simples de Euler $h_{t+1} = h_t + f_\theta(h_t)$, el bloque RK4 calcula cuatro evaluaciones intermedias y las combina con los pesos clásicos de RK4. Para evitar gradientes que se desvanecen, la arquitectura introduce \textbf{compuertas aprendidas} que interpolar entre aproximaciones intermedias:

\begin{algorithmic}[1]
\State $g \gets \sigma(\text{Lineal}([k_1, k_2, k_3, k_4]))$ \Comment{compuerta aprendible}
\State $h_{t+1} \gets h_t + g \cdot k_1 + (1-g) \cdot k_2$ \Comment{refinamiento interpolado}
\end{algorithmic}

Esta formulación reduce el número efectivo de parámetros mediante el uso compartido de pesos---la misma $f_\theta$ se evalúa múltiples veces---proporcionando un refinamiento de trayectoria más rico.

\begin{itemize}[leftmargin=*]
    \item \textbf{Complejidad de entrenamiento}: $\mathcal{O}(k \cdot n^2 \cdot d)$ donde $k$ es el orden RK (4 para RK4).
    \item \textbf{Complejidad de inferencia}: $\mathcal{O}(k \cdot n)$ por token (mayor sobrecarga constante).
    \item \textbf{Fortalezas}: Precisión significativamente mayor en tareas de generación (BLEU de última generación); eficiente en parámetros mediante uso compartido de pesos.
    \item \textbf{Debilidades}: Alto costo de cómputo por paso; la latencia de inferencia aumenta por un factor constante $k$; requiere compuertas complejas para la estabilidad.
\end{itemize}

\subsection{Memoria en Tiempo de Prueba: Titans}

La arquitectura Titans \citep{behrouz_titans_2025} introduce una dimensión ortogonal: en lugar de pesos estáticos, el modelo mantiene una memoria aprendible que se \textit{actualiza durante la inferencia} basada en una señal impulsada por sorpresa:

\begin{algorithmic}[1]
\State \textbf{Atención local a corto plazo:} Aplicar atención estándar o dispersa dentro de una ventana de contexto fija $c$
\State $y_t^{\text{local}} \gets \text{Atencion}(q_t, k_{[t-c:t]}, v_{[t-c:t]})$
\State \textbf{Señal de actualización de memoria (sorpresa):}
\State $S_t \gets \eta_t S_{t-1} - \theta_t \nabla_\ell(\mathcal{M}_{t-1}; x_t)$ \Comment{sorpresa como gradiente de la pérdida respecto a la memoria}
\State $\mathcal{M}_t \gets (1 - \alpha_t) \mathcal{M}_{t-1} + S_t$ \Comment{memoria actualizada mediante EMA}
\State \textbf{Recuperación de memoria a largo plazo:}
\State $y_t^{\text{memoria}} \gets \mathcal{M}_t^* (q_t)$ \Comment{consultar módulo de memoria para información recuperada}
\State \textbf{Combinación de salida:}
\State $y_t \gets y_t^{\text{local}} + \text{compuerta}(y_t^{\text{memoria}})$ \Comment{combinar memoria local y recuperada}
\end{algorithmic}

El módulo de memoria $\mathcal{M}$ se actualiza literalmente durante el paso hacia adelante calculando gradientes de una pérdida asociativa y aplicando pasos de SGD con momento. Las tasas de decaimiento $\eta_t, \theta_t, \alpha_t$ son en sí mismas dependientes de la entrada, permitiendo al modelo cambiar de paradigma de memoria cuando el contexto cambia.

\begin{itemize}[leftmargin=*]
    \item \textbf{Complejidad de entrenamiento}: Aproximadamente $\mathcal{O}(c^2 + n \cdot f)$ donde $c$ es la ventana local y $f$ es la sobrecarga de actualización de memoria.
    \item \textbf{Complejidad de inferencia}: $\mathcal{O}(c^2)$ atención local más $\mathcal{O}(1)$ recuperación de memoria por token.
    \item \textbf{Fortalezas}: Maneja longitudes de contexto extremas; permite recuperación asociativa real; la memoria se adapta a la entrada.
    \item \textbf{Debilidades}: Significativamente más complejo; la inferencia incluye cálculo de gradientes; mayor sobrecarga de coordinación.
\end{itemize}

\subsection{Comparación y Síntesis Arquitectónica}

\small
\begin{longtable}{p{2.0cm}p{1.8cm}p{1.6cm}p{1.6cm}p{4.5cm}}
\toprule
\textbf{Arquitectura} & \textbf{Entrenar} & \textbf{Inferir} & \textbf{Estado} & \textbf{Característica Clave} \\
\midrule
\endhead
Atención Estándar & $O(n^2 d)$ & $O(n)$ caché & $O(nd)$ & Expresividad perfecta, enrutamiento completo de tokens, cuello de botella KV. \\
Atención Sigmoide & $O(n^2 d)$ & $O(n)$ caché & $O(nd)$ & Compuerta elemento a elemento, eficiente en hardware, estable a escala con norma híbrida. \\
RetNet & $O(n^2 d)$ & $O(1)$ & $O(d^2)$ & Decaimiento multi-escala, triple modo de cómputo, patrón de olvido fijo. \\
Mamba & $O(nd)$ & $O(1)$ & $O(d)$ & Dinámica selectiva dependiente de la entrada, entrenamiento e inferencia lineales, escaneo de hardware. \\
Transformer EDO & $O(kn^2 d)$ & $O(kn)$ & $O(nd)$ & Refinamiento por integración numérica, uso compartido de pesos por etapas, mayor precisión. \\
Titans & $O(c^2 + nf)$ & $O(c^2 + 1)$ & $O(1)$ & Adaptación de memoria en tiempo de prueba, contexto extremo, actualizaciones de parámetros durante inferencia. \\
\bottomrule
\end{longtable}
\normalsize

El campo exhibe una progresión clara a lo largo de dos ejes: (i) \textbf{eficiencia computacional}, pasando de $O(n^2)$ a $O(n)$ u $O(1)$, y (ii) \textbf{adaptabilidad de memoria}, pasando de pesos estáticos a modelos dinámicos actualizados en tiempo de prueba. La atención estándar sigue siendo la línea base de expresividad; Mamba y RetNet representan la frontera práctica de eficiencia; Titans introduce un eje de innovación ortogonal (aprendizaje en tiempo de prueba). La elección de arquitectura refleja la restricción fundamental de ingeniería: expresividad versus costo de despliegue.

\section{Annex C: Comprehensive Sparse Attention Mechanisms}
\label{sec:annex-sparse}

\subsection{Resumen Ejecutivo}

Los mecanismos de atención dispersa abordan la prohibitiva complejidad computacional y de memoria $\mathcal{O}(n^2)$ de la atención de producto punto escalado estándar al restringir el conjunto de pares clave-valor atendidos. La atención dispersa moderna abarca un rico espacio de diseño distinguido por cuatro dimensiones ortogonales: (i) \textbf{patrón de dispersidad} (geométrico fijo, dependiente de datos o basado en frecuencia), (ii) \textbf{capacidad de entrenamiento} (entrenado de extremo a extremo o solo inferencia), (iii) \textbf{unidad de dispersidad} (tokens individuales, ventanas locales o bloques), y (iv) \textbf{mecanismo} (enmascaramiento, selección, filtrado o compresión). Este anexo sintetiza siete familias principales de atención dispersa---Sparse Transformer, Longformer, BigBird, SparseK, NSA, SpargeAttn y FASA---proporcionando formulaciones matemáticas, pseudocódigo algorítmico y comparaciones arquitectónicas exhaustivas.

\subsection{Sparse Transformer: Patrones Estrideados y Fijos Factorizados}

\subsubsection{Formulación Matemática}

Sparse Transformer \citep{child_sparse_transformer_2019} reduce la complejidad cuadrática a $\mathcal{O}(n\sqrt{n})$ factorizando la atención dispersa en dos cabezas dispersas complementarias. Sea $n$ la longitud de la secuencia y $l = \lfloor\sqrt{n}\rfloor$ el stride.

\paragraph{Cabeza de atención estrideada}: Cada posición $i$ atiende a cada $l$-ésima posición anterior:
\[
\mathcal{A}_i^{(\text{estrideada})} = \{j : j \leq i,\; (i - j) \bmod l = 0\}
\]

\paragraph{Cabeza de atención fija}: Cada posición $i$ atiende a posiciones dentro de su bloque local más columnas de resumen fijas:
\[
\mathcal{A}_i^{(\text{fija})} = \{j : \lfloor j/l \rfloor = \lfloor i/l \rfloor\} \cup \{j : j \bmod l \in \{l-c, \ldots, l-1\}\}
\]
donde $c$ es un hiperparámetro que controla el número de columnas de resumen (típicamente $c=1$ o $c=2$).

El cálculo de atención para cada cabeza $h$ sigue las reglas estándar de producto punto escalado restringido al conjunto disperso:
\[
\text{Attn}_h(Q, K, V)_i = \text{softmax}\left(\frac{q_i^h K_{\mathcal{A}_i}^h\top}{\sqrt{d_k}}\right) V_{\mathcal{A}_i}^h
\]

La idea clave es que con las dos cabezas factorizadas operando en paralelo a través de una profundidad de $L$ capas transformer, cada posición puede alcanzar cualquier otra posición a través de un camino de longitud máxima $L + 1$, preservando la alcanzabilidad de largo alcance con una fracción del costo de la atención densa.

\subsubsection{Pseudocódigo Algorítmico}

\begin{algorithm}[H]
\caption{Atención Sparse Transformer: Cabezas Duales Estrideada + Fija}
\begin{algorithmic}[1]
\Require Consulta $\mathbf{Q} \in \mathbb{R}^{n \times d}$, Clave $\mathbf{K} \in \mathbb{R}^{n \times d}$, Valor $\mathbf{V} \in \mathbb{R}^{n \times d}$
\Require Stride $l = \lfloor\sqrt{n}\rfloor$, columnas de resumen $c \in \{1, 2\}$
\Ensure Salida $\mathbf{Y} \in \mathbb{R}^{n \times d}$
\State \textbf{Cabeza 1: Atención Estrideada}
\For{$i = 1, \ldots, n$}
    \State $\mathcal{A}_i \gets \{j : j \leq i \text{ y } (i-j) \bmod l = 0\}$ \Comment{Cada $l$-ésima posición}
    \State $\text{puntajes}_{i} \gets \mathbf{q}_i \mathbf{K}_{\mathcal{A}_i}^\top / \sqrt{d_k}$
    \State $\text{attn}_{i} \gets \text{softmax}(\text{puntajes}_{i})$
    \State $\mathbf{y}_i^{(1)} \gets \text{attn}_{i} \mathbf{V}_{\mathcal{A}_i}$
\EndFor
\State \textbf{Cabeza 2: Atención de Bloque Fijo + Resumen}
\For{$i = 1, \ldots, n$}
    \State $\text{inicio\_bloque} \gets \lfloor i/l \rfloor \cdot l$, $\text{fin\_bloque} \gets \min(i+1, \text{inicio\_bloque}+l)$
    \State $\mathcal{A}_i \gets [\text{inicio\_bloque}, \text{fin\_bloque}) \cup \{\text{columnas de las últimas } c \text{ columnas}\}$
    \State $\text{puntajes}_{i} \gets \mathbf{q}_i \mathbf{K}_{\mathcal{A}_i}^\top / \sqrt{d_k}$
    \State $\text{attn}_{i} \gets \text{softmax}(\text{puntajes}_{i})$
    \State $\mathbf{y}_i^{(2)} \gets \text{attn}_{i} \mathbf{V}_{\mathcal{A}_i}$
\EndFor
\State Concatenar: $\mathbf{Y} = [\mathbf{y}^{(1)} \| \mathbf{y}^{(2)}]$ y proyectar mediante capa lineal de salida
\Return $\mathbf{Y}$
\end{algorithmic}
\end{algorithm}

\subsubsection{Características Clave}

\begin{itemize}[leftmargin=*]
    \item \textbf{Complejidad}: $\mathcal{O}(n\sqrt{n} \cdot d)$ durante el entrenamiento; $\mathcal{O}(n)$ por token durante la generación.
    \item \textbf{Entrenable}: Sí; retropropagación completa a través de las posiciones seleccionadas.
    \item \textbf{Patrón de dispersidad}: Fijo e independiente de los datos; los patrones no se adaptan al contenido.
    \item \textbf{Fortaleza}: Alcanzabilidad estructurada que garantiza dependencias de largo alcance; demostrado efectivo en secuencias largas (Enwik8, imágenes de texto).
    \item \textbf{Limitación}: Los patrones fijos pueden pasar por alto relaciones importantes pero no contiguas; requiere kernels CUDA personalizados para eficiencia práctica.
\end{itemize}

\subsection{Longformer: Ventana Deslizante, Dilatación y Tokens Globales}

\subsubsection{Formulación Matemática}

Longformer \citep{beltagy_longformer_2020} logra una complejidad lineal $\mathcal{O}(n \cdot w)$ combinando tres patrones de atención complementarios.

\paragraph{Atención de ventana deslizante}: Vecindario local de tamaño $w$:
\[
\mathcal{A}_i^{(\text{ventana})} = \{j : |i - j| \leq w/2\}
\]

\paragraph{Ventana deslizante dilatada}: Atención con ventana y dilatación $d$, saltando stride $d+1$:
\[
\mathcal{A}_i^{(\text{dilatada})} = \{j : |i - j| \leq (w/2)(d+1) \text{ y } (i-j) \bmod (d+1) = 0\}
\]

\paragraph{Atención global}: Tokens globales designados (ej., $[\text{CLS}]$) atienden a todas las posiciones y todos atienden a ellos:
\[
\mathcal{A}_i^{(\text{global})} = \begin{cases}
\{1, \ldots, n\} & \text{si } i \in \mathcal{G} \\
\mathcal{A}_i^{(\text{ventana})} \cup \mathcal{G} & \text{en otro caso}
\end{cases}
\]

Los tres patrones se aplican mediante proyecciones separadas. La atención local y global pueden usar las mismas proyecciones o unas separadas específicas para cada tarea.

\subsubsection{Pseudocódigo Algorítmico}

\begin{algorithm}[H]
\caption{Longformer: Ventana Deslizante + Dilatada + Global}
\begin{algorithmic}[1]
\Require Consulta $\mathbf{Q} \in \mathbb{R}^{n \times d}$, Clave $\mathbf{K}$, Valor $\mathbf{V}$, tamaño de ventana $w$, dilatación $d$, índices de tokens globales $\mathcal{G}$
\Ensure Salida $\mathbf{Y}$
\For{$i = 1, \ldots, n$}
    \If{$i \in \mathcal{G}$}
        \State $\mathcal{A}_i \gets \{1, \ldots, n\}$ \Comment{Token global atiende a todos}
    \Else
        \State Ventana local: $\mathcal{A}^{(\text{local})} \gets [i - w/2, i + w/2] \cap [1, n]$
        \State Desplazamientos dilatados: $\mathcal{A}^{(\text{dilatada})} \gets \{j : (i-j) \bmod (d+1) = 0\} \cap \mathcal{A}^{(\text{rango dilatado})}$
        \State Combinar: $\mathcal{A}_i \gets \mathcal{A}^{(\text{local})} \cup \mathcal{A}^{(\text{dilatada})} \cup \mathcal{G}$
    \EndIf
    \State $\text{puntajes}_i \gets \mathbf{q}_i \mathbf{K}_{\mathcal{A}_i}^\top / \sqrt{d_k}$
    \State $\text{attn}_i \gets \text{softmax}(\text{puntajes}_i)$
    \State $\mathbf{y}_i \gets \text{attn}_i \mathbf{V}_{\mathcal{A}_i}$
\EndFor
\Return $\mathbf{Y}$
\end{algorithmic}
\end{algorithm}

\subsubsection{Características Clave}

\begin{itemize}[leftmargin=*]
    \item \textbf{Complejidad}: $\mathcal{O}(n \cdot w)$ donde $w$ es el tamaño de ventana (lineal cuando $w \ll n$).
    \item \textbf{Entrenable}: Sí; reemplazo directo para la atención estándar.
    \item \textbf{Cobertura}: Con $L$ capas y tamaño de ventana $w$, el campo receptivo crece a $L \times w$, cubriendo toda la secuencia con redes poco profundas.
    \item \textbf{Fortaleza}: Práctico para tareas a nivel de documentos; la dilatación expande los campos receptivos locales con sobrecarga mínima; los tokens globales proporcionan correspondencia consulta-documento.
    \item \textbf{Limitación}: El tamaño de ventana sigue siendo un hiperparámetro de diseño; subóptimo para tareas que requieren patrones de atención densos.
\end{itemize}

\subsection{BigBird: Grafo Disperso Aleatorio, Local y Global}

\subsubsection{Formulación Matemática}

BigBird \citep{zaheer_bigbird_2020} preserva la aproximación universal y la completitud de Turing utilizando una combinación fundamentada de tres patrones de dispersidad.

\paragraph{Conexiones aleatorias}: Cada posición se conecta a $r$ posiciones muestreadas aleatoriamente:
\[
\mathcal{A}_i^{(\text{aleatoria})} = \text{MuestraAleatoria}(\{1, \ldots, n\}, r)
\]

\paragraph{Ventana local}: Vecindario de ventana deslizante:
\[
\mathcal{A}_i^{(\text{local})} = \{j : |i-j| \leq w/2\}
\]

\paragraph{Tokens globales}: Anclas globales específicas de la tarea:
\[
\mathcal{A}_i^{(\text{global})} = \begin{cases}
\{1, \ldots, n\} & \text{si } i \in \mathcal{G} \\
\mathcal{A}_i^{(\text{aleatoria})} \cup \mathcal{A}_i^{(\text{local})} \cup \mathcal{G} & \text{en otro caso}
\end{cases}
\]

La máscara dispersa compuesta $M$ es:
\[
M_{ij} = \mathbf{1}[j \in \mathcal{A}_i^{(\text{aleatoria})} \cup \mathcal{A}_i^{(\text{local})} \cup \mathcal{A}_i^{(\text{global})}]
\]

En la práctica, BigBird agrupa tokens en bloques de tamaño $b$ y aplica la decisión de dispersidad a nivel de bloque, permitiendo implementaciones eficientes de bloques dispersos.

\subsubsection{Garantía Teórica}

Los autores demuestran que con $O(1)$ conexiones aleatorias y tokens globales, el grafo disperso mantiene las mismas propiedades de aproximación universal y completitud de Turing que la atención densa. Formalmente, cualquier función computable puede representarse y cualquier secuencia de operaciones puede simularse dentro del grafo de conectividad de BigBird.

\subsubsection{Pseudocódigo Algorítmico}

\begin{algorithm}[H]
\caption{BigBird: Atención Dispersa por Bloques Aleatoria + Local + Global}
\begin{algorithmic}[1]
\Require Consulta $\mathbf{Q}$, Clave $\mathbf{K}$, Valor $\mathbf{V}$, tamaño de bloque $b$, enlaces aleatorios por bloque $r$, índices globales $\mathcal{G}$
\Ensure Salida $\mathbf{Y}$
\State Reformar en bloques: $n_b = \lceil n / b \rceil$ bloques
\For{bloque $i = 0, \ldots, n_b-1$}
    \For{posición $j$ en bloque $i$}
        \State \textbf{Ventana local}: $\mathcal{A}_j^{(\text{local})} = \text{posiciones cercanas en el bloque} \cup \text{posiciones frontera}$
        \State \textbf{Muestreo aleatorio de bloques}: Muestrear $r$ bloques uniformemente al azar, agregar todas las posiciones de esos bloques
        \State \textbf{Global}: Agregar todas las posiciones de tokens globales
        \State $\mathcal{A}_j \gets \mathcal{A}_j^{(\text{local})} \cup \mathcal{A}_j^{(\text{aleatoria})} \cup \mathcal{G}$
    \EndFor
\EndFor
\For{$i = 1, \ldots, n$}
    \State $\text{puntajes}_i \gets \mathbf{q}_i \mathbf{K}_{\mathcal{A}_i}^\top / \sqrt{d_k}$
    \State $\text{attn}_i \gets \text{softmax}(\text{puntajes}_i)$
    \State $\mathbf{y}_i \gets \text{attn}_i \mathbf{V}_{\mathcal{A}_i}$
\EndFor
\Return $\mathbf{Y}$
\end{algorithmic}
\end{algorithm}

\subsubsection{Características Clave}

\begin{itemize}[leftmargin=*]
    \item \textbf{Complejidad}: $\mathcal{O}(n)$ o casi lineal en implementación óptima de bloques dispersos.
    \item \textbf{Entrenable}: Sí.
    \item \textbf{Base teórica}: Mantiene demostrablemente la aproximación universal y la completitud de Turing con conectividad dispersa.
    \item \textbf{Fortaleza}: Rendimiento sólido en documentos largos para preguntas-respuestas, resúmenes y tareas de genómica.
    \item \textbf{Limitación}: El muestreo aleatorio introduce varianza y no determinismo; el ajuste de los hiperparámetros $r, w, g$ sigue dependiendo de la tarea.
\end{itemize}

\subsection{Atención SparseK: Selección Diferencial Top-$k$}

\subsubsection{Formulación Matemática}

SparseK \citep{lou_sparsek_2024} permite dispersidad aprendible mediante un \textbf{operador top-$k$ diferenciable}. Una red de puntuación $\phi_\theta$ evalúa la importancia de cada par clave-valor:

\[
u_j = \phi_\theta(\mathbf{k}_j, \mathbf{q}_i) \in \mathbb{R}
\]

El \textbf{operador SparseK} selecciona los puntajes top-$k$ manteniéndose diferenciable. Calcula un umbral $\tau(u)$ tal que la suma de los puntajes activos sea igual a $k$:

\[
\text{SparseK}(u, k)_j = \max(u_j - \tau(u), 0), \quad \text{donde} \quad \sum_j \max(u_j - \tau, 0) = k
\]

El umbral $\tau$ se encuentra mediante bisección. La atención se convierte en:

\[
m_j = \text{SparseK}(u, k)_j, \quad \text{Attn}_i = \text{softmax}\left(\frac{\mathbf{q}_i K_{\text{sel}}^\top}{\sqrt{d_k}}\right) V_{\text{sel}}
\]

donde $K_{\text{sel}}, V_{\text{sel}}$ contienen solo las entradas top-$k$ (aquellas con $m_j > 0$).

\subsubsection{Pseudocódigo Algorítmico}

\begin{algorithm}[H]
\caption{SparseK: Atención Top-$k$ Diferenciable}
\begin{algorithmic}[1]
\Require Consulta $\mathbf{q} \in \mathbb{R}^d$, Matriz de clave $\mathbf{K} \in \mathbb{R}^{n \times d}$, Matriz de valor $\mathbf{V} \in \mathbb{R}^{n \times d}$
\Require Red de puntuación $\phi_\theta$, dispersidad objetivo $k$
\Ensure Salida de atención $\mathbf{y}$
\State \textbf{Calcular puntajes de importancia:}
\For{$j = 1, \ldots, n$}
    \State $u_j \gets \phi_\theta(\mathbf{k}_j, \mathbf{q})$
\EndFor
\State \textbf{Calcular top-$k$ diferenciable mediante umbral:}
\State $u_{\text{ordenado}} \gets \text{ordenar}(u, \text{descendente}=\text{Verdadero})$
\State $\text{suma\_acum} \gets \text{suma\_acumulada}(u_{\text{ordenado}})$
\State Encontrar $\rho = \max\{i : u_{\text{ordenado}}[i] > 0 \text{ y suma\_acum}[i] \leq k\}$
\State $\tau \gets (u_{\text{ordenado}}[\rho] - k) / (\rho + 1)$ \Comment{Umbral para $k$ elementos activos}
\State $m \gets \max(u - \tau, 0)$ \Comment{Máscara de selección diferenciable}
\State \textbf{Aplicar máscara y calcular atención:}
\State $K_{\text{sel}} \gets K[m > 0]$, $V_{\text{sel}} \gets V[m > 0]$
\State $\text{puntajes} \gets \mathbf{q} K_{\text{sel}}^\top / \sqrt{d_k}$
\State $\text{attn} \gets \text{softmax}(\text{puntajes})$
\State $\mathbf{y} \gets \text{attn} \cdot V_{\text{sel}}$
\Return $\mathbf{y}$
\end{algorithmic}
\end{algorithm}

\subsubsection{Características Clave}

\begin{itemize}[leftmargin=*]
    \item \textbf{Complejidad}: $\mathcal{O}(n \cdot d)$ durante el entrenamiento (lineal en la longitud de la secuencia); $\mathcal{O}(k)$ por token durante la generación.
    \item \textbf{Entrenable}: Sí; flujo de gradiente de extremo a extremo a través del operador SparseK.
    \item \textbf{Generación incremental}: Soporta generación autorregresiva eficiente con memoria constante.
    \item \textbf{Fortaleza}: Se integra perfectamente en arquitecturas LLM existentes; mínimo ajuste fino necesario.
    \item \textbf{Limitación}: El acceso a memoria disperso de la selección top-$k$ puede limitar la eficiencia de caché en algunos hardwares; la red de puntuación añade sobrecarga.
\end{itemize}

\subsection{NSA (Native Sparse Attention): Ramas Jerárquicas Alineadas con el Hardware}

\subsubsection{Formulación Matemática}

NSA \citep{yuan_nsa_2025} descompone la atención dispersa en tres ramas paralelas que se combinan mediante compuertas aprendidas.

\paragraph{Rama de compresión}: Un MLP aprendible $\varphi$ comprime bloques de clave-valor:

\[
\tilde{K}_t^{\text{cmp}} = \left[\varphi(k_{id+1:id+l})\right]_{\substack{1 \leq i \leq \lfloor(t-l)/d\rfloor}} \quad, \quad \tilde{V}_t^{\text{cmp}} = \left[\varphi(v_{id+1:id+l})\right]
\]

\paragraph{Rama de selección}: Los bloques de alta importancia se seleccionan basándose en puntajes comprimidos:

\[
p_t^{\text{cmp}} = \text{softmax}(q_t^\top \tilde{K}_t^{\text{cmp}} / \sqrt{d}), \quad I_t = \mathrm{TopK}(p_t^{\text{cmp}}, n), \quad \tilde{K}_t^{\text{sel}} = \mathrm{Recolectar}(K, I_t)
\]

\paragraph{Rama de ventana deslizante}: Contexto local fijo:

\[
\tilde{K}_t^{\text{win}} = K_{t-w:t}, \quad \tilde{V}_t^{\text{win}} = V_{t-w:t}
\]

\paragraph{Combinación con compuertas}: Las tres salidas de atención se combinan mediante compuertas aprendidas:

\[
\alpha_t^{(\text{cmp})} = \sigma(\text{Lineal}_\text{cmp}(q_t)), \quad \alpha_t^{(\text{sel})} = \sigma(\text{Lineal}_\text{sel}(q_t)), \quad \alpha_t^{(\text{win})} = \sigma(\text{Lineal}_\text{win}(q_t))
\]

\[
y_t = \alpha_t^{(\text{cmp})} \cdot \text{Attn}(q_t, \tilde{K}^{\text{cmp}}, \tilde{V}^{\text{cmp}}) + \alpha_t^{(\text{sel})} \cdot \text{Attn}(q_t, \tilde{K}^{\text{sel}}, \tilde{V}^{\text{sel}}) + \alpha_t^{(\text{win})} \cdot \text{Attn}(q_t, \tilde{K}^{\text{win}}, \tilde{V}^{\text{win}})
\]

\subsubsection{Pseudocódigo Algorítmico}

\begin{algorithm}[H]
\caption{NSA: Ramas Comprimida + Seleccionada + Ventana con Compuertas Aprendidas}
\begin{algorithmic}[1]
\Require Consulta $\mathbf{q}_t$, secuencias de Clave/Valor almacenadas en caché, tamaño de bloque $l$, stride $d$, conteo de selección $n$, ventana $w$
\Require MLP de compresión $\varphi$, redes de compuerta $\text{Lineal}_{\text{cmp}}, \text{Lineal}_{\text{sel}}, \text{Lineal}_{\text{win}}$
\Ensure Salida $\mathbf{y}_t$
\State \textbf{Rama de compresión:}
\For{$i = 1$ to $\lfloor t/d \rfloor$}
    \State $k_i^{\text{cmp}} \gets \varphi(\mathbf{K}_{id+1:id+l})$ \Comment{Comprimir bloque mediante MLP}
    \State $v_i^{\text{cmp}} \gets \varphi(\mathbf{V}_{id+1:id+l})$
\EndFor
\State $\tilde{\mathbf{K}}^{\text{cmp}} \gets [k_1^{\text{cmp}}, \ldots, k_{\lfloor t/d \rfloor}^{\text{cmp}}]$, $\tilde{\mathbf{V}}^{\text{cmp}} \gets [v_1^{\text{cmp}}, \ldots, v_{\lfloor t/d \rfloor}^{\text{cmp}}]$
\State $\mathbf{y}_t^{(\text{cmp})} \gets \text{SDPA}(\mathbf{q}_t, \tilde{\mathbf{K}}^{\text{cmp}}, \tilde{\mathbf{V}}^{\text{cmp}})$
\State \textbf{Rama de selección:}
\State Calcular puntajes sobre comprimidos: $\mathbf{s} \gets \text{softmax}(\mathbf{q}_t \tilde{\mathbf{K}}^{\text{cmp}\top} / \sqrt{d})$
\State Seleccionar los $n$ bloques más importantes: $I \gets \mathrm{TopK}(\mathbf{s}, n)$
\State Recolectar bloques completos: $\tilde{\mathbf{K}}^{\text{sel}} \gets [\mathbf{K}_{I_i d+1:I_i d+l}]_i$, $\tilde{\mathbf{V}}^{\text{sel}} \gets [\mathbf{V}_{I_i d+1:I_i d+l}]_i$
\State $\mathbf{y}_t^{(\text{sel})} \gets \text{SDPA}(\mathbf{q}_t, \tilde{\mathbf{K}}^{\text{sel}}, \tilde{\mathbf{V}}^{\text{sel}})$
\State \textbf{Rama de ventana:}
\State $\mathbf{y}_t^{(\text{win})} \gets \text{SDPA}(\mathbf{q}_t, \mathbf{K}_{t-w:t}, \mathbf{V}_{t-w:t})$
\State \textbf{Fusión con compuertas:}
\State $\alpha^{(\text{cmp})} \gets \sigma(\text{Lineal}_{\text{cmp}}(\mathbf{q}_t))$, $\alpha^{(\text{sel})} \gets \sigma(\text{Lineal}_{\text{sel}}(\mathbf{q}_t))$, $\alpha^{(\text{win})} \gets \sigma(\text{Lineal}_{\text{win}}(\mathbf{q}_t))$
\State $\mathbf{y}_t \gets \alpha^{(\text{cmp})} \mathbf{y}_t^{(\text{cmp})} + \alpha^{(\text{sel})} \mathbf{y}_t^{(\text{sel})} + \alpha^{(\text{win})} \mathbf{y}_t^{(\text{win})}$
\Return $\mathbf{y}_t$
\end{algorithmic}
\end{algorithm}

\subsubsection{Características Clave}

\begin{itemize}[leftmargin=*]
    \item \textbf{Complejidad}: $\mathcal{O}(t/d + n \cdot l + w)$ tokens atendidos por paso (significativamente reducido frente a $\mathcal{O}(t)$ para densa).
    \item \textbf{Entrenable}: Sí; entrenamiento directo con todas las ramas diferenciables.
    \item \textbf{Alineación con hardware}: Diseñado para utilización eficiente de núcleos tensoriales; la compresión y selección reducen la presión sobre el ancho de banda de memoria.
    \item \textbf{Fortaleza}: 11.6$\times$ de aceleración en decodificación y 9$\times$ de aceleración hacia adelante en secuencias de 64k; supera a la atención completa en muchas tareas.
    \item \textbf{Limitación}: Requiere kernels Triton/CUDA personalizados; la arquitectura multi-rama compleja aumenta la sobrecarga de ingeniería.
\end{itemize}

\subsection{FASA: Atención Dispersa Consciente de la Frecuencia}

\subsubsection{Formulación Matemática}

FASA \citep{wang_fasa_2026} es un \textbf{método en tiempo de inferencia sin entrenamiento} que explota la estructura del embedding de posición rotatorio (RoPE). Bajo RoPE, el embedding del token se rota por $\theta_i = B^{-2(i-1)/d}$, descomponiendo el espacio $d$-dimensional en $d/2$ fragmentos de frecuencia (FC).

\paragraph{Idea clave}: Solo un subconjunto pequeño ($ < 1\%$) de los FC importa para la conciencia contextual; la mayoría codifica patrones posicionales.

\paragraph{Identificación de fragmento de frecuencia dominante}: Para cada capa $l$ y cabeza $h$, identificar FC dominantes mediante concordancia contextual (CA):

\[
\text{CA}^{l,h,i} = \frac{|\text{TopK}(\alpha^{l,h}) \cap \text{TopK}(\alpha^{l,h,i})|}{K}
\]

donde $\alpha^{l,h}$ es la máscara de atención completa y $\alpha^{l,h,i}$ es la máscara usando solo el FC $i$. Los FC dominantes tienen CA alta---contribuyen significativamente al patrón de atención final.

\paragraph{Etapa de predicción de importancia de tokens (TIP)}: Usando solo FC dominantes, calcular importancia ligera por token:

\[
S_t^{l,h} = \sum_{i \in \mathcal{I}_{\text{dom}}^{l,h}} \alpha^{l,h,i}(\mathbf{q}_t, \mathbf{K}_{1:t}), \quad \mathcal{T}_t = \text{TopK-Índices}(S_t, N_{\text{fac}})
\]

\paragraph{Etapa de cálculo de atención enfocada (FAC)}: Calcular atención de precisión completa en los tokens seleccionados:

\[
\hat{\alpha}_{\text{FAC}} = \text{softmax}\left(\frac{\mathbf{q}_t \mathbf{K}_{\mathcal{T}_t}^\top}{\sqrt{d}}\right), \quad \mathbf{o}_t = \hat{\alpha}_{\text{FAC}} \mathbf{V}_{\mathcal{T}_t}
\]

\subsubsection{Pseudocódigo Algorítmico}

\begin{algorithm}[H]
\caption{FASA: Atención Dispersa Consciente de la Frecuencia (en Tiempo de Inferencia)}
\begin{algorithmic}[1]
\Require Consulta actual $\mathbf{q}_t$, Clave/Valor almacenados en caché $\mathbf{K}_{1:t}, \mathbf{V}_{1:t}$, base RoPE $B$, FC dominantes $\mathcal{I}_{\text{dom}}$
\Require Presupuesto TIP $N_{\text{tip}}$, presupuesto FAC $N_{\text{fac}}$
\Ensure Salida $\mathbf{o}_t$
\State \textbf{Etapa 1: Predicción de Importancia de Tokens (TIP)}
\For{capa $l$ y cabeza $h$}
    \For{posición $i = 1$ to $t$}
        \State Calcular importancia desde FC dominantes: $s_i \gets 0$
        \For{$fc \in \mathcal{I}_{\text{dom}}^{l,h}$}
            \State Rotar consulta y clave en FC $fc$: $\mathbf{q}^{(fc)}, \mathbf{k}_i^{(fc)}$
            \State $s_i^{(fc)} \gets \text{softmax}(\mathbf{q}^{(fc)} \mathbf{k}_i^{(fc)\top} / \sqrt{d}) $
            \State $s_i \gets s_i + s_i^{(fc)}$
        \EndFor
    \EndFor
    \State $\mathcal{T}_t \gets \mathrm{TopK-Índices}([s_1, \ldots, s_t], N_{\text{fac}})$ \Comment{Seleccionar tokens principales}
\EndFor
\State \textbf{Etapa 2: Cálculo de Atención Enfocada (FAC)}
\State $\text{puntajes}_{\text{fac}} \gets \mathbf{q}_t \mathbf{K}_{\mathcal{T}_t}^\top / \sqrt{d}$ \Comment{Producto punto de precisión completa}
\State $\alpha_{\text{fac}} \gets \text{softmax}(\text{puntajes}_{\text{fac}})$ \Comment{Softmax completo}
\State $\mathbf{o}_t \gets \alpha_{\text{fac}} \mathbf{V}_{\mathcal{T}_t}$ \Comment{Suma ponderada de valores}
\Return $\mathbf{o}_t$
\end{algorithmic}
\end{algorithm}

\subsubsection{Características Clave}

\begin{itemize}[leftmargin=*]
    \item \textbf{Complejidad}: TIP es $\mathcal{O}(t \cdot N_{\text{tip}})$; FAC es $\mathcal{O}(N_{\text{fac}} \cdot d)$.
    \item \textbf{Entrenable}: No; optimización de inferencia sin entrenamiento.
    \item \textbf{Aplicabilidad}: Requiere modelos basados en RoPE; extendido a ALiBi y MLA con modificaciones.
    \item \textbf{Fortaleza}: Precisión cercana a la óptima con $\leq 256$ tokens de entre millones; hasta 2.56$\times$ de aceleración en decodificación; ortogonal a la cuantización.
    \item \textbf{Limitación}: Requiere calibración fuera de línea; la identificación de FC dominantes añade sobrecarga de preprocesamiento.
\end{itemize}

\subsection{SpargeAttn: Filtrado de Dos Etapas a Nivel de Bloques}

\subsubsection{Formulación Matemática}

SpargeAttn \citep{zhang_spargeattn_2025} es un \textbf{método universal sin entrenamiento} que predice qué bloques de la matriz de atención contendrán valores insignificantes y omite el cálculo para esos bloques.

\paragraph{Etapa 1 --- Predicción dispersa}: Calcular importancia proxy de bajo costo $\hat{s}_{ij}$ para el bloque $(i,j)$:

\[
\hat{s}_{ij} = f_{\text{pred}}(\mathbf{Q}_i, \mathbf{K}_j; \text{hiperparámetros}), \quad \text{omitir si } \hat{s}_{ij} < \epsilon_1
\]

Los predictores comunes incluyen similitud de media de bloque o estadísticas de autosimilitud.

\paragraph{Etapa 2 --- Filtrado consciente de softmax}: Después de calcular $\tilde{P}_{ij} = \text{softmax}(Q_i K_j^\top/\sqrt{d})$, verificar si la probabilidad máxima del bloque es insignificante en relación con el máximo de softmax en ejecución:

\[
\text{omitir } P_{ij} V_j \quad \text{si} \quad \max(\tilde{P}_{ij}) < e^{m_{\text{antiguo}} - m_{\text{nuevo}}} \cdot \epsilon_2
\]

donde $m_{\text{antiguo}}, m_{\text{nuevo}}$ son máximos del softmax en línea en ejecución (calculados mediante numéricos estilo FlashAttention).

\subsubsection{Pseudocódigo Algorítmico}

\begin{algorithm}[H]
\caption{SpargeAttn: Filtrado Disperso de Dos Etapas por Bloques}
\begin{algorithmic}[1]
\Require Bloques de consulta $\mathbf{Q}_i$ para $i = 1, \ldots, n_b$, Bloques de clave $\mathbf{K}_j$ para $j = 1, \ldots, n_b$, Bloques de valor $\mathbf{V}_j$
\Require Umbral de predicción $\epsilon_1$, umbral de softmax $\epsilon_2$, tamaño de bloque $b_s$
\Ensure Salida $\mathbf{Y}$
\State Inicializar: máximo de softmax en línea $m_{\text{global}} \gets -\infty$
\For{fila de bloque $i = 1$ to $n_b$}
    \State Inicializar: salida de fila de bloque $\mathbf{Y}_i \gets 0$, máximo local $m_i \gets -\infty$
    \For{columna de bloque $j = 1$ to $n_b$}
        \State \textbf{Etapa 1: Predicción Dispersa}
        \State Calcular importancia proxy: $\hat{s}_{ij} \gets f_{\text{pred}}(\mathbf{Q}_i, \mathbf{K}_j)$
        \If{$\hat{s}_{ij} < \epsilon_1$}
            \State \textbf{omitir} este bloque $[i,j]$  \Comment{Terminación temprana}
            \State \textbf{continuar} \Comment{Pasar al siguiente bloque}
        \EndIf
        \State \textbf{Etapa 2: Calcular y Filtrar}
        \State $\tilde{P}_{ij} \gets \text{softmax}(\mathbf{Q}_i \mathbf{K}_j^\top / \sqrt{d})$
        \State $m_{\text{bloque}} \gets \max(\tilde{P}_{ij})$
        \If{$m_{\text{bloque}} < e^{m_i - m_{\text{global}}} \cdot \epsilon_2$}
            \State El bloque contribuye insignificantemente; omitir  \Comment{Poda consciente de softmax}
            \State \textbf{continuar}
        \EndIf
        \State Actualizar máximo global: $m_i \gets \max(m_i, m_{\text{bloque}})$
        \State Acumular: $\mathbf{Y}_i \gets \mathbf{Y}_i + \tilde{P}_{ij} \mathbf{V}_j$
    \EndFor
    \State $m_{\text{global}} \gets \max(m_{\text{global}}, m_i)$
\EndFor
\Return $\mathbf{Y}$
\end{algorithmic}
\end{algorithm}

\subsubsection{Características Clave}

\begin{itemize}[leftmargin=*]
    \item \textbf{Complejidad}: Empírica $\mathcal{O}(n^2 \cdot s)$ donde $s$ es la fracción de bloques no omitidos (típicamente 0.2--0.5).
    \item \textbf{Entrenable}: No; aceleración plug-and-play para modelos existentes.
    \item \textbf{Universalidad}: Funciona en modelos de lenguaje, modelos de difusión de imágenes y generación de video.
    \item \textbf{Fortaleza}: 2.5--5$\times$ de aceleración comparado con métodos densos o dispersos anteriores; compatible con cuantización (integración SageAttention).
    \item \textbf{Limitación}: La aceleración depende de la dispersidad inherente; la granularidad a nivel de bloque puede pasar por alto patrones más finos; se requiere ajuste de umbral por familia de modelos.
\end{itemize}

\subsection{Resumen Comparativo}

\footnotesize
\begin{longtable}{p{1.9cm}p{1.7cm}p{1.6cm}p{1.6cm}p{1.1cm}p{3.5cm}}
\toprule
\textbf{Método} & \textbf{Complejidad} & \textbf{Entrenable} & \textbf{Unidad} & \textbf{Año} & \textbf{Contribución Principal} \\
\midrule
\endhead
Sparse Transformer & $O(n\sqrt{n}d)$ & Sí & Patrón de token & 2019 & Máscaras fijas y estrideadas factorizadas \\
Longformer & $O(nwd)$ & Sí & Local + global & 2020 & Escalado lineal con dilatación \\
BigBird & $O(n)$ & Sí & Grafo disperso & 2020 & Prueba de completitud teórica \\
SparseK & $O(nd)$ & Sí & Tokens top-$k$ & 2024 & Selección diferenciable \\
NSA & $O((t/d+nl'+w)d)$ & Sí & Multi-rama & 2025 & Diseño de 3 ramas alineado con hardware \\
FASA & $O(tN_{\text{tip}}+N_{\text{fac}}d)$ & No & Fragmentos de frecuencia & 2026 & Conocimiento de frecuencia RoPE \\
SpargeAttn & $O(n^2 s \cdot d)$ & No & Filtro de bloques & 2025 & Filtrado universal de dos etapas \\
\bottomrule
\end{longtable}
\normalsize

\subsection{Espacio de Diseño y Criterios de Selección}

Los siete métodos de atención dispersa ocupan posiciones complementarias en un espacio de diseño multidimensional:

\begin{itemize}[leftmargin=*]
    \item \textbf{Patrones geométricos/fijos} (Sparse Transformer, Longformer): Simples de implementar y analizar, pero los patrones no se adaptan al contenido.
    \item \textbf{Selección aprendida} (SparseK, NSA): Permiten adaptación mediante redes de selección entrenables, a costa de mayor complejidad de implementación.
    \item \textbf{Aceleración sin entrenamiento} (FASA, SpargeAttn): Permiten aceleración directa de modelos existentes sin reentrenamiento, adecuados para sistemas desplegados.
    \item \textbf{Garantías teóricas} (BigBird): Proporcionan demostraciones formales de completitud expresiva, importantes para comprender los márgenes de seguridad.
\end{itemize}

\textbf{Orientación de selección:}
\begin{enumerate}[leftmargin=*]
    \item \textbf{Entrenamiento de nuevos modelos} donde los recursos lo permitan: NSA o SparseK para máxima aceleración de inferencia; BigBird para garantía teórica.
    \item \textbf{Comprensión de documentos de contexto largo}: Longformer o FASA por simplicidad práctica; NSA para escala extrema.
    \item \textbf{Aceleración de modelos existentes}: FASA para LLMs basados en RoPE; SpargeAttn para cualquier arquitectura.
    \item \textbf{Entornos con recursos limitados}: Sparse Transformer por simplicidad y eficiencia probada a escala moderada.
\end{enumerate}

\section{Annex D: Gated Attention Families---Complete Literature Analysis}
\label{sec:annex-gated}

\subsection{Resumen Ejecutivo: Compuertas para el Control de Memoria}

El campo emergente de los \emph{mecanismos de atención con compuerta} aborda un desafío fundamental en el modelado de secuencias: ¿cómo debe retenerse, olvidarse y actualizarse la información a medida que el modelo procesa nuevos tokens? Los estados recurrentes aditivos tradicionales acumulan toda la información sin borrado, lo que lleva a la saturación de memoria y a la incapacidad de adaptarse a contextos cambiantes. La atención softmax proporciona expresividad completa pero usa un caché KV de $\mathcal{O}(Ld)$ durante la inferencia, limitando el tamaño de la ventana de contexto. Los mecanismos con compuerta proporcionan un punto intermedio: control selectivo sobre la supervivencia de la información.

El principio unificador entre las arquitecturas con compuerta es que las compuertas controlan \textit{qué información sobrevive}. Las compuertas operan en tres niveles distintos:
\begin{enumerate}[leftmargin=*]
    \item \textbf{Decaimiento del estado recurrente} (GLA, HGRN2): Compuertas multiplicativas sobre la matriz de memoria $S_t$ controlan cuánto del estado anterior persiste en el siguiente paso.
    \item \textbf{Intensidad de escritura en actualizaciones recurrentes} (DeltaNet, Gated DeltaNet, Gated DeltaNet-2): Las compuertas controlan la magnitud y la dirección canalizada de la nueva información escrita en la memoria.
    \item \textbf{Sesgo de logits softmax} (FoX): Las compuertas inyectan sesgo de actualidad a nivel de token en el cálculo de logits de atención.
    \item \textbf{Dispersidad post-atención} (Gated Softmax): Las compuertas escalan selectivamente los canales de salida de SDPA.
\end{enumerate}

Este apéndice sintetiza ocho arquitecturas principales de atención con compuerta publicadas entre 2023 y 2026, analizando sus fundamentos matemáticos, características de eficiencia en hardware y compensaciones empíricas.

\subsection{1. Atención Lineal con Compuerta (GLA)}

\subsubsection{Fundamento Matemático}

La Atención Lineal con Compuerta \citep{yang_gla_2023} aumenta la atención lineal estándar (que usa un estado recurrente con valor de matriz) con decaimiento multiplicativo dependiente de los datos. La atención lineal estándar reformula el mecanismo softmax tradicional como una recurrencia sobre estados ocultos:
\[
\mathbf{S}_t = \mathbf{S}_{t-1} + \mathbf{v}_t \mathbf{k}_t^\top, \quad \mathbf{o}_t = \mathbf{S}_t \mathbf{q}_t
\]
donde el estado $\mathbf{S}_t \in \mathbb{R}^{d_v \times d_k}$ acumula productos externos. Esta formulación puramente aditiva sufre de ``sobrecarga de memoria'': la información se acumula monótonamente y no puede borrarse, degradando el rendimiento en tareas que requieren selección de contexto.

GLA introduce una \textbf{matriz de compuerta diagonal} $\mathbf{G}_t \in [0,1]^{d_k \times d_k}$ dependiente de los datos:
\[
\mathbf{S}_t = \mathbf{G}_t \odot \mathbf{S}_{t-1} + \mathbf{v}_t \mathbf{k}_t^\top, \quad \mathbf{o}_t = \mathbf{S}_t \mathbf{q}_t
\]
La compuerta $\mathbf{G}_t$ se parametriza con una estructura de producto externo:
\[
\mathbf{G}_t = \mathbf{1} \boldsymbol{\alpha}_t^\top, \quad \boldsymbol{\alpha}_t = \text{logsigmoid}\left(\mathbf{W}_{gk} \mathbf{x}_t\right) / c
\]
donde $\mathbf{W}_{gk}$ es una proyección de rango bajo ($\mathbb{R}^{d \to d_k}$) y $c \approx 16$ es un normalizador. La activación logsigmoid mapea $\alpha_t \in \mathbb{R}$ a aproximadamente $[-1/2, 0]$, y la división por $c$ contrae esto aún más para que esté cerca de la identidad en la inicialización. La compuerta resultante $G_t = 1 + \alpha_t \approx 1$ cerca de la inicialización, luego aprende a decaer $(\alpha_t \to -\infty)$ información histórica.

\begin{algorithm}[H]
\caption{Atención Lineal con Compuerta (Forma Recurrente)}
\begin{algorithmic}[1]
\Require Entrada $\mathbf{x}_t$; estado anterior $\mathbf{S}_{t-1}$
\Require Proyecciones: $W_q, W_k, W_v$ (estándar); $W_{gk}$ (proyección de compuerta)
\Ensure Salida $\mathbf{o}_t$ y estado actualizado $\mathbf{S}_t$
\State $\mathbf{q}_t \gets W_q \mathbf{x}_t$
\State $\mathbf{k}_t \gets W_k \mathbf{x}_t$
\State $\mathbf{v}_t \gets W_v \mathbf{x}_t$
\State Calcular compuerta: $\boldsymbol{\alpha}_t \gets \text{logsigmoid}(W_{gk} \mathbf{x}_t) / 16$ \Comment{Decaimiento por dimensión de clave}
\State Aplicar compuerta de producto externo: $\mathbf{G}_t \gets \mathbf{1} \boldsymbol{\alpha}_t^\top$  \Comment{producto externo diagonal}
\State Decaer estado anterior: $\mathbf{S}_t^\text{decaimiento} \gets \mathbf{G}_t \odot \mathbf{S}_{t-1}$ \Comment{producto elemento a elemento}
\State Agregar nueva asociación: $\mathbf{S}_t \gets \mathbf{S}_t^\text{decaimiento} + \mathbf{v}_t \mathbf{k}_t^\top$
\State Recuperar mediante consulta: $\mathbf{o}_t^\text{crudo} \gets \mathbf{S}_t \mathbf{q}_t$
\State Aplicar compuerta de salida: $\mathbf{o}_t \gets \text{GroupNorm}(\mathbf{o}_t^\text{crudo}) \otimes \text{SiLU}(W_g \mathbf{x}_t)$
\Return $\mathbf{o}_t$, $\mathbf{S}_t$
\end{algorithmic}
\end{algorithm}

\subsubsection{Entrenamiento Eficiente en Hardware}

Para el entrenamiento paralelo, GLA usa un algoritmo por fragmentos que agrupa tokens en fragmentos $C$ y calcula transiciones recurrentes en paralelo:
\[
\mathbf{O}_{[t]} = \overleftarrow{\mathbf{Q}}_{[t]} \mathbf{S}_{[t]}^\top + \left(\mathbf{Q}_{[t]} \mathbf{K}_{[t]}^\top \odot \mathbf{\Gamma}_{[t]}\right) \mathbf{V}_{[t]}
\]
donde $\overleftarrow{\mathbf{Q}}_{[t]}$ atiende al estado en el límite del fragmento (retrospectiva), y $\mathbf{\Gamma}_{[t]}$ es la máscara causal con escalado consciente del decaimiento:
\[
\mathbf{\Gamma}_{[t],ij} = \begin{cases} \prod_{l=j}^{i-1} \alpha_l & \text{si } i > j \text{ (retrospectiva)} \\ \prod_{l=1}^{i-j} \alpha_l & \text{si } i \le j \text{ (en-fragmento)} \end{cases}
\]
Este diseño maximiza las operaciones matmul susceptibles de aceleración mediante núcleos tensoriales, logrando una complejidad total subcuadrática $\mathcal{O}(nLd^2/C)$ donde $L$ es el número de cabezas de atención y $C$ es el tamaño del fragmento.

\subsubsection{Fortalezas y Limitaciones}

\begin{table}[H]
\centering
\small
\begin{tabularx}{0.95\linewidth}{Xp{4cm}p{4cm}}
\toprule
\textbf{Aspecto} & \textbf{Fortalezas} & \textbf{Limitaciones} \\
\midrule
Eficiencia computacional & Entrenamiento subcuadrático mediante paralelismo por fragmentos. Inferencia con memoria constante $O(d^2)$. & Requiere kernels CUDA/Triton personalizados; no disponible en todos los frameworks. \\
Generalización de contexto & Extiende entrenamiento de 2K tokens a más de 20K tokens con degradación de perplejidad insignificante. & Aún rinde por debajo de softmax en benchmarks con mucha recuperación (aguja en el pajar). \\
Selectividad & Decaimiento dependiente de los datos por dimensión de clave. & La estructura de compuerta limita la selectividad por par clave-valor; no hay control directo sobre qué asociaciones específicas borrar. \\
\bottomrule
\end{tabularx}
\end{table}

\subsection{2. DeltaNet: Atención Lineal con Corrección de Errores}

\subsubsection{Fundamento Matemático}

DeltaNet \citep{yang_deltanet_2024} aplica una \textbf{regla de aprendizaje delta} clásica a las actualizaciones de estado de la atención lineal. En lugar de acumulación aditiva, DeltaNet calcula la diferencia entre el valor predicho (recuperado de la memoria) y el valor objetivo, luego realiza una actualización de corrección de errores:
\[
\mathbf{S}_t = \mathbf{S}_{t-1}(\mathbf{I} - \beta_t \mathbf{k}_t \mathbf{k}_t^\top) + \beta_t \mathbf{v}_t \mathbf{k}_t^\top
\]
donde $\beta_t \in (0,1)$ es un escalar de \textbf{intensidad de escritura} dependiente de los datos. Esta actualización puede descomponerse como una operación de borrado-escritura:
\[
\mathbf{v}_t^\text{anterior} = \mathbf{S}_{t-1} \mathbf{k}_t \quad \text{(predicción actual)}, \quad
\mathbf{v}_t^\text{actualizado} = \beta_t \mathbf{v}_t + (1-\beta_t)\mathbf{v}_t^\text{anterior} \quad \text{(mezcla anterior/nuevo)}
\]
de modo que:
\[
\mathbf{S}_t = \mathbf{S}_{t-1} - \mathbf{v}_t^\text{anterior}\mathbf{k}_t^\top + \mathbf{v}_t^\text{actualizado}\mathbf{k}_t^\top
\]

La idea clave es que esta regla de actualización minimiza una pérdida de ECM en línea en cada paso temporal:
\[
\mathcal{L}_t(\mathbf{S}) = \tfrac{1}{2}\|\mathbf{S}\mathbf{k}_t - \mathbf{v}_t\|^2 \Rightarrow \nabla_{\mathbf{S}} \mathcal{L}_t = (\mathbf{S}\mathbf{k}_t - \mathbf{v}_t)\mathbf{k}_t^\top
\]

Un paso de SGD con tasa de aprendizaje $\beta_t$ produce la actualización de la regla delta. Esta conexión con el entrenamiento en tiempo de prueba (TTT) proporciona fundamentación teórica: DeltaNet optimiza directamente la predicción de valores en tiempo de inferencia.

\subsubsection{Entrenamiento Paralelo Eficiente}

La matriz de transición de DeltaNet $\mathbf{I} - \beta_t \mathbf{k}_t \mathbf{k}_t^\top$ es una matriz de Householder generalizada (actualización de rango 1 a la identidad). Para el entrenamiento por fragmentos, DeltaNet usa la \textbf{representación WY}, que representa el producto de $m$ de estas actualizaciones de rango 1 de forma compacta:
\[
\prod_{i=1}^{m} (\mathbf{I} - \beta_i \mathbf{k}_i \mathbf{k}_i^\top) = \mathbf{I} - \mathbf{W}\mathbf{Y}^\top
\]
donde $\mathbf{W}, \mathbf{Y} \in \mathbb{R}^{d \times m}$ se construyen recursivamente. Esta factorización permite paralelismo rico en matmul sin materializar el producto completo, manteniendo el entrenamiento por fragmentos eficiente.

\begin{algorithm}[H]
\caption{Actualización de Estado de DeltaNet con Representación WY}
\begin{algorithmic}[1]
\Require Fragmento de entrada $\mathbf{X}_{[t]}$; estado anterior $\mathbf{S}_{t-1}$
\Require Consultas normalizadas L2 $\hat{\mathbf{Q}}_t$, claves $\hat{\mathbf{K}}_t$, valores $\mathbf{V}_t$
\Ensure Salida $\mathbf{O}_{[t]}$ y estado actualizado $\mathbf{S}_t$
\State \textbf{Fase por fragmentos:}
\For{posición $i = 1$ to tamaño\_fragmento}
    \State Normalizar: $\hat{\mathbf{k}}_i \gets \hat{\mathbf{K}}_t[i] / \|\hat{\mathbf{K}}_t[i]\|$, $\hat{\mathbf{q}}_i \gets \hat{\mathbf{Q}}_t[i] / \|\hat{\mathbf{Q}}_t[i]\|$
    \State Calcular intensidad de escritura: $\beta_i \gets \text{sigmoid}(W_\beta \mathbf{x}_i)$
    \State Predicción: $\hat{\mathbf{v}}^\text{pred} \gets \mathbf{S}_{i-1}^\text{en-fragmento} \hat{\mathbf{k}}_i$
    \State Mezcla consciente del error: $\mathbf{v}_i^\text{actualización} \gets \beta_i(\mathbf{v}_i - \hat{\mathbf{v}}^\text{pred})$
    \State Borrar-escribir: $\mathbf{S}_i^\text{en-fragmento} \gets \mathbf{S}_{i-1}^\text{en-fragmento} - \hat{\mathbf{v}}^\text{pred}\hat{\mathbf{k}}_i^\top + (\beta_i \mathbf{v}_i + (1-\beta_i)\hat{\mathbf{v}}^\text{pred})\hat{\mathbf{k}}_i^\top$
    \State Salida: $\mathbf{o}_i \gets \mathbf{S}_i^\text{en-fragmento} \hat{\mathbf{q}}_i$
\EndFor
\State \textbf{Fase entre fragmentos:} Usar representación WY de los productos de transición
\Return $\mathbf{O}_{[t]}$, $\mathbf{S}_{t}$
\end{algorithmic}
\end{algorithm}

\subsubsection{Fortalezas y Limitaciones}

\begin{table}[H]
\centering
\small
\begin{tabularx}{0.95\linewidth}{Xp{4cm}p{4cm}}
\toprule
\textbf{Aspecto} & \textbf{Fortalezas} & \textbf{Limitaciones} \\
\midrule
Recuperación asociativa & Recuperación perfecta en el benchmark MQAR (Recuperación Asociativa Multi-Consulta); la semántica de corrección de errores se ajusta naturalmente a patrones de recuperación. & Sin olvido global, la memoria aún se satura en longitudes de secuencia extremas; requiere mecanismos de reinicio periódico para contexto ilimitado. \\
Teoría & Fundamentada en optimización de ECM en línea; conexión clara con el paradigma de entrenamiento en tiempo de prueba. & Requiere claves normalizadas L2 para estabilidad numérica; la normalización de doble ancho añade sobrecarga. \\
Rendimiento & La representación WY permite entrenamiento eficiente por fragmentos. & Ligeramente más lento que Mamba2 por token debido a matrices de transición más ricas (rango 1 en lugar de diagonales). \\
\bottomrule
\end{tabularx}
\end{table}

\subsection{3. Gated DeltaNet: Síntesis de Compuerta y Corrección de Errores}

\subsubsection{Fundamento Matemático}

Gated DeltaNet \citep{yang_gated_deltanet_2024} sintetiza las fortalezas de GLA y DeltaNet combinando una \textbf{compuerta de decaimiento} $\alpha_t$ (de GLA) con la \textbf{compuerta de intensidad de escritura} $\beta_t$ (de DeltaNet):
\[
\mathbf{S}_t = \mathbf{S}_{t-1}\left(\alpha_t (\mathbf{I} - \beta_t \mathbf{k}_t \mathbf{k}_t^\top)\right) + \beta_t \mathbf{v}_t \mathbf{k}_t^\top
\]
Reescribiendo para mayor claridad:
\[
\mathbf{S}_t = \alpha_t \mathbf{S}_{t-1}(\mathbf{I} - \beta_t \mathbf{k}_t \mathbf{k}_t^\top) + \beta_t \mathbf{v}_t \mathbf{k}_t^\top
\]
Las dos compuertas son complementarias:
\begin{itemize}[leftmargin=*]
    \item $\alpha_t \to 0$ borra rápidamente el estado histórico: $\mathbf{S}_t \approx \beta_t \mathbf{v}_t \mathbf{k}_t^\top$ (reinicio de contexto).
    \item $\alpha_t \to 1$ recupera el comportamiento puro de regla delta: $\mathbf{S}_t \approx \mathbf{S}_{t-1}(\mathbf{I} - \beta_t \mathbf{k}_t \mathbf{k}_t^\top) + \beta_t \mathbf{v}_t \mathbf{k}_t^\top$ (actualizaciones dirigidas).
    \item $\beta_t \to 0$ omite la escritura pero decae: $\mathbf{S}_t \approx \alpha_t \mathbf{S}_{t-1}$ (olvido puro).
    \item $\beta_t \to 1$ reemplaza la memoria agresivamente: $\mathbf{S}_t \approx \alpha_t \mathbf{S}_{t-1} + \mathbf{v}_t \mathbf{k}_t^\top$ (reemplazo).
\end{itemize}

Desde una perspectiva de aprendizaje en línea, Gated DeltaNet corresponde a:
\[
\min_{\mathbf{S}_t} \|\mathbf{S}_t - \alpha_t \mathbf{S}_{t-1}\|_F^2 - 2\langle \mathbf{S}_t \mathbf{k}_t, \beta_t(\mathbf{v}_t - \alpha_t \mathbf{S}_{t-1} \mathbf{k}_t)\rangle
\]
que introduce un decaimiento de pesos adaptativo $\alpha_t$ en una actualización tipo SGD---análogo al decaimiento de pesos desacoplado en la optimización de aprendizaje profundo.

\begin{algorithm}[H]
\caption{Paso Adelante Paralelo por Fragmentos de Gated DeltaNet}
\begin{algorithmic}[1]
\Require Fragmento de entrada $\mathbf{X}_{[i]}$; estado anterior $\mathbf{S}_{i-1}$; tamaño de fragmento $C$
\Require Proyecciones: $W_q, W_k, W_v, W_\alpha, W_\beta$
\Ensure Salida $\mathbf{O}_{[i]}$ y estado $\mathbf{S}_i$
\State \textbf{Cálculo en-fragmento:}
\For{$j = 1$ to $C$}
    \State $\mathbf{q}_j \gets W_q \mathbf{x}_j$, $\mathbf{k}_j \gets W_k \mathbf{x}_j$, $\mathbf{v}_j \gets W_v \mathbf{x}_j$
    \State $\alpha_j \gets \text{sigmoid}(W_\alpha \mathbf{x}_j)$, $\beta_j \gets \text{sigmoid}(W_\beta \mathbf{x}_j)$
    \State Aplicar compuerta combinada: $\mathbf{S}_j \gets \alpha_j \mathbf{S}_{j-1}(\mathbf{I} - \beta_j \mathbf{k}_j \mathbf{k}_j^\top) + \beta_j \mathbf{v}_j \mathbf{k}_j^\top$
    \State $\mathbf{o}_j \gets \mathbf{S}_j \mathbf{q}_j$
\EndFor
\State \textbf{Recurrencia entre fragmentos:}
\State $\mathbf{S}_i \gets \mathbf{S}_C$ del en-fragmento; propagar entre fragmentos mediante máscaras $\prod \alpha$ acumulativas
\State \textbf{Compuerta de salida:} $\mathbf{O}_{[i]} \gets \text{GroupNorm}(\mathbf{O}^\text{crudo}) \otimes \text{SiLU}(W_g \mathbf{X}_{[i]})$
\Return $\mathbf{O}_{[i]}$, $\mathbf{S}_i$
\end{algorithmic}
\end{algorithm}

\subsubsection{Rendimiento Empírico y Compensaciones}

Gated DeltaNet se ha integrado en los modelos de producción Qwen3-Next y Qwen3.5 de Alibaba. Empíricamente:
\[
\begin{array}{l|cccc}
\textbf{Tarea} & \textbf{Gated DeltaNet} & \textbf{Mamba2} & \textbf{DeltaNet} & \textbf{GLA} \\
\hline
\text{Modelado de Lenguaje} & \mathbf{1.0\times} & 1.08\times & 1.05\times & 1.12\times \\
\text{Recuperación Asociativa} & \mathbf{1.0\times} & --(sin perfecto) & \mathbf{1.0\times} & 0.75 \\
\text{Extrapolación de Longitud} & \mathbf{1.0\times} & 0.98\times & 0.96\times & 0.94\times \\
\text{Rendimiento (tokens/seg)} & 0.85\times & \mathbf{1.0\times} & 0.82\times & 0.88\times
\end{array}
\]
Gated DeltaNet logra el mejor equilibrio entre diversas tareas a costa de un rendimiento ligeramente reducido debido a matrices de transición más ricas.

\subsection{4. HGRN2: Compuerta Jerárquica con Expansión de Producto Externo}

\subsubsection{Fundamento Matemático}

HGRN2 \citep{qin_hgrn2_2024} usa una expansión de estado basada en producto externo con \textbf{compuertas de olvido acotadas inferiormente de forma jerárquica}. La actualización de estado es:
\[
\mathbf{S}_t = \text{diag}(\mathbf{g}_t) \cdot \mathbf{S}_{t-1} + \mathbf{v}_t \mathbf{k}_t^\top, \quad \mathbf{o}_t = \mathbf{S}_t \mathbf{q}_t
\]
donde $\mathbf{g}_t \in [b_\ell, 1]^d$ es una compuerta de olvido con cota inferior para la capa $\ell$. Las cotas satisfacen:
\[
0 \le b_{\ell_{\text{poco profunda}}} < b_{\ell_{\text{media}}} < b_{\ell_{\text{profunda}}} \le 1
\]
es decir, las cotas aumentan (se vuelven menos restrictivas) en las capas más profundas. La compuerta en sí se calcula como:
\[
\mathbf{g}_t = b_\ell + (1 - b_\ell) \cdot \sigma(W_g \mathbf{x}_t)
\]
donde $\sigma$ es sigmoide. Cuando $W_g$ produce logits débilmente negativos, $\mathbf{g}_t \approx b_\ell$ (retención forzada en capas superficiales). Cuando las salidas son fuertemente positivas, $\mathbf{g}_t \approx 1$ (refresco de memoria en cualquier capa).

La estructura jerárquica anima a diferentes capas a especializarse en diferentes escalas temporales: las capas superficiales modelan dependencias locales (alta retención debido a la cota $b_\ell$), mientras que las capas profundas capturan estructura de largo alcance (compuerta flexible con cota superior alta).

\begin{algorithm}[H]
\caption{Paso Adelante de HGRN2 con Cotas Jerárquicas}
\begin{algorithmic}[1]
\Require Entrada $\mathbf{x}_t$; estado anterior $\mathbf{S}_{t-1}$; índice de capa $\ell \in [0, L-1]$
\Require Cotas no decrecientes: $0 \le b_0 < b_1 < \cdots < b_{L-1} \le 1$
\Ensure Salida $\mathbf{o}_t$ y estado $\mathbf{S}_t$
\State Cota de retención específica de capa: $b \gets b_\ell$
\State Proyectar: $\mathbf{q}_t \gets W_q \mathbf{x}_t$, $\mathbf{k}_t \gets W_k \mathbf{x}_t$, $\mathbf{v}_t \gets W_v \mathbf{x}_t$
\State Calcular compuerta de olvido con cota: $\mathbf{g}_t \gets b + (1-b) \cdot \sigma(W_g \mathbf{x}_t)$ \Comment{Confinado a $[b, 1]$}
\State Multiplicación diagonal (vectorizada): $\mathbf{S}_t \gets \mathbf{S}_{t-1} \circ \text{diag}(\mathbf{g}_t) + \mathbf{v}_t \mathbf{k}_t^\top$
\State Recuperar: $\mathbf{o}_t \gets \mathbf{S}_t \mathbf{q}_t$
\State Normalización opcional: $\mathbf{o}_t \gets \text{RMSNorm}(\mathbf{o}_t)$
\Return $\mathbf{o}_t$, $\mathbf{S}_t$
\end{algorithmic}
\end{algorithm}

\subsubsection{Escalado y Resultados Empíricos}

A escala de 3B en 100B tokens, HGRN2 supera ligeramente a Mamba2 y a los Transformers con arquitectura LLaMA en modelado de lenguaje. La compuerta jerárquica proporciona modelado temporal multi-escala sin variantes arquitectónicas explícitas por capa. Sin embargo, la compuerta diagonal con valor vectorial (en oposición a la selección elemento a elemento) proporciona menos flexibilidad por elemento que las variantes de regla delta.

\subsection{5. Forgetting Transformer (FoX): Compuerta en el Espacio de Logits Softmax}

\subsubsection{Fundamento Matemático}

Forgetting Transformer (FoX), propuesto por Lin et al.~\citep{lin_forgetting_transformer_2025}, incrusta una compuerta de olvido directamente en los logits de atención softmax. En lugar de reemplazar softmax con recurrencia lineal, FoX preserva la expresividad completa de softmax mientras añade control de actualidad mediante un sesgo de logit dependiente de los datos.

Para cada posición de token $t$ y todas las claves $j \le t$, calcular un escalar de compuerta de olvido:
\[
f_t = \sigma(w_f^\top \mathbf{x}_t + b_f)
\]
El sesgo de logit en la posición $(i,j)$ es el producto acumulativo log-olvido:
\[
d_{ij} = \sum_{l=j+1}^{i} \log f_l = \log \prod_{l=j+1}^{i} f_l
\]
La salida de atención completa se convierte en:
\[
\mathbf{O} = \text{softmax}\left(\frac{\mathbf{Q}\mathbf{K}^\top}{\sqrt{d_k}} + \mathbf{D}\right) \mathbf{V}
\]
donde $\mathbf{D}_{ij} = d_{ij}$. Equivalentemente:
\[
\text{Atencion}(i,j) = \frac{\exp(\mathbf{q}_i^\top \mathbf{k}_j / \sqrt{d_k} + d_{ij})}{\sum_{j'=1}^{i} \exp(\mathbf{q}_i^\top \mathbf{k}_{j'} / \sqrt{d_k} + d_{ij'})}
\]

Esta ponderación reduce el peso de los tokens más antiguos multiplicativamente: un token de hace 10 pasos se escala por $\prod_{l=j+1}^{i} f_l$, que es típicamente mucho menor que 1 si $f_t < 1$ en promedio.

El mecanismo es matemáticamente equivalente a una \textbf{variante aprendible dependiente de los datos de ALiBi (Atención con Sesgos Lineales)}, donde trabajos anteriores usaban $d_{ij} = -\infty \cdot |i-j|$ fijo o $d_{ij} = -(i-j)$.

\begin{algorithm}[H]
\caption{FoX: Atención con Olvido y Sesgo de Actualidad}
\begin{algorithmic}[1]
\Require Consultas $\mathbf{Q}$, claves $\mathbf{K}$, valores $\mathbf{V}$; entrada $\mathbf{X}$
\Require Parámetros de compuerta de olvido: $w_f \in \mathbb{R}^d$, $b_f \in \mathbb{R}$
\Ensure Salida de atención $\mathbf{O}$
\State \textbf{Calcular compuertas de olvido:}
\For{$t = 1$ to $T$}
    \State $f_t \gets \sigma(w_f^\top \mathbf{x}_t + b_f)$ \Comment{Probabilidad de olvido por token}
\EndFor
\State \textbf{Calcular sesgos acumulativos de log-olvido:}
\For{$i = 1$ to $T$}
    \For{$j = 1$ to $i$}
        \State $d_{ij} \gets \sum_{l=j+1}^{i} \log f_l$ \Comment{Producto acumulativo en espacio logarítmico}
    \EndFor
\EndFor
\State \textbf{SDPA estándar con sesgo:}
\State Calcular puntajes: $\mathbf{S} \gets \mathbf{Q}\mathbf{K}^\top / \sqrt{d_k}$
\State Agregar sesgo: $\mathbf{S} \gets \mathbf{S} + \mathbf{D}$
\State Aplicar softmax: $\mathbf{A} \gets \text{softmax}(\mathbf{S}, \text{dim}=-1)$
\State Ponderar valores: $\mathbf{O} \gets \mathbf{A} \mathbf{V}$
\Return $\mathbf{O}$
\end{algorithmic}
\end{algorithm}

\subsubsection{Integración con FlashAttention}

La ventaja clave de FoX es la compatibilidad con FlashAttention y las implementaciones optimizadas existentes. Los sesgos de olvido se añaden en el cálculo de logits softmax, que ya es una operación central en el algoritmo por bloques de FlashAttention. La sobrecarga es:
\[
\text{Sobrecarga} \approx 0.5--2\% \text{ del tiempo de pared}
\]
ya que la adición de sesgo se amortiza entre las operaciones matmul.

\subsubsection{Fortalezas y Limitaciones}

FoX logra resultados de última generación en extrapolación de longitud, recuperación casi perfecta de aguja en el pajar y comprensión superior de contexto largo en comparación con los Transformers. Sin embargo, sigue siendo $\mathcal{O}(L^2 d)$ en entrenamiento y $\mathcal{O}(L d)$ por paso en inferencia con caché KV, limitando las longitudes de secuencia extremas.

\subsection{6. Atención con Compuerta (Post-SDPA Sigmoide)}

\subsubsection{Fundamento Matemático}

El Mejor Artículo de NeurIPS 2025 del equipo de Qwen propone aplicar una compuerta sigmoide \textbf{después} de la atención de producto punto escalado (SDPA):
\[
\mathbf{Y} = \text{SDPA}(\mathbf{Q}, \mathbf{K}, \mathbf{V}) = \text{softmax}\left(\frac{\mathbf{Q}\mathbf{K}^\top}{\sqrt{d_k}}\right) \mathbf{V}
\]
\[
\mathbf{Y}' = \mathbf{Y} \odot \sigma\left(\mathbf{X} \mathbf{W}_\theta\right)
\]
donde el puntaje de compuerta $\sigma(\mathbf{X} \mathbf{W}_\theta)$ puede tener forma $\mathbb{R}^{n \times q \times d_k}$ (compuerta elemento a elemento por cabeza de consulta) o $\mathbb{R}^{n \times q}$ (compuerta fija por cabeza). En más de 30 variantes y modelos MoE de 15B + densos de 1.7B entrenados en 3.5T tokens, el equipo encontró:

\begin{itemize}[leftmargin=*]
    \item La compuerta por cabeza añade solo $\sim 1.6M$ parámetros a un modelo de 15B (sobrecarga insignificante).
    \item Las compuertas efectivas son \textbf{dispersas}: activación media $\approx 0.116$ (es decir, el 88\% son cero o casi cero).
    \item Las compuertas actúan como filtros dependientes de la consulta, suprimiendo canales de bajo valor mientras preservan señales de alto valor.
    \item Las compuertas eliminan demostrablemente el fenómeno de \textbf{sumidero de atención}, donde el patrón de atención se vuelve dominado por un único token temprano.
\end{itemize}

\begin{algorithm}[H]
\caption{Atención Softmax con Compuerta (Post-SDPA)}
\begin{algorithmic}[1]
\Require Consultas $\mathbf{Q}$, Claves $\mathbf{K}$, Valores $\mathbf{V}$; entrada $\mathbf{X}$
\Require Matriz de proyección de compuerta $W_g \in \mathbb{R}^{d \to d_{\text{compuerta}}}$ donde $d_{\text{compuerta}} \in \{1, d_k\}$
\Ensure Salida con compuerta $\mathbf{Y}'$
\State \textbf{SDPA estándar:}
\State Calcular atención: $\mathbf{A} \gets \text{softmax}(\mathbf{Q}\mathbf{K}^\top / \sqrt{d_k})$ 
\State Ponderar valores: $\mathbf{Y} \gets \mathbf{A} \mathbf{V}$
\State \textbf{Calcular compuertas:}
\State Proyectar entrada: $\mathbf{G} \gets \mathbf{X} W_g$  \Comment{forma: $[n, q, d_{\text{compuerta}}]$}
\State Aplicar sigmoide: $\mathbf{G} \gets \sigma(\mathbf{G})$ \Comment{Compuerta elemento a elemento}
\State \textbf{Aplicar compuerta:}
\If{$d_{\text{compuerta}} = 1$}
    \State Transmitir y escalar: $\mathbf{Y}' \gets \mathbf{Y} \cdot \mathbf{G}$ \Comment{Escalado por cabeza}
\Else
    \Comment{$d_{\text{compuerta}} = d_k$}
    \State Modulación elemento a elemento: $\mathbf{Y}' \gets \mathbf{Y} \odot \mathbf{G}$ \Comment{Compuerta por dimensión}
\EndIf
\Return $\mathbf{Y}'$
\end{algorithmic}
\end{algorithm}

\subsubsection{Hallazgos Clave}

\begin{itemize}[leftmargin=*]
    \item \textbf{Supresión del sumidero de atención}: La compuerta rompe el ciclo de retroalimentación donde softmax se concentra en el primer token, permitiendo una atención más equilibrada.
    \item \textbf{Estabilidad de entrenamiento}: Los modelos con compuerta toleran tasas de aprendizaje más grandes ($+30\%$ mayor $\eta_{max}$ estable).
    \item \textbf{Sobrecarga mínima}: El impacto en latencia es solo del $1.6\%$ en GPUs H100; el rendimiento cae como máximo $2-3\%$.
    \item \textbf{Mejora en la ley de escalado}: Mejora la pérdida en todas las escalas de modelo, desde 800M hasta 110B parámetros de forma consistente.
\end{itemize}

\normalsize

\subsection{7. Gated DeltaNet-2: Borrado y Escritura Canalizados Decouplados}

\subsubsection{Fundamento Matemático}

Gated DeltaNet-2 \citep{hatamizadeh_gated_deltanet2_2026} aborda una limitación compartida de Gated DeltaNet y Kimi Delta Attention (KDA): ambos usan una única \emph{compuerta delta escalar} $\beta_t$ para controlar simultáneamente (i) cuánto contenido antiguo se borra en la dirección de lectura del eje de claves y (ii) cuánto valor nuevo se escribe en el eje de valores. Gated DeltaNet-2 \textbf{desacopla} estos dos roles en una compuerta de borrado canalizada $\mathbf{b}_t$ (eje de claves) y una compuerta de escritura canalizada $\mathbf{w}_t$ (eje de valores), heredando además el decaimiento canalizado $\boldsymbol{\alpha}_t$ de KDA. Con estado $\mathbf{S}_t \in \mathbb{R}^{d_k \times d_v}$, decaimiento $\mathbf{D}_t = \mathrm{Diag}(\boldsymbol{\alpha}_t)$, dirección de borrado $\mathbf{e}_t = \mathbf{b}_t \odot \mathbf{k}_t$ y objetivo de escritura $\mathbf{z}_t = \mathbf{w}_t \odot \mathbf{v}_t$, la actualización de estado es:
\[
\mathbf{S}_t = \left(\mathbf{I} - \mathbf{k}_t \mathbf{e}_t^\top\right) \mathbf{D}_t \, \mathbf{S}_{t-1} + \mathbf{k}_t \mathbf{z}_t^\top, \qquad \mathbf{o}_t = \mathbf{S}_t^\top \mathbf{q}_t
\]
o, expandiendo los productos de las compuertas:
\[
\mathbf{S}_t = \left(\mathbf{I} - \mathbf{k}_t (\mathbf{b}_t \odot \mathbf{k}_t)^\top\right) \mathbf{D}_t \, \mathbf{S}_{t-1} + \mathbf{k}_t (\mathbf{w}_t \odot \mathbf{v}_t)^\top
\]
Las dos compuertas se producen mediante proyecciones lineales independientes seguidas de sigmoide:
\[
\mathbf{b}_t = \sigma(\mathbf{W}_b \mathbf{x}_t), \qquad \mathbf{w}_t = \sigma(\mathbf{W}_w \mathbf{x}_t)
\]
y el decaimiento canalizado sigue una parametrización de log-decaimiento (calculada en fp32 para evitar pérdida de precisión en productos acumulativos largos):
\[
\mathbf{g}_t = \mathbf{a} - \mathrm{softplus}(\boldsymbol{\delta}), \qquad \boldsymbol{\alpha}_t = \exp(\mathbf{g}_t)
\]
Los roles estructurales de las dos compuertas son complementarios:
\begin{itemize}[leftmargin=*]
    \item $\mathbf{b}_t$ (eje de claves) hace que la \textbf{dirección de lectura} sea selectiva por canal: determina qué canales de clave se eliminan antes de escribir.
    \item $\mathbf{w}_t$ (eje de valores) hace que el \textbf{objetivo de escritura} sea selectivo por canal: determina qué canales de valor se depositan en memoria.
    \item $\boldsymbol{\alpha}_t$ (por canal de clave) proporciona decaimiento global, como en KDA.
\end{itemize}
\textbf{Reducciones.} Fijar $\mathbf{b}_t = \beta_t \mathbf{1}_{d_k}$ y $\mathbf{w}_t = \beta_t \mathbf{1}_{d_v}$ recupera KDA; colapsar además $\boldsymbol{\alpha}_t$ a un escalar recupera Gated DeltaNet. Desde la perspectiva de aprendizaje en línea/pesos rápidos, la actualización es el minimizador de:
\[
\mathcal{L}_t(\mathbf{S}) = \tfrac{1}{2}\|\mathbf{S} - \mathbf{D}_t \mathbf{S}_{t-1}\|_F^2 - 2\left\langle \mathbf{S}^\top \mathbf{k}_t,\; \mathbf{z}_t - (\mathbf{D}_t \mathbf{S}_{t-1})^\top \mathbf{e}_t \right\rangle
\]

\begin{algorithm}[H]
\caption{Avance Recurrente de Gated DeltaNet-2 (Forma de Referencia)}
\begin{algorithmic}[1]
\Require Entrada $\mathbf{x}_t$; estado previo $\mathbf{S}_{t-1} \in \mathbb{R}^{d_k \times d_v}$
\Require Proyecciones: $W_q, W_k, W_v, W_b, W_w$; parámetros de log-decaimiento $\mathbf{a}, \boldsymbol{\delta}$
\Ensure Salida $\mathbf{o}_t$ y estado $\mathbf{S}_t$
\State $\mathbf{q}_t \gets \mathrm{normalize}(W_q \mathbf{x}_t)$, $\mathbf{k}_t \gets \mathrm{normalize}(W_k \mathbf{x}_t)$, $\mathbf{v}_t \gets \mathrm{silu}(W_v \mathbf{x}_t)$
\State $\mathbf{b}_t \gets \sigma(W_b \mathbf{x}_t)$ \Comment{compuerta de borrado eje-clave, $\mathbb{R}^{d_k}$}
\State $\mathbf{w}_t \gets \sigma(W_w \mathbf{x}_t)$ \Comment{compuerta de escritura eje-valor, $\mathbb{R}^{d_v}$}
\State $\boldsymbol{\alpha}_t \gets \exp(\mathbf{a} - \mathrm{softplus}(\boldsymbol{\delta}))$ \Comment{decaimiento canalizado, fp32}
\State Dirección de borrado: $\mathbf{e}_t \gets \mathbf{b}_t \odot \mathbf{k}_t$; objetivo de escritura: $\mathbf{z}_t \gets \mathbf{w}_t \odot \mathbf{v}_t$
\State Decaer estado: $\mathbf{S}_t \gets \mathrm{Diag}(\boldsymbol{\alpha}_t)\, \mathbf{S}_{t-1}$
\State Lectura: $\mathbf{r}_t \gets \mathbf{S}_t^\top \mathbf{e}_t$ \Comment{suma sobre eje de claves}
\State Actualización delta: $\mathbf{S}_t \gets \mathbf{S}_t - \mathbf{k}_t \mathbf{r}_t^\top + \mathbf{k}_t \mathbf{z}_t^\top$
\State Recuperar: $\mathbf{o}_t \gets \mathbf{S}_t^\top \mathbf{q}_t$
\State Compuerta de salida: $\mathbf{o}_t \gets \mathrm{GroupNorm}(\mathbf{o}_t) \odot \mathrm{silu}(W_g \mathbf{x}_t)$
\State \Return $\mathbf{o}_t$, $\mathbf{S}_t$
\end{algorithmic}
\end{algorithm}

El algoritmo de entrenamiento por fragmentos extiende la representación WY de KDA absorbiendo el decaimiento canalizado en una recurrencia delta asimétrica sobre un estado normalizado por decaimiento, permitiendo paralelismo rico en matmuls en tensor cores. Como $\mathbf{b}_t$ y $\mathbf{w}_t$ son operadores diagonales por canal que varían por fila, el atajo de post-escalado escalar de KDA se rompe: el paso hacia atrás requiere una \textbf{inversión WY consciente de las compuertas} que incorpora las compuertas en cada producto escalar que acumula el gradiente.

\subsubsection{Rendimiento Empírico y Compromisos}

A 1.3B parámetros entrenados con 100B tokens FineWeb-Edu, Gated DeltaNet-2 supera a Mamba-2, Mamba-3, Gated DeltaNet y KDA en modelado de lenguaje, razonamiento de sentido común y---más notablemente---recuperación de contexto largo del mundo real y RULER (especialmente multi-key needle-in-a-haystack). Su rendimiento en H100 se mantiene casi plano ($38.0 \to 36.1$ Kt/s) al crecer la longitud de secuencia de 2K a 16K, donde el de un Transformer cae bruscamente. Los estudios de ablación muestran que la compuerta de borrado $\mathbf{b}_t$ explica la mayor parte de la ganancia; la compuerta de escritura $\mathbf{w}_t$ aporta una mejora adicional menor. La variante híbrida (Gated DeltaNet-2 recurrente + atención de ventana deslizante) reduce aún más la brecha con los modelos de atención pura en tareas que requieren agregación de evidencia local.

\begin{table}[H]
\centering
\footnotesize
\begin{tabularx}{\linewidth}{lXX}
\toprule
\textbf{Aspecto} & \textbf{Fortaleza} & \textbf{Limitación} \\
\midrule
Recuperación & La mejor entre modelos delta recurrentes en multi-key NIAH. & La variante recurrente aún supera al Transformer en NQ/DROP. \\
Rendimiento & Escalado casi plano 2K$\to$16K; pequeña sobrecarga vs KDA. & Requiere un nuevo kernel Triton WY consciente de las compuertas. \\
Memoria & Estado de inferencia constante $\mathcal{O}(d^2)$. & Rango de borrado de autovalores negativos $[0,2]$ no da ganancia consistente a 1.3B. \\
\bottomrule
\end{tabularx}
\caption{Compromisos de Gated DeltaNet-2 a escala 1.3B / 100B tokens \citep{hatamizadeh_gated_deltanet2_2026}.}
\label{tab:gated-deltanet2-tradeoffs-es}
\end{table}

\normalsize

\subsection{Análisis Comparativo: Taxonomía de Mecanismos con Compuerta}

\footnotesize
\begin{longtable}{p{1.8cm}p{1.7cm}p{1.8cm}p{1.6cm}p{2.0cm}p{2.4cm}p{1.3cm}}
\toprule
\textbf{Arquitectura} & \textbf{Año de Publicación} & \textbf{Representación de Estado} & \textbf{Tipo de Compuerta} & \textbf{Complejidad de Entrenamiento} & \textbf{Complejidad de Inferencia} & \textbf{Ref} \\
\midrule
\endhead
GLA & 2023 & Matriz (recurrente) & Decaimiento diagonal & $\mathcal{O}(Ld^2)$ & $\mathcal{O}(d^2)$ memoria & \citep{yang_gla_2023} \\
DeltaNet & 2024 & Matriz (recurrente) & Escritura escalar ($\beta$) & $\mathcal{O}(Ld^2)$ & $\mathcal{O}(d^2)$ memoria & \citep{yang_deltanet_2024} \\
Gated DeltaNet & 2024 & Matriz (recurrente) & Decaimiento + escritura & $\mathcal{O}(Ld^2)$ & $\mathcal{O}(d^2)$ memoria & \citep{yang_gated_deltanet_2024} \\
RetNet & 2023 & Matriz (recurrente) & Exponencial fija & $\mathcal{O}(Ld^2)$ & $\mathcal{O}(d)$ por paso & \citep{sun_retentive_2023} \\
HGRN2 & 2024 & Matriz (recurrente) & Diagonal (acotada) & $\mathcal{O}(Ld^2)$ & $\mathcal{O}(d^2)$ memoria & \citep{qin_hgrn2_2024} \\
FoX & 2025 & Atención completa & Sesgo en espacio de logits & $\mathcal{O}(L^2d)$ & $\mathcal{O}(L)$ caché KV & \citep{lin_forgetting_transformer_2025} \\
Gated Softmax & 2025 & Atención completa & Sigmoide post-SDPA & $\mathcal{O}(L^2d)$ & $\mathcal{O}(L)$ caché KV & \citep{qiu_gated_attention_2025} \\
Gated DeltaNet-2 & 2026 & Matriz (recurrente) & Borrado + escritura decouplados (canalizado) & $\mathcal{O}(Ld^2)$ & $\mathcal{O}(d^2)$ memoria & \citep{hatamizadeh_gated_deltanet2_2026} \\
\bottomrule
\end{longtable}
\normalsize

\medskip

\footnotesize
\begin{longtable}{p{1.8cm}p{2.0cm}p{2.2cm}p{2.2cm}p{2.2cm}p{1.8cm}}
\toprule
\textbf{Arquitectura} & \textbf{Entrenable} & \textbf{Recuperación en Contexto} & \textbf{Recuperación Asociativa} & \textbf{Extrapolación de Longitud} & \textbf{Ventaja Clave} \\
\midrule
\endhead
GLA & Sí & Moderada & Débil & Buena (20K) & Eficiencia en hardware; paralelismo por fragmentos \\
DeltaNet & Sí & Fuerte & Perfecta (MQAR) & Buena & Semántica de corrección de errores; recuperación sólida \\
Gated DeltaNet & Sí & Muy fuerte & Perfecta & Excelente & Equilibrada: olvido + actualizaciones dirigidas \\
RetNet & Sí & Débil & Débil & Buena & Simplicidad; no necesita caché KV \\
HGRN2 & Sí & Moderada & Débil & Moderada & Modelado temporal multi-escala \\
FoX & Sí & Muy fuerte & Muy fuerte (casi perfecta) & Excelente & Preserva softmax; sobrecarga mínima \\
Gated Softmax & Sí & Fuerte & Buena & Buena & Simplicidad; no reemplaza SDPA \\
Gated DeltaNet-2 & Sí & Muy fuerte & Muy fuerte & Excelente & Borrado/escritura canalizados decouplados; mejor RULER multi-key \\
\bottomrule
\end{longtable}
\normalsize

\medskip

\footnotesize
\begin{longtable}{p{2.0cm}p{2.4cm}p{2.8cm}p{3.4cm}}
\toprule
\textbf{Arquitectura} & \textbf{Rendimiento Típico} & \textbf{Memoria por Inferencia} & \textbf{Caso de Uso Principal} \\
\midrule
\endhead
GLA & Alto (CUDA por fragmentos) & $O(d^2)$ constante & Contexto largo con restricciones de memoria \\
DeltaNet & Medio-Alto & $O(d^2)$ constante & Recuperación y razonamiento asociativo \\
Gated DeltaNet & Medio & $O(d^2)$ constante & Producción: Qwen3-Next, rendimiento equilibrado \\
RetNet & Alto & $O(d)$ por paso & Inferencia extremadamente rápida; modelos simples \\
HGRN2 & Medio-Alto & $O(d^2)$ constante & Patrones temporales multi-escala \\
FoX & Moderado (cuadrático) & $O(Ld)$ caché KV & Extrapolación de longitud; generación de contexto largo \\
Gated Softmax & Alto (sobrecarga mínima) & $O(Ld)$ caché KV & Mejora directa para modelos existentes \\
Gated DeltaNet-2 & Alto (escalado casi plano) & $O(d^2)$ constante & Recuperación de contexto largo; multi-key NIAH \\
\bottomrule
\end{longtable}
\normalsize

\subsection{Principio Unificado de Compuerta}

Las ocho arquitecturas comparten un principio común: \textbf{la compuerta es un mecanismo para controlar el flujo de información y la retención selectiva de memoria}. Las decisiones de diseño específicas divergen a lo largo de varias dimensiones:

\begin{enumerate}[leftmargin=*]
    \item \textbf{Ubicación de la compuerta}: Actualizaciones de estado recurrente (GLA, DeltaNet, Gated DeltaNet, Gated DeltaNet-2, HGRN2) versus espacio de logits/salida softmax (FoX, Gated Softmax, RetNet).
    \item \textbf{Granularidad de control}: Por dimensión (GLA, HGRN2 diagonal), por clave (DeltaNet), por token (FoX), por cabeza (Gated Softmax) o borrado/escritura canalizados decouplados (Gated DeltaNet-2).
    \item \textbf{Dependencia de datos}: Completamente dependiente de datos (GLA, DeltaNet, Gated DeltaNet, Gated DeltaNet-2, FoX, Gated Softmax) versus esquemas fijos (RetNet con decaimiento exponencial).
    \item \textbf{Compensación de expresividad}: Eficiencia recurrente lineal (GLA, DeltaNet, Gated DeltaNet, Gated DeltaNet-2, HGRN2, RetNet) versus expresividad completa de softmax (FoX, Gated Softmax).
\end{enumerate}

\subsection{Implementación y Orientación Práctica}

\subsubsection{Cuándo Usar Cada Arquitectura}

\begin{enumerate}[leftmargin=*]
    \item \textbf{GLA}: Inferencia rápida con contextos moderados a largos (2K--20K tokens), cuando los kernels CUDA personalizados son aceptables y el rendimiento de recuperación no es crítico.
    \item \textbf{DeltaNet}: Fuerte recuperación asociativa y tareas de aprendizaje en contexto; bueno para tareas sintéticas (MQAR) y pruebas de asociación clave-valor.
    \item \textbf{Gated DeltaNet}: Modelos de producción que requieren el mejor equilibrio de recuperación, olvido y rendimiento (probado en Qwen3-Next; ICLR 2025).
    \item \textbf{Gated DeltaNet-2}: Cargas de recuperación de contexto largo (multi-key NIAH) donde el borrado/escritura canalizado decouplado da el recuerdo más fuerte; cuando se dispone de un kernel WY consciente de las compuertas.
    \item \textbf{RetNet}: Inferencia extremadamente eficiente sin almacenamiento en caché KV; adecuado para despliegues en dispositivos móviles/de borde o modelos de juguete.
    \item \textbf{HGRN2}: Jerarquías temporales multi-escala; cuando los límites de olvido específicos por capa son deseables.
    \item \textbf{FoX}: Extrapolación de longitud y comprensión de contexto largo; cuando el entrenamiento cuadrático es aceptable y no se desea reemplazar softmax.
    \item \textbf{Gated Softmax}: Mejora rápida para modelos softmax existentes con cambios mínimos de código; mejor para profesionales que desean ganancias inmediatas sin una reestructuración importante.
\end{enumerate}

\subsubsection{Consideraciones de Hardware}

\begin{table}[H]
\centering
\small
\begin{tabularx}{0.95\linewidth}{Xp{3.5cm}p{3.5cm}}
\toprule
\textbf{Objetivo de Hardware} & \textbf{Arquitecturas Recomendadas} & \textbf{Notas} \\
\midrule
GPU (H100/A100) & Gated Softmax, FoX, Gated DeltaNet, Gated DeltaNet-2 & Implementaciones maduras de softmax/lineal. GLA/DeltaNet requieren kernels personalizados. \\
TPU (alta memoria) & GLA, DeltaNet, Gated DeltaNet, Gated DeltaNet-2, HGRN2 & Hardware soporta matmuls eficientes; los métodos recurrentes lineales brillan. \\
Móvil/Borde & RetNet, Gated Softmax ligero & Inferencia con memoria constante es crítica. \\
\bottomrule
\end{tabularx}
\end{table}

\section{Annex E: Conceptual Introduction---Transformers and Attention for Beginners}
\label{sec:annex-tutorial}

Este apéndice proporciona una introducción accesible a los conceptos fundamentales de los transformers y el mecanismo de atención, dirigido a lectores sin conocimientos profundos previos en aprendizaje profundo. Los transformers son el motor detrás de los sistemas modernos de inteligencia artificial como ChatGPT, pero su funcionamiento puede entenderse mediante analogías del mundo real.

\subsection{¿Qué es un Transformer?}

Un transformer es una arquitectura neuronal que procesa texto o secuencias de datos interpretando el significado de cada palabra mientras considera todas las demás palabras en contexto. Imagina que lees una oración:

\begin{center}
\textit{``El banco estaba lleno de gente esperando. María fue al banco.''}
\end{center}

¿Qué significa ``banco'' en cada caso? En la primera oración, se refiere a una institución financiera (o un asiento). En la segunda, no está claro sin contexto. Un transformer resuelve esto automáticamente al mirar todas las palabras circundantes. En la segunda oración, al ver ``María fue'', el modelo comprende mejor que probablemente se trata de una ribera (donde va a nadar) o una institución financiera (donde va a realizar una transacción).

Esta capacidad de entender el significado completo de una palabra considerando toda la oración se llama \emph{contextualización}, y es el corazón de cómo funcionan los transformers modernos.

\subsection{El Mecanismo de Atención: Una Analogía}

El mecanismo de atención es el ``truco mágico'' que permite a los transformers entender el contexto. Piensa en ello como participar en una conversación importante en una habitación ruidosa:

\begin{enumerate}[leftmargin=*]
    \item \textbf{Mucho ruido}: Alguien está hablando y es muy importante para ti.
    \item \textbf{Ignoras el resto}: Tu cerebro automáticamente enfoca la atención en esa persona, ignorando otros sonidos.
    \item \textbf{Entiendes mejor}: Al concentrarte, captas cada palabra claramente.
\end{enumerate}

El mecanismo de atención neuronal funciona de la misma manera: cada palabra ``pregunta'' qué tan importante es cada otra palabra para entender su significado, luego ``enfoca'' su atención en las más relevantes.

\subsection{Ejemplo Práctico Paso a Paso}

Veamos cómo un transformer entiende la palabra ``come'' en dos contextos diferentes:

\subsubsection{Contexto 1: ``El gato come pescado''}

El transformer, al procesar ``come'', pregunta internamente:
\begin{itemize}[leftmargin=*]
    \item ¿Qué tan relevante es ``El''? Poco (es solo un artículo). \textbf{Atención baja}.
    \item ¿Qué tan relevante es ``gato''? Muy relevante (es el sujeto, quien come). \textbf{Atención alta}.
    \item ¿Qué tan relevante es ``pescado''? Muy relevante (es lo que se come). \textbf{Atención alta}.
\end{itemize}

Por lo tanto, la representación mental de ``come'' se enfoca principalmente en ``gato'' y ``pescado''.

\subsubsection{Contexto 2: ``El restaurante come los márgenes de beneficio''}

Aquí el transformer pregunta en un contexto diferente:
\begin{itemize}[leftmargin=*]
    \item ¿Qué tan relevante es ``restaurante''? Muy relevante (es el sujeto). \textbf{Atención alta}.
    \item ¿Qué tan relevante es ``beneficio''? Muy relevante (se ve afectado). \textbf{Atención alta}.
    \item ¿Qué tan relevante es ``márgenes''? Muy relevante (es lo que se consume). \textbf{Atención alta}.
\end{itemize}

La palabra ``come'' obtiene una representación completamente diferente porque ``atiende'' a palabras distintas en contextos distintos.

\subsection{Cómo Funciona la Atención Matemáticamente (Versión Simple)}

Detrás del cambio de atención hay matemáticas. Aquí está la versión simplificada sin demasiado detalle técnico:

\subsubsection{Paso 1: Consultas, Claves y Valores}

Cada palabra en la oración se transforma en tres versiones:
\begin{itemize}[leftmargin=*]
    \item \textbf{Consulta}: ``¿Qué información necesito saber?''
    \item \textbf{Clave}: ``¿Qué información tengo?''
    \item \textbf{Valor}: ``Aquí está mi información importante.''
\end{itemize}

Imagina en una biblioteca, cada visitante es una ``consulta'', cada libro es una ``clave'' y el contenido es el ``valor''. El bibliotecario (atención) empareja las consultas con las claves más relevantes para acceder al valor correcto.

\subsubsection{Paso 2: Compatibilidad}

El sistema examina qué tan \emph{compatible} es cada consulta con cada clave. Las consultas similares a las claves obtienen un puntaje de compatibilidad \emph{alto}. Esto se calcula multiplicando la consulta por la clave (matemáticamente, el producto punto).

\subsubsection{Paso 3: Enfoque}

Los puntajes de compatibilidad se convierten en ``pesos de enfoque''. Las compatibilidades altas significan ``enfocar atención aquí'', y las bajas significan ``ignorar esto''. Matemáticamente, una función llamada \texttt{softmax} convierte los puntajes en porcentajes (como: 40\% atención aquí, 35\% aquí, 25\% allá).

\subsubsection{Paso 4: Combinar}

Finalmente, el sistema combina todos los valores, ponderados por los pesos de enfoque. Si una palabra tiene un 80\% de atención en el sujeto, el 80\% de la información del sujeto se mezcla en la representación de esa palabra.

\subsection{Múltiples Cabezas de Atención: Múltiples Perspectivas}

Una característica clave es que los transformers no usan una única atención sino \emph{múltiples cabezas de atención} en paralelo. Es como si tuvieras 8 personas analizando la oración \emph{simultáneamente}, cada una prestando atención a diferentes aspectos:

\begin{itemize}[leftmargin=*]
    \item \textbf{Persona 1}: ``Me enfoco en las relaciones sujeto-verbo.''
    \item \textbf{Persona 2}: ``Busco el objeto directo.''
    \item \textbf{Persona 3}: ``Sigo la información sobre el tiempo verbal.''
    \item \textbf{Persona 4}: ``Busco modificadores y adjetivos.''
\end{itemize}

Cada ``cabeza'' aprende a enfocarse en diferentes patrones del lenguaje. Juntas, capturan una comprensión mucho más rica que una sola cabeza.

\subsection{Apilar Capas}

Los transformers procesan información en \emph{múltiples capas} (generalmente de 12 a 48 para modelos prácticos). Imagina editar un ensayo:
\begin{enumerate}[leftmargin=*]
    \item \textbf{Primera pasada}: Corriges ortografía y gramática básica.
    \item \textbf{Segunda pasada}: Mejoras la claridad y la estructura de las oraciones.
    \item \textbf{Tercera pasada}: Aseguras consistencia y flujo narrativo.
\end{enumerate}

Los transformers funcionan de la misma manera: cada capa refina la comprensión del texto. La primera capa captura características básicas (palabras simples, géneros), mientras que las capas posteriores entienden conceptos complejos (relaciones entre palabras, significados abstractos).

\subsection{¿Por Qué es Importante?}

El mecanismo de atención fue revolucionario porque:
\begin{enumerate}[leftmargin=*]
    \item \textbf{Paralelismo}: Encuentra contexto para \emph{cualquier palabra con cualquier otra palabra}, sin procesarlas secuencialmente (a diferencia de sistemas anteriores). Esto lo hace muy rápido.
    \item \textbf{Flexibilidad}: Aprende qué patrones buscar automáticamente a partir de los datos, sin necesidad de programar reglas manualmente.
    \item \textbf{Escalabilidad}: Funciona desde textos pequeños hasta contextos con millones de palabras.
    \item \textbf{Capacidades generales}: El mismo mecanismo funciona para traducción, resumen, preguntas-respuestas, generación de texto, visión por computadora y más.
\end{enumerate}

\subsection{Desafíos Prácticos}

A pesar del extraordinario éxito de los transformers, presentan desafíos:

\subsubsection{Complejidad Computacional}

Si una oración tiene 1,000 palabras, el mecanismo de atención estándar debe comparar cada palabra con las otras 1,000 palabras. Eso es $1,000 \times 1,000 = 1,000,000$ de comparaciones. Para documentos largos con millones de palabras, esto se vuelve prohibitivamente costoso en tiempo y memoria.

\subsubsection{Requisitos de Memoria}

Especialmente durante la inferencia (cuando se usan modelos entrenados), almacenar la matriz de atención completa puede consumir enormes cantidades de RAM o memoria de GPU, limitando las longitudes de secuencia que se pueden procesar.

\subsubsection{Eficiencia del Dispositivo}

Los transformers fueron diseñados para TPUs y GPUs potentes. Ejecutarlos en teléfonos o dispositivos integrados es un desafío.

\subsection{Soluciones Modernas}

Para resolver estos desafíos, la investigación ha propuesto muchas variantes:

\begin{itemize}[leftmargin=*]
    \item \textbf{Atención Dispersa}: En lugar de comparar cada palabra con TODAS las demás, compara solo con un subconjunto estratégico (vecinos cercanos, patrones periódicos). Reduce la complejidad de $\mathcal{O}(n^2)$ a $\mathcal{O}(n \log n)$ o $\mathcal{O}(n)$.
    \item \textbf{Modelos Recurrentes}: Como Mamba, incorporan aspectos de las antiguas redes neuronales recurrentes pero con mejor eficiencia moderna.
    \item \textbf{Atención con Compuerta}: Mecanismos que aprenden selectivamente qué información pasar hacia adelante, reduciendo las necesidades de almacenamiento.
    \item \textbf{Cuantización}: Usar números más pequeños (enteros en lugar de decimales) para reducir la memoria sin perder demasiada precisión.
\end{itemize}

Estas innovaciones son lo que este kit de herramientas (Frankenstein) te permite experimentar fácilmente.

\subsection{Conclusiones Clave}

\begin{itemize}[leftmargin=*]
    \item Un \textbf{transformer} es una red neuronal que entiende el contexto del lenguaje muy bien.
    \item El \textbf{mecanismo de atención} permite que cada palabra ``se enfoque'' en qué otras palabras son relevantes para entender su significado.
    \item La atención funciona calculando la \textbf{compatibilidad} entre cada palabra (consulta) y todas las demás (claves), luego combinando información basada en esa compatibilidad (valores).
    \item \textbf{Múltiples cabezas} de atención exploran diferentes patrones en paralelo.
    \item \textbf{Múltiples capas} refinan progresivamente la comprensión.
    \item El principal desafío es que la atención estándar tiene complejidad cuadrática, lo que es costoso para textos largos.
    \item Muchas \textbf{soluciones modernas} (atención dispersa, modelos recurrentes, cuantización) abordan estos desafíos, y este kit de herramientas te permite explorarlas todas.
\end{itemize}

\end{document}
